\chapter{Sistem PDB} \label{sys:chapter}

%%%%%%%%%%%%%%%%%%%%%%%%%%%%%%%%%%%%%%%%%%%%%%%%%%%%%%%%%%%%%%%%%%%%%%%%%%%%%%

\section{Pengantar sistem PDB} \label{sec:introtosys}

%mbxINTROSUBSECTION

\sectionnotes{1 hingga 1,5 perkuliahan\EPref{, \S4.1 dalam \cite{EP}}\BDref{,
\S7.1 dalam \cite{BD}}}

\subsection{Sistem}

Sering kali kita tidak hanya memiliki satu variabel dependen dan satu persamaan.
Kita akan melihat bahwa kita mungkin memperoleh sistem yang terdiri atas beberapa
persamaan dan beberapa variabel dependen meskipun kita memulai dengan
satu persamaan.

Diberikan beberapa variabel dependen,
misalkan $y_1$, $y_2$, \ldots, $y_n$,
kita dapat memiliki persamaan diferensial yang melibatkan semuanya beserta
turunannya terhadap satu variabel independen $x$.
Sebagai contoh, $y_1'' = f(y_1',y_2',y_1,y_2,x)$.
Biasanya, ketika kita memiliki dua variabel dependen, kita memiliki dua persamaan
seperti
\begin{equation*}
\begin{aligned}
y_1'' & = f_1(y_1',y_2',y_1,y_2,x) , \\
y_2'' & = f_2(y_1',y_2',y_1,y_2,x) ,
\end{aligned}
\end{equation*}
untuk sejumlah fungsi $f_1$ dan $f_2$ yang diketahui.  Bentuk di atas kita sebut
\emph{\myindex{sistem persamaan diferensial}}.
Lebih tepatnya, bentuk di atas adalah \emph{\myindex{sistem orde kedua}}
PDB karena turunan orde kedua muncul.
Contoh suatu \emph{\myindex{sistem orde pertama}} ialah
\begin{equation*}
\begin{aligned}
x_1' & = g_1(x_1,x_2,x_3,t) , \\
x_2' & = g_2(x_1,x_2,x_3,t) , \\
x_3' & = g_3(x_1,x_2,x_3,t) ,
\end{aligned}
\end{equation*}
dengan $x_1,x_2,x_3$ sebagai
variabel dependen,
dan $t$ sebagai variabel independen.

Terminologi untuk sistem pada dasarnya sama dengan terminologi untuk
persamaan tunggal.
Untuk sistem di atas, suatu
\emph{solusi}\index{solusi sistem}
adalah himpunan tiga fungsi $x_1(t)$, $x_2(t)$, $x_3(t)$, sedemikian sehingga
\begin{equation*}
\begin{aligned}
x_1'(t) &= g_1\bigl(x_1(t),x_2(t),x_3(t),t\bigr) , \\
x_2'(t) &= g_2\bigl(x_1(t),x_2(t),x_3(t),t\bigr) , \\
x_3'(t) &= g_3\bigl(x_1(t),x_2(t),x_3(t),t\bigr) .
\end{aligned}
\end{equation*}

Kita juga mungkin memiliki
\emph{syarat awal}\index{syarat awal untuk sistem}.
Seperti pada persamaan tunggal, kita menentukan $x_1$, $x_2$, dan $x_3$ untuk suatu $t$ tetap.
Sebagai contoh, $x_1(0) = a_1$, $x_2(0) = a_2$, $x_3(0) = a_3$,
dengan $a_1$, $a_2$, dan $a_3$ berupa konstanta.
Untuk sistem orde kedua,
kita juga harus menentukan turunan pertama di titik awal.
Jika kita menemukan solusi yang memuat konstanta sebarang,
dengan penyelesaian terhadap konstanta-konstanta tersebut memberikan solusi bagi
setiap syarat awal, solusi ini kita sebut
\emph{solusi umum}\index{solusi umum sistem}.
Sebaiknya kita melihat suatu contoh sederhana.

\begin{example} \label{example:simplesystem}
Kadang-kadang suatu sistem mudah diselesaikan
dengan menyelesaikan satu variabel lalu variabel kedua.
Ambil
sistem orde pertama
\begin{equation*}
\begin{aligned}
y_1' & = y_1 , \\
y_2' & = y_1 - y_2 ,
\end{aligned}
\end{equation*}
dengan $y_1$, $y_2$ sebagai variabel dependen dan $x$ sebagai variabel
independen.  Tinjau syarat awal
$y_1(0) = 1$, $y_2(0) = 2$.

Kita perhatikan bahwa $y_1 = C_1 e^x$ adalah solusi umum persamaan pertama.
Kemudian kita substitusikan $y_1$ ini ke dalam persamaan kedua
dan memperoleh persamaan $y_2' = C_1e^x - y_2$, yang merupakan persamaan linear
orde pertama
yang mudah diselesaikan untuk $y_2$.
Dengan metode faktor integrasi, kita peroleh
\begin{equation*}
e^x y_2 = \frac{C_1}{2}e^{2x} + C_2 ,
\end{equation*}
atau $y_2 = \frac{C_1}{2}e^{x} + C_2e^{-x}$.  Oleh karena itu, solusi umum sistem tersebut
adalah
\begin{equation*}
y_1 = C_1 e^x , \qquad
y_2 = \frac{C_1}{2}e^{x} + C_2e^{-x} .
\end{equation*}
Kita selesaikan $C_1$ dan $C_2$ berdasarkan syarat awal.
Kita substitusikan $x=0$ untuk memperoleh
$1 = y_1(0) = C_1$ dan
$2 = y_2(0) = \nicefrac{C_1}{2} + C_2$,
atau dengan kata lain,
$C_1=1$ dan $C_2=\nicefrac{3}{2}$.  Jadi solusinya adalah
$y_1 = e^x$, dan
$y_2 = (\nicefrac{1}{2}) e^x + (\nicefrac{3}{2}) e^{-x}$.
\end{example}

Pada umumnya, kita tidak akan cukup beruntung untuk dapat menyelesaikan
setiap variabel secara terpisah seperti dalam
contoh tersebut---kita perlu menyelesaikan semua variabel sekaligus.
Walaupun kita tidak selalu dapat menyelesaikan satu variabel demi satu,
kita akan berusaha mempertahankan sebanyak mungkin manfaat teknik ini.
Dalam pengertian tertentu, untuk menyelesaikan sistem, kita tetap akan (berusaha) menyelesaikan
sekumpulan persamaan tunggal dan menggabungkan solusi-solusinya.
Untuk saat ini, jangan dahulu memikirkan cara menyelesaikan sistem.

Kita terutama akan meninjau \emph{\myindex{sistem linear}}.
\exampleref{example:simplesystem}
adalah suatu \emph{\myindex{sistem linear orde pertama}}.
Sistem itu linear karena tidak satu pun variabel dependen atau turunannya
muncul dalam fungsi nonlinear atau dengan pangkat
lebih tinggi daripada satu ($y_1$, $y_2$, $y_1'$, $y_2'$, konstanta, dan fungsi dari $x$
boleh muncul, tetapi bukan $y_1y_2$ atau ${(y_2')}^2$ atau $y_1^3$).
Contoh lain yang lebih
rumit dari sistem linear (orde kedua) ialah
\begin{equation*}
\begin{aligned}
y_1'' &= e^x y_1' + x^2 y_1 + 5 y_2 + \sin(x), \\
y_2'' &= x y_1'-y_2' + 2 y_1 + \cos(x).
\end{aligned}
\end{equation*}

\subsection{Penerapan}

Kita meninjau beberapa penerapan sederhana sistem dan cara menyusun
persamaannya.

\begin{example} \label{sintro:closedbrine-example}
Pertama, kita meninjau tangki garam dan larutan garam, tetapi kali ini air mengalir
dari satu tangki ke tangki lain dan kembali lagi.
Bayangkan kita memiliki dua tangki, masing-masing berisi larutan garam bervolume $V$ liter.
Jumlah garam dalam tangki pertama adalah $x_1$ gram, dan jumlah garam
dalam tangki kedua adalah $x_2$ gram.
Misalkan $t$ menyatakan waktu---variabel independen.
Cairan dalam setiap tangki terus-menerus
diaduk hingga tercampur sempurna dan
mengalir atau dipompa dengan laju $r$ liter per detik dari setiap tangki ke tangki lainnya.
Lihat \figurevref{sintro:closedbrine}.

\begin{myfig}
\capstart
\inputpdft{lin-tank-sys}{%
Diagram dua tangki berisi cairan.  Zat dalam tangki pertama
diberi label x sub 1 untuk menyatakan jumlah garam
dan zat dalam tangki kedua diberi label x sub 2.
Kedua tangki memiliki label Vol. sama dengan V untuk menyatakan bahwa volumenya
sama dan bernilai V dalam kedua tangki.  Terdapat dua pipa,
masing-masing membawa cairan dari satu tangki ke tangki lainnya,
dan keduanya diberi tanda r yang menyatakan aliran.}
\caption{Sistem tertutup dua tangki larutan garam.\label{sintro:closedbrine}}
\end{myfig}

Laju perubahan $x_1$,
yakni $x_1'$, adalah
laju garam yang masuk dikurangi laju garam yang keluar.
Kerapatan garam dalam tangki 2 adalah $\nicefrac{x_2}{V}$,
sehingga laju garam yang masuk ke tangki 1 adalah
$\nicefrac{x_2}{V}$ dikalikan $r$.
Kerapatan garam dalam tangki 1 adalah $\nicefrac{x_1}{V}$,
sehingga laju garam yang keluar dari tangki 1 adalah
$\nicefrac{x_1}{V}$ dikalikan $r$.
Dengan kata lain,
\begin{equation*}
x_1' = \frac{x_2}{V} r - \frac{x_1}{V} r =
\frac{r}{V} x_2 - \frac{r}{V} x_1  = \frac{r}{V} (x_2-x_1).
\end{equation*}
Demikian pula, untuk menentukan laju $x_2'$, peran $x_1$ dan $x_2$
dibalik.  Secara keseluruhan, sistem PDB untuk masalah ini adalah
\begin{equation*}
\begin{aligned}
x_1' & = \frac{r}{V} (x_2-x_1), \\
x_2' & = \frac{r}{V} (x_1-x_2).
\end{aligned}
\end{equation*}
Dalam sistem ini, kita tidak dapat menyelesaikan $x_1$ atau $x_2$ secara terpisah.  Kita harus
menyelesaikan $x_1$ dan $x_2$ sekaligus,
yang jelas secara intuitif---jumlah garam dalam satu tangki
memengaruhi jumlah garam dalam tangki lainnya.
Kita tidak dapat mengetahui $x_1$ sebelum mengetahui $x_2$, dan sebaliknya.

Kita belum mengetahui cara menemukan semua solusi, tetapi
secara intuitif setidaknya kita dapat menemukan beberapa solusi.  Andaikan kita
mengetahui bahwa pada awalnya kedua tangki memiliki jumlah garam yang sama.  Artinya,
kita memiliki syarat awal seperti $x_1(0)=x_2(0) = C$.  Maka jelas bahwa
jumlah garam yang masuk dan keluar dari setiap tangki sama, sehingga jumlah tersebut
tidak berubah.  Dengan kata lain, $x_1 = C$ dan $x_2 = C$ (fungsi-fungsi
konstan) merupakan solusi: $x_1' = x_2' = 0$, dan
$\frac{r}{V}(x_2-x_1) = \frac{r}{V}(x_1-x_2) = 0$, sehingga persamaan-persamaannya dipenuhi.

Mari kita pikirkan penyusunan tersebut sedikit lebih jauh tanpa menyelesaikannya.  Andaikan
syarat awalnya adalah $x_1(0) = A$ dan $x_2(0) = B$, untuk dua konstanta
berbeda $A$ dan $B$.  Karena tidak ada garam yang masuk atau keluar dari sistem tertutup ini,
jumlah total garam konstan.  Artinya, $x_1+x_2$ konstan,
sehingga nilainya sama dengan $A+B$.
Secara intuitif, jika $A > B$, lebih banyak garam akan mengalir keluar dari tangki
satu daripada masuk ke dalamnya.  Setelah waktu yang lama, kita kemudian mengharapkan
jumlah garam dalam setiap tangki menjadi sama. Dengan kata lain,
solusi $x_1$ dan $x_2$ keduanya seharusnya
menuju $\frac{A+B}{2}$ ketika
$t$ menuju $\infty$.
Setelah Anda mengetahui cara menyelesaikan sistem,
Anda dapat memeriksa
bahwa hal ini memang benar.
\end{example}

\begin{example} \label{sintro:carts-example}
Mari kita lihat contoh orde kedua.
Kita kembali ke susunan massa dan pegas, tetapi kali ini kita
meninjau dua massa.

\begin{mywrapfigsimp}{2.0in}{2.3in}
\noindent
\inputpdft{cartsfig}{%
Diagram dua kereta beroda yang dihubungkan oleh sebuah pegas.
Kereta-kereta tersebut diberi label m sub 1 dan m sub 2 yang menyatakan massanya.
Pegas diberi label k yang menyatakan konstanta pegas.
Terdapat tanda garis vertikal dan panah horizontal ke kanan di bawah
setiap kereta, masing-masing berlabel x sub 1 dan x sub 2, yang menyatakan bahwa
keadaan yang ditampilkan adalah posisi setimbang dan x sub 1 serta x sub 2 memberikan
simpangan setiap kereta dari posisi setimbangnya.}
\end{mywrapfigsimp}
Tinjau satu pegas dengan konstanta $k$ dan dua massa $m_1$
dan $m_2$.  Bayangkan massa-massa tersebut sebagai kereta pada
lintasan lurus tanpa gesekan.  Misalkan $x_1$ adalah simpangan kereta pertama
dan $x_2$ adalah simpangan kereta kedua.
Artinya, kita menempatkan kedua
kereta di suatu tempat tanpa tegangan pada pegas,
kita menandai posisi
kereta pertama dan kedua,
dan kita menyebutnya posisi nol.
Kemudian $x_1$ mengukur seberapa jauh kereta pertama dari posisi nolnya,
dan $x_2$ mengukur seberapa jauh kereta kedua dari posisi nolnya.
Gaya yang dikerahkan pegas pada kereta pertama adalah
$k(x_2-x_1)$
karena $x_2-x_1$ adalah besarnya peregangan (atau pemampatan) pegas dari
posisi setimbang.
Gaya yang dikerahkan pada kereta kedua berlawanan,
jadi besarnya sama dengan tanda negatif.
\myindex{Hukum kedua Newton} menyatakan bahwa
gaya sama dengan massa dikalikan percepatan, yakni,
\begin{equation*}
\begin{aligned}
m_1 x_1'' & = k(x_2-x_1) , \\
m_2 x_2'' & = - k(x_2-x_1) .
\end{aligned}
\end{equation*}

Sekali lagi, kita tidak dapat menyelesaikan variabel $x_1$ atau $x_2$
satu demi satu.  Ke mana kereta pertama bergerak
bergantung tepat pada gerak kereta kedua dan sebaliknya.
\end{example}

\subsection{Mengubah ke sistem orde pertama}

Sampai taraf tertentu,
kita hanya perlu mampu menyelesaikan sistem orde pertama.
Tinjau persamaan diferensial orde $n^{\text{th}}$
\begin{equation*}
y^{(n)} = F(y^{(n-1)},\ldots,y',y,x) .
\end{equation*}
Kita definisikan variabel baru $u_1, u_2, \ldots, u_n$ dan tuliskan sistem
\begin{align*}
u_1' & = u_2 , \\
u_2' & = u_3 , \\
& ~\, \vdots \\
u_{n-1}' & = u_n , \\
u_n' & = F(u_n,u_{n-1},\ldots,u_2,u_1,x) .
\end{align*}
Kita selesaikan
sistem ini untuk $u_1$, $u_2$, \ldots, $u_n$.  Setelah kita memperoleh solusi
untuk $u$,
kita dapat membuang $u_2$ hingga $u_n$ dan menetapkan $y = u_1$.
$y$ ini menyelesaikan persamaan semula.

\begin{example}
Ambil $x''' = 2x''+ 8x' + x + t$.  Dengan menetapkan $u_1 = x$, $u_2 = x'$, $u_3
= x''$, kita memperoleh sistem:
\begin{equation*}
u_1' = u_2, \qquad u_2' = u_3, \qquad u_3' = 2u_3 + 8u_2 + u_1 + t .
\end{equation*}
Perhatikan mengapa $x=u_1$ menyelesaikan persamaan semula: Dua persamaan
pertama dalam sistem memberikan
$u_3' = u_2'' = u_1''' = x'''$.
Jika kita substitusikan $u_3' = x'''$,
$u_1 = x$, $u_2 = x'$, dan $u_3
= x''$ ke dalam
persamaan ketiga sistem tersebut, kita
memperoleh kembali persamaan orde ketiga yang menjadi titik awal kita.
\end{example}

Gagasan yang sama berlaku untuk sistem persamaan diferensial
orde tinggi.
Suatu sistem $k$ persamaan diferensial dalam $k$
variabel tak diketahui, yang semuanya berorde $n$, dapat diubah menjadi
sistem orde pertama dengan $n \times k$
persamaan dan $n \times k$ variabel tak diketahui.

\begin{example}
Tinjau sistem dari contoh kereta, \exampleref{sintro:carts-example}:
\begin{equation*}
m_1 x_1''  = k(x_2-x_1), \qquad m_2 x_2'' = - k(x_2-x_1) .
\end{equation*}
Misalkan $u_1 = x_1$, $u_2 = x_1'$,
$u_3 = x_2$, $u_4 = x_2'$.  Sistem orde kedua tersebut menjadi
sistem orde pertama
\begin{equation*}
u_1' = u_2, \qquad
m_1 u_2'  = k(u_3-u_1), \qquad
u_3' = u_4, \qquad
m_2 u_4' = - k(u_3-u_1) .
\end{equation*}
\end{example}

\begin{example}
Gagasan tersebut juga berlaku dalam arah sebaliknya.
Andaikan kita ingin menyelesaikan
sistem
\begin{equation*}
x' = 2y-x , \qquad
y' = x,
\end{equation*}
untuk syarat
awal $x(0) = 1$, $y(0) =0$.
Misalkan variabel independennya adalah $t$.

Jika kita mendiferensialkan persamaan kedua, kita memperoleh
$y''=x'$.  Kita mengetahui $x'$ dalam $x$ dan $y$, serta
mengetahui bahwa $x=y'$.  Jadi,
\begin{equation*}
y'' = x' = 2y-x = 2y-y' .
\end{equation*}
Sekarang kita memiliki persamaan $y''+y'-2y = 0$.  Kita mengetahui cara menyelesaikan
persamaan ini dan memperoleh $y = C_1 e^{-2t} + C_2 e^t$.  Setelah memperoleh $y$,
kita menggunakan persamaan $y' = x$ untuk memperoleh $x$.
\begin{equation*}
x = y' = -2 C_1 e^{-2t} + C_2 e^t .
\end{equation*}
Kita selesaikan syarat awal $1 = x(0) = -2 C_1 + C_2$
dan $0 = y(0) = C_1 + C_2$.  Jadi, $C_1 = -C_2$ dan $1 = 3C_2$.
Dengan demikian $C_1 = \nicefrac{-1}{3}$ dan $C_2 = \nicefrac{1}{3}$.  Solusi kita adalah
\begin{equation*}
x = \frac{2e^{-2t} + e^t}{3} ,\qquad
y = \frac{-e^{-2t} + e^t}{3} .
\end{equation*}
\end{example}

\begin{exercise}
Substitusikan dan periksa bahwa bentuk tersebut memang merupakan solusinya.
\end{exercise}

Beralih antara sistem
dan persamaan orde tinggi
berguna untuk alasan lain.  Sebagai contoh, perangkat lunak untuk menyelesaikan PDB secara numerik
(pendekatan) umumnya dirancang bagi sistem orde pertama.  Untuk menggunakannya,
ambil PDB apa pun yang ingin Anda selesaikan dan ubahlah menjadi
sistem orde pertama.  Tidak terlalu sulit
untuk mengadaptasi kode komputer bagi metode Euler atau Runge--Kutta untuk
persamaan orde pertama agar
dapat menangani sistem orde pertama.
Kita cukup memperlakukan variabel dependen bukan sebagai
bilangan, melainkan sebagai vektor.  Dalam banyak bahasa pemrograman matematika, hampir
tidak ada perbedaan sintaksis.

%Sebenarnya, inilah yang dilakukan IODE ketika Anda memintanya menyelesaikan persamaan orde kedua
%secara numerik dalam Proyek IODE III jika Anda telah mengerjakan proyek tersebut.

\subsection{Sistem otonom dan medan vektor}

Sistem yang persamaannya tidak bergantung pada variabel independen
disebut \emph{\myindex{sistem otonom}}.  Sebagai contoh,
sistem $x'=2y-x$, $y'=x$ bersifat otonom karena
variabel independennya,
misalnya $t$, tidak muncul dalam persamaan.

Untuk sistem otonom, kita dapat menggambar apa yang disebut
\emph{\myindex{medan arah}} atau \emph{\myindex{medan vektor}},
suatu plot yang serupa dengan medan kemiringan, tetapi
alih-alih memberikan kemiringan di setiap titik, kita memberikan arah (dan
besar).  Contoh sebelumnya, $x' = 2y-x$, $y' = x$, menyatakan
bahwa di titik $(x,y)$ arah yang harus kita tempuh untuk memenuhi
persamaan tersebut harus searah dengan vektor $( 2y-x, x )$
dengan kecepatan yang sama dengan besar vektor ini.  Jadi kita gambar
vektor $(2y-x,x)$ di titik $(x,y)$ dan melakukannya untuk
banyak titik pada bidang-$xy$.
Sebagai contoh, di titik $(1,2)$, kita menggambar vektor
$\bigl(2(2)-1,1\bigr) = (3,1)$,
sebuah vektor yang menunjuk ke kanan dan sedikit ke atas,
sedangkan di titik $(2,1)$ kita menggambar vektor $\bigl(2(1)-2,2\bigr) = (0,2)$,
sebuah vektor yang menunjuk lurus ke atas.
Ketika menggambar vektor-vektor tersebut, kita akan memperkecil
ukurannya agar banyak vektor dapat dimuat pada medan arah yang sama.
Jika kita menggambar panah sesuai ukuran sebenarnya, diagram akan menjadi
kacau setelah kita menggambar lebih dari beberapa panah.
Jadi kita skalakan semuanya agar bahkan panah terpanjang pun tidak mengganggu
yang lain.
Kita terutama
tertarik pada arah dan ukuran relatifnya.  Lihat
\figurevref{sintro-vectorfield:fig}.
Diagram yang kita gambar dalam
\sectionref{auteq:section}
untuk persamaan otonom dalam satu dimensi
serupa,
tetapi perhatikan betapa jauh lebih rumit keadaannya ketika kita
hanya menambahkan satu dimensi lagi.

Kita dapat menggambar lintasan solusi di bidang.  Andaikan
solusi diberikan oleh $x = f(t)$, $y=g(t)$.  Kita pilih suatu interval
$t$ (misalnya $0 \leq t \leq 2$ untuk contoh kita) dan memplot semua titik
$\bigl(f(t),g(t)\bigr)$ untuk $t$ dalam rentang yang dipilih.  Gambar yang dihasilkan
disebut \emph{\myindex{potret fase}}
(atau \myindex{potret bidang fase}).
Kurva tertentu yang diperoleh
disebut \emph{\myindex{trajektori}} atau \emph{\myindex{kurva solusi}}.
Lihat contoh plot dalam \figurevref{sintro-vectorfield-sol:fig}.
Dalam gambar tersebut, solusi berawal di $(1,0)$ dan bergerak sepanjang medan vektor
sejauh 2 satuan $t$.  Kita telah menyelesaikan sistem ini secara eksak, sehingga
kita menghitung $x(2)$ dan $y(2)$ untuk memperoleh
$x(2) \approx 2.475$ dan $y(2) \approx 2.457$.  Titik ini bersesuaian
dengan ujung kanan atas kurva solusi yang diplot dalam gambar.

\begin{myfig}
\parbox[t]{3.0in}{
 \capstart
 \diffyincludegraphics{width=3.0in}{width=4.5in}{sintro-vectorfield}{%
Diagram suatu medan arah, yaitu kisi panah-panah kecil.
Panah di bagian atas diagram umumnya menunjuk ke kanan, sedangkan panah
di sisi kanan umumnya menunjuk ke atas.  Sepanjang diagonal x sama dengan y,
panah-panah menunjuk menjauhi titik asal dan
panah di tempat lain cenderung menunjuk menuju garis ini lalu perlahan berubah
untuk menunjuk ke arah yang sama dengan panah pada diagonal.}
 \caption{Medan arah untuk $x' = 2y-x$, $y' = x$.%
 \label{sintro-vectorfield:fig}}
}
\quad
\parbox[t]{3.0in}{
 \capstart
 \diffyincludegraphics{width=3.0in}{width=4.5in}{sintro-vectorfield-sol}{%
Diagram medan arah yang sama seperti gambar sebelumnya, tetapi kini dengan
solusi yang berawal di (1,0), mengalir sejenak menuju diagonal
x sama dengan y, lalu berbelok tajam ke kanan dan kemudian mengalir di sepanjang diagonal menuju
sudut kanan atas, serta berakhir sebelum mencapai sudut gambar.}
 \caption{Medan arah untuk $x' = 2y-x$, $y' = x$ dengan
 trajektori solusi yang berawal di $(1,0)$
 untuk $0 \leq t \leq 2$.%
 \label{sintro-vectorfield-sol:fig}}
}
\end{myfig}

Kita dapat menggambar potret fase dan trajektori pada bidang-$xy$
meskipun sistemnya tidak otonom.  Namun, dalam kasus ini kita tidak dapat menggambar
medan arah karena medan tersebut berubah seiring perubahan $t$.  Untuk
setiap $t$, kita akan memperoleh medan arah yang berbeda.

\subsection{Teorema Picard}

Sebelum melangkah lebih jauh, kita sebutkan bahwa teorema Picard tentang
keberadaan dan ketunggalan juga berlaku bagi sistem PDB.
Mari kita nyatakan kembali teorema ini dalam konteks tersebut.
Sistem umum orde pertama
berbentuk
\begin{equation} \label{eq:general-system}
\begin{aligned}
x_1' & = F_1(x_1,x_2,\ldots,x_n,t) , \\
x_2' & = F_2(x_1,x_2,\ldots,x_n,t) , \\
& \vdots \\
x_n' & = F_n(x_1,x_2,\ldots,x_n,t) .
\end{aligned}
\end{equation}

\begin{theorem}[Teorema Picard tentang keberadaan dan ketunggalan untuk sistem]%
\label{sys:picardthm}%
\index{keberadaan dan ketunggalan untuk sistem}\index{teorema Picard}
Jika untuk setiap $j=1,2,\ldots,n$ dan setiap
$k = 1,2,\ldots,n$
setiap $F_j$ kontinu
dan turunan
$\frac{\partial F_j}{\partial x_k}$ ada serta
kontinu di dekat suatu $(x_1^0,x_2^0,\ldots,x_n^0,t^0)$, maka solusi
\eqref{eq:general-system}
yang memenuhi syarat awal
$x_1(t^0) = x_1^0$,
$x_2(t^0) = x_2^0$, \ldots,
$x_n(t^0) = x_n^0$
ada (setidaknya untuk $t$ dalam suatu interval kecil) dan tunggal.
\end{theorem}

Artinya, terdapat solusi tunggal untuk setiap syarat awal
asalkan sistemnya wajar (setiap $F_j$ dan turunan parsialnya
terhadap variabel $x$ kontinu).  Seperti pada persamaan tunggal,
kita mungkin tidak memiliki solusi untuk semua waktu $t$, tetapi solusi
dijamin ada setidaknya selama suatu
selang waktu singkat.

Karena kita dapat mengubah setiap PDB orde $n^{\text{th}}$ menjadi
sistem orde pertama,
teorema ini juga memberikan
keberadaan dan ketunggalan solusi bagi persamaan orde tinggi.

\subsection{Latihan}

\begin{exercise}
Carilah solusi umum $x_1' = x_2 - x_1 + t$, $x_2' = x_2$.
\end{exercise}

\begin{exercise}
Carilah solusi umum $x_1' = 3 x_1 - x_2 + e^t$, $x_2' = x_1$.
\end{exercise}

\begin{exercise}
Tuliskan $ay'' + by' + cy = f(x)$
sebagai sistem PDB orde pertama.
\end{exercise}

\begin{exercise}
Tuliskan $x'' + y^2 y' - x^3 = \sin(t)$,
$y'' + {(x'+y')}^2 -x = 0$ sebagai sistem PDB orde pertama.
\end{exercise}

\begin{exercise}
Andaikan dua massa pada kereta di permukaan tanpa gesekan memiliki
simpangan $x_1$ dan $x_2$ seperti dalam \examplevref{sintro:carts-example}.
Andaikan sebuah roket memberikan gaya $F$ dalam arah positif pada kereta
$x_1$.  Susunlah sistem persamaannya.
\end{exercise}

\begin{exercise}
Andaikan tangki-tangki tersebut seperti dalam
\examplevref{sintro:closedbrine-example}, keduanya mula-mula bervolume $V$,
tetapi kini laju aliran dari tangki 1 ke tangki 2 adalah $r_1$, dan
laju aliran dari tangki 2 ke tangki satu adalah $r_2$.
Perhatikan bahwa volumenya kini tidak konstan.  Susunlah sistem persamaannya.
\end{exercise}

\setcounter{exercise}{100}

\begin{exercise}
Carilah solusi umum untuk $y_1' = 3 y_1$, $y_2' = y_1 + y_2$,
$y_3' = y_1 + y_3$.
\end{exercise}
\exsol{%
$y_1 = C_1 e^{3x}$,
$y_2 = C_2 e^x+ \frac{C_1}{2} e^{3 x}$,
$y_3 = C_3 e^x+ \frac{C_1}{2} e^{3 x}$
}

\begin{exercise}
Selesaikan $y'=2x$, $x'=x+y$, $x(0)=1$, $y(0)=3$.
\end{exercise}
\exsol{%
$x=\frac{5}{3} e^{2t} - \frac{2}{3} e^{-t}$,
$y=\frac{5}{3} e^{2t} + \frac{4}{3} e^{-t}$
}

\begin{exercise}
Tuliskan $x''' = x+t$ sebagai sistem orde pertama.
\end{exercise}
\exsol{%
$u_1' = u_2$,
$u_2' = u_3$,
$u_3' = u_1+t$, (di sini $u_1=x$)
}

\begin{exercise}
Tuliskan $y_1'' + y_1 + y_2 = t$,
$y_2'' + y_1 - y_2 = t^2$ sebagai sistem orde pertama.
\end{exercise}
\exsol{%
$y_1' = y_3$,
$y_2' = y_4$,
%$y_3' + y_1 + y_2 = t$,
$y_3' = - y_1 - y_2 + t$,
%$y_4' + y_1 - y_2 = t^2$
$y_4' = - y_1 + y_2 + t^2$
}

\begin{exercise}
Andaikan dua massa pada kereta di permukaan tanpa gesekan memiliki
simpangan $x_1$ dan $x_2$ seperti dalam \examplevref{sintro:carts-example}.
Andaikan simpangan awalnya $x_1(0)=x_2(0)=0$, dan kecepatan awalnya $x_1'(0) = x_2'(0) = a$ untuk suatu bilangan $a$.
Gunakan intuisi Anda
untuk menyelesaikan sistem tersebut dan jelaskan penalaran Anda.
\end{exercise}
\exsol{%
$x_1 = x_2 = at$.  Penjelasan intuisinya diserahkan kepada pembaca.
}

\begin{exercise}
Andaikan tangki-tangki tersebut seperti dalam
\examplevref{sintro:closedbrine-example}, kecuali bahwa air bersih mengalir masuk
dengan laju $s$ liter per detik ke tangki 1, dan larutan garam mengalir keluar dari tangki 2
menuju saluran pembuangan, juga dengan laju $s$ liter per detik.
Laju aliran dari tangki 1 ke tangki 2 tetap $r$, tetapi laju aliran
dari tangki 2 kembali ke tangki 1 adalah $r-s$ (andaikan $r > s$).
\begin{tasks}
\task Gambarlah diagramnya.
\task Susunlah sistem persamaannya.
\task Secara intuitif, apa yang terjadi ketika $t$ menuju tak hingga? Jelaskan.
\end{tasks}
\end{exercise}
\exsol{%
a) Diserahkan kepada pembaca.
\quad b)
$x_1' = \frac{r-s}{V} x_2- \frac{r}{V}x_1$,
$x_2' = \frac{r}{V} (x_1- x_2)$.
\quad c) Ketika $t$ menuju tak hingga, $x_1$ dan $x_2$ keduanya menuju nol;
penjelasannya diserahkan kepada pembaca.
}

%%%%%%%%%%%%%%%%%%%%%%%%%%%%%%%%%%%%%%%%%%%%%%%%%%%%%%%%%%%%%%%%%%%%%%%%%%%%%%

\sectionnewpage
\section{Matriks dan sistem linear} \label{sec:matrix}

%mbxINTROSUBSECTION

\sectionnotes{1,5 perkuliahan\EPref{,
bagian pertama \S5.1 dalam \cite{EP}}\BDref{,
\S7.2 dan \S7.3 dalam \cite{BD}}, lihat pula \appendixref{linalg:appendix}}

\subsection{Matriks dan vektor}

Sebelum mulai membahas sistem linear PDB, kita perlu
membahas matriks, jadi mari kita tinjau secara singkat.  Suatu \emph{\myindex{matriks}}
adalah susunan bilangan berukuran $m
\times n$ ($m$ baris dan $n$ kolom).  Sebagai contoh, kita menyatakan
matriks $3 \times 5$ sebagai berikut
\begin{equation*}
A =
\begin{bmatrix}
a_{11} & a_{12} & a_{13} & a_{14} & a_{15} \\
a_{21} & a_{22} & a_{23} & a_{24} & a_{25} \\
a_{31} & a_{32} & a_{33} & a_{34} & a_{35}
\end{bmatrix} .
\end{equation*}
Bilangan $a_{ij}$ disebut \emph{elemen}\index{elemen matriks}
atau \emph{entri}\index{entri matriks}.
Kita mengatakan suatu matriks \emph{persegi}\index{matriks persegi}
jika $m=n$, yaitu matriks tersebut memiliki jumlah baris dan kolom yang sama.

Dalam konteks matriks,
dengan \emph{\myindex{vektor}}, biasanya yang kita maksud ialah
\emph{\myindex{vektor kolom}}, yaitu matriks $m \times 1$.
Jika yang kita maksud adalah \emph{\myindex{vektor baris}},
kita akan menyatakannya secara eksplisit (vektor baris adalah matriks $1 \times n$).
Biasanya kita menyatakan
matriks dengan huruf kapital dan vektor dengan huruf kecil yang diberi
panah seperti $\vec{x}$ atau $\vec{b}$.  Kita menuliskan $\vec{0}$ untuk vektor yang semua entrinya nol.

Kita definisikan beberapa operasi pada matriks.  Kita
menginginkan matriks $1 \times 1$ benar-benar berlaku seperti bilangan, sehingga operasi kita
harus sesuai dengan sudut pandang ini.
Pertama, kita dapat mengalikan\index{perkalian skalar} suatu matriks dengan
\emph{\myindex{skalar}} (suatu bilangan).
Kita cukup mengalikan setiap entri dalam matriks dengan skalar tersebut.  Sebagai contoh,
\begin{equation*}
2
\begin{bmatrix}
1 & 2 & 3 \\
4 & 5 & 6
\end{bmatrix} =
\begin{bmatrix}
2 & 4 & 6 \\
8 & 10 & 12
\end{bmatrix} .
\end{equation*}
Penjumlahan matriks\index{penjumlahan matriks} juga sederhana.
Kita menjumlahkan matriks elemen demi elemen.
Sebagai contoh,
\begin{equation*}
\begin{bmatrix}
1 & 2 & 3 \\
4 & 5 & 6
\end{bmatrix} +
\begin{bmatrix}
1 & 1 & -1 \\
0 & 2 & 4
\end{bmatrix}
=
\begin{bmatrix}
2 & 3 & 2 \\
4 & 7 & 10
\end{bmatrix} .
\end{equation*}
Jika ukurannya tidak sesuai, penjumlahan tidak terdefinisi.
Jika kita menyatakan matriks dengan semua entri nol sebagai 0,
$c$, $d$ sebagai skalar, dan $A$, $B$, $C$ sebagai matriks, kita
memiliki aturan-aturan yang sudah dikenal berikut:
\begin{align*}
A + 0 & = A = 0 + A , \\
A + B & = B + A , \\
(A + B) + C & = A + (B + C) , \\
c(A+B) & = cA+cB, \\
(c+d)A & = cA + dA.
\end{align*}

Operasi lain yang berguna untuk matriks adalah apa yang disebut
\emph{\myindex{transpos}}.  Operasi ini hanya menukar baris dan kolom suatu
matriks.
Transpos $A$ dinyatakan dengan $A^T$.  Contoh:
\begin{equation*}
\begin{bmatrix}
1 & 2 & 3 \\
4 & 5 & 6
\end{bmatrix}^T =
\begin{bmatrix}
1 & 4 \\
2 & 5 \\
3 & 6
\end{bmatrix}
\end{equation*}

\subsection{Perkalian matriks}

Perkalian matriks sedikit lebih rumit.
Pertama, kita definisikan apa yang disebut
\emph{\myindex{hasil kali titik}} (atau \emph{\myindex{hasil kali dalam}}) dari dua vektor.
Biasanya ini berupa vektor baris yang dikalikan
dengan vektor kolom berukuran sama.  Untuk hasil kali titik,
kita mengalikan
setiap pasangan entri dari vektor pertama dan kedua, lalu menjumlahkan
hasil-hasil kali tersebut.  Hasilnya adalah satu bilangan.
Sebagai contoh,
\begin{equation*}
\begin{bmatrix}
a_1 & a_2 & a_3
\end{bmatrix}
\cdot
\begin{bmatrix}
b_1 \\
b_2 \\
b_3
\end{bmatrix}
= a_1 b_1 + a_2 b_2 + a_3 b_3 .
\end{equation*}
Hal yang sama berlaku untuk vektor yang lebih besar (atau lebih kecil).

Dengan hasil kali titik, kita mendefinisikan
\emph{\myindex{hasil kali matriks}}.
Kita menyatakan dengan $\operatorname{row}_i(A)$ baris $i^{\text{th}}$
dari $A$ dan dengan
$\operatorname{column}_j(A)$ kolom $j^{\text{th}}$ dari $A$.
Untuk matriks $m \times n$ $A$ dan matriks $n \times p$ $B$,
kita dapat mendefinisikan hasil kali $AB$.  Kita tetapkan $AB$ sebagai matriks $m \times p$
yang entri $ij^{\text{th}}$-nya adalah hasil kali titik
\begin{equation*}
\operatorname{row}_i(A) \cdot
\operatorname{column}_j(B) .
\end{equation*}
Perhatikan bagaimana ukurannya bersesuaian: $m \times n$ dikalikan $n \times p$ menjadi
$m \times p$.  Contoh:
\begin{multline*}
\begin{bmatrix}
1 & 2 & 3 \\
4 & 5 & 6
\end{bmatrix}
\begin{bmatrix}
1 & 0 & -1 \\
1 & 1 & 1 \\
1 & 0 & 0
\end{bmatrix}
= \\ =
\begin{bmatrix}
1\cdot 1 + 2\cdot 1 + 3 \cdot 1 & ~~
1\cdot 0 + 2\cdot 1 + 3 \cdot 0 & ~~
1\cdot (-1) + 2\cdot 1 + 3 \cdot 0 \\
4\cdot 1 + 5\cdot 1 + 6 \cdot 1 & ~~
4\cdot 0 + 5\cdot 1 + 6 \cdot 0 & ~~
4\cdot (-1) + 5\cdot 1 + 6 \cdot 0
\end{bmatrix}
=
\begin{bmatrix}
6 & 2 & 1 \\
15 & 5 & 1
\end{bmatrix}
\end{multline*}

Untuk perkalian, kita menginginkan analog dari bilangan 1.  Analog ini adalah
apa yang disebut \emph{\myindex{matriks identitas}}.
Matriks identitas adalah matriks persegi dengan 1 pada
diagonal dan nol di semua tempat lain.  Matriks ini biasanya dinyatakan dengan $I$.
Untuk setiap ukuran, kita memiliki matriks identitas yang berbeda.
Karena itu, terkadang kita menyatakan
ukurannya sebagai subskrip.  Sebagai contoh, $I_3$ adalah matriks identitas
$3 \times 3$
\begin{equation*}
I = I_3 =
\begin{bmatrix}
1 & 0 & 0 \\
0 & 1 & 0 \\
0 & 0 & 1
\end{bmatrix} .
\end{equation*}

\pagebreak[2]
Kita memiliki aturan-aturan berikut untuk perkalian matriks.  Andaikan
$A$, $B$, $C$ adalah matriks dengan ukuran yang sesuai sehingga pernyataan berikut
terdefinisi.  Misalkan $\alpha$ menyatakan suatu skalar (bilangan).
Maka,
\begin{align*}
A(BC) & = (AB)C, \\
A(B+C) & = AB + AC, \\
(B+C)A & = BA + CA, \\
\alpha(AB) & = (\alpha A)B = A(\alpha B), \\
IA & = A = AI .
\end{align*}

\pagebreak[2]
Beberapa peringatan perlu disampaikan.
\begin{enumerate}[(i)]
\item Secara umum $AB \neq BA$ (kadang-kadang mungkin benar secara kebetulan).  Artinya,
matriks tidak \myindex{komutatif}.
Sebagai contoh, ambil
$A = \left[ \begin{smallmatrix} 1 & 1 \\ 1 & 1 \end{smallmatrix} \right]$
dan
$B = \left[ \begin{smallmatrix} 1 & 0 \\ 0 & 2 \end{smallmatrix} \right]$.
\item $AB = AC$ tidak selalu mengakibatkan $B=C$, bahkan jika $A$ bukan 0.
\item $AB = 0$ tidak selalu berarti $A=0$ atau $B=0$.
Cobalah, misalnya,
$A = B = \left[ \begin{smallmatrix} 0 & 1 \\ 0 & 0 \end{smallmatrix}
\right]$.
\end{enumerate}

Agar dua butir terakhir berlaku, kita perlu \myquote{membagi} dengan
suatu matriks.  Di sinilah \emph{\myindex{invers matriks}} berperan.
Andaikan $A$ dan $B$ adalah matriks $n \times n$ sedemikian sehingga
\begin{equation*}
AB = I = BA .
\end{equation*}
Maka kita menyebut $B$ invers dari $A$ dan menyatakan $B$ dengan $A^{-1}$.
Jika invers $A$ ada, kita menyebut $A$
\emph{dapat dibalik\index{matriks dapat dibalik}}.
Jika $A$ tidak dapat dibalik, terkadang kita mengatakan $A$
\emph{singular\index{matriks singular}}.

Jika $A$ dapat dibalik, maka $AB = AC$ memang mengakibatkan
$B = C$ (khususnya, invers $A$ tunggal).
Kita cukup mengalikan kedua ruas dengan $A^{-1}$ (dari kiri) untuk memperoleh
$A^{-1}AB = A^{-1}AC$ atau $IB=IC$ atau $B=C$.
Juga tidak sulit dilihat bahwa ${(A^{-1})}^{-1} = A$.

\subsection{Determinan}

Untuk matriks persegi, kita mendefinisikan besaran berguna yang disebut
\emph{\myindex{determinan}}.  Kita mendefinisikan
determinan matriks $1 \times 1$ sebagai nilai satu-satunya entri.
Untuk matriks $2 \times 2$, kita definisikan
\begin{equation*}
\det \left(
\begin{bmatrix}
a & b \\
c & d
\end{bmatrix}
\right)
\overset{\text{def}}{=}
ad-bc .
\end{equation*}

Sebelum mencoba mendefinisikan
determinan untuk matriks yang lebih besar, mari kita perhatikan
makna determinan.  Tinjau matriks $n \times n$
sebagai pemetaan dari ruang Euklides berdimensi $n$, ${\mathbb{R}}^n$, ke
dirinya sendiri, dengan $\vec{x}$ dipetakan ke $A \vec{x}$.
Secara khusus, matriks $2 \times 2$ $A$ adalah pemetaan dari
bidang ke dirinya sendiri.  Determinan
$A$ adalah faktor perubahan luas objek.
Jika kita mengambil persegi satuan (persegi dengan sisi 1) di bidang, maka
$A$ memetakan persegi tersebut menjadi jajaran genjang berluas $\lvert\det(A)\rvert$.  Tanda
$\det(A)$ menyatakan perubahan orientasi (negatif jika sumbu-sumbunya terbalik).  Sebagai
contoh, misalkan
\begin{equation*}
A =
\begin{bmatrix}
1 & 1 \\
-1 & 1
\end{bmatrix} .
\end{equation*}
Maka $\det(A) = 1+1 = 2$.  Mari kita lihat ke mana persegi (satuan) dengan titik sudut
$(0,0)$, $(1,0)$, $(0,1)$, dan $(1,1)$ dipetakan.  Jelas $(0,0)$ dipetakan
ke $(0,0)$.
\begin{equation*}
\begin{bmatrix}
1 & 1 \\
-1 & 1
\end{bmatrix}
\begin{bmatrix}
1 \\ 0
\end{bmatrix} =
\begin{bmatrix}
1 \\
-1
\end{bmatrix}
,
\qquad
\begin{bmatrix}
1 & 1 \\
-1 & 1
\end{bmatrix}
\begin{bmatrix}
0 \\ 1
\end{bmatrix} =
\begin{bmatrix}
1 \\
1
\end{bmatrix}
,
\qquad
\begin{bmatrix}
1 & 1 \\
-1 & 1
\end{bmatrix}
\begin{bmatrix}
1 \\ 1
\end{bmatrix} =
\begin{bmatrix}
2 \\
0
\end{bmatrix}
.
\end{equation*}
Citra persegi tersebut adalah persegi lain dengan titik sudut $(0,0)$, $(1,-1)$,
$(1,1)$, dan $(2,0)$.  Persegi citra itu
memiliki
sisi sepanjang $\sqrt{2}$ dan karena itu luasnya 2.

Jika Anda mengingat kembali geometri sekolah menengah, Anda mungkin pernah melihat rumus untuk
menghitung luas suatu \myindex{jajaran genjang}
dengan titik sudut $(0,0)$, $(a,c)$, $(b,d)$
dan $(a+b,c+d)$.  Rumus itu tepatnya adalah
\begin{equation*}
\left\lvert \, \det \left(
\begin{bmatrix} a & b \\ c & d \end{bmatrix}
\right) \, \right\lvert.
\end{equation*}
Garis vertikal di atas menyatakan nilai mutlak.
Matriks $\left[ \begin{smallmatrix} a & b \\ c & d \end{smallmatrix}
\right]$
memetakan persegi satuan ke jajaran genjang yang diberikan.

\medskip

Mari kita tinjau determinan untuk matriks yang lebih besar.  Kita definisikan $A_{ij}$ sebagai
matriks $A$ dengan baris $i^{\text{th}}$ dan kolom $j^{\text{th}}$
dihapus.  Untuk menghitung determinan suatu matriks, pilih satu baris, misalnya
baris $i^{\text{th}}$, lalu hitung:
\begin{equation*}
\mybxbg{~~
\det (A) =
\sum_{j=1}^n
{(-1)}^{i+j}
a_{ij} \det (A_{ij}) .
~~}
\end{equation*}
Untuk baris pertama, kita peroleh
\begin{equation*}
\det (A) =
a_{11} \det (A_{11}) -
a_{12} \det (A_{12}) +
a_{13} \det (A_{13}) -
\cdots
\begin{cases}
+ a_{1n} \det (A_{1n}) & \text{jika } n \text{ ganjil,} \\
- a_{1n} \det (A_{1n}) & \text{jika } n \text{ genap.}
\end{cases}
\end{equation*}
Secara bergantian kita menjumlahkan dan mengurangkan determinan submatriks
$A_{ij}$ yang dikalikan dengan $a_{ij}$ untuk $i$ tetap dan semua $j$.
Untuk matriks $3 \times 3$,
dengan memilih baris pertama, kita memperoleh $\det (A) = a_{11} \det (A_{11}) -
a_{12} \det (A_{12}) + a_{13} \det (A_{13})$.  Sebagai contoh,
\begin{equation*}
\begin{split}
\det \left(
\begin{bmatrix}
1 & 2 & 3 \\
4 & 5 & 6 \\
7 & 8 & 9
\end{bmatrix}
\right)
& =
1 \cdot
\det \left(
\begin{bmatrix}
5 & 6 \\
8 & 9
\end{bmatrix}
\right)
-
2 \cdot
\det \left(
\begin{bmatrix}
4 & 6 \\
7 & 9
\end{bmatrix}
\right)
+
3 \cdot
\det \left(
\begin{bmatrix}
4 & 5 \\
7 & 8
\end{bmatrix}
\right) \\
& =
1 (5 \cdot 9 - 6 \cdot 8)
-
2 (4 \cdot 9 - 6 \cdot 7)
+
3 (4 \cdot 8 - 5 \cdot 7)
= 0 .
\end{split}
\end{equation*}

Bilangan ${(-1)}^{i+j}\det(A_{ij})$ disebut
\emph{kofaktor\index{kofaktor}}
matriks tersebut dan
cara menghitung determinan ini disebut
\emph{\myindex{ekspansi kofaktor}}.
Baris mana pun yang Anda pilih selalu memberikan bilangan yang sama.
Determinan juga dapat dihitung dengan melakukan ekspansi
sepanjang kolom (memilih kolom alih-alih baris di atas).
Berlaku bahwa $\det(A) = \det(A^T)$.

Notasi umum untuk determinan adalah sepasang garis
vertikal:
\begin{equation*}
\begin{vmatrix}
a & b \\
c & d
\end{vmatrix}
=
\det \left(
\begin{bmatrix}
a & b \\
c & d
\end{bmatrix}
\right) .
\end{equation*}
Secara pribadi, saya menganggap notasi ini membingungkan karena garis vertikal biasanya
menyatakan besaran positif, sedangkan determinan dapat bernilai negatif.  Pikirkan pula
cara menuliskan nilai mutlak suatu determinan.  Saya tidak akan
menggunakan notasi ini dalam buku ini.

\medskip

Bayangkan determinan memberi tahu Anda penskalaan suatu pemetaan.
Jika $B$ menggandakan ukuran objek geometris dan $A$ melipat-tigakannya,
maka $AB$ (yang menerapkan $B$ pada objek lalu $A$) seharusnya membuat ukuran
bertambah dengan faktor $6 = 3 \cdot 2$.  Hal ini berlaku secara umum:
\begin{equation*}
\det(AB) = \det(A)\det(B) .
\end{equation*}
Sifat ini merupakan salah satu yang paling berguna dan dapat digunakan untuk
benar-benar menghitung determinan.  Akibat yang sangat menarik ialah
maknanya bagi keberadaan invers.
Ambil $A$ dan $B$ sebagai invers satu sama lain, yakni $AB=I$.  Maka
\begin{equation*}
\det(A)\det(B) = \det(AB) = \det(I) = 1 .
\end{equation*}
Baik $\det(A)$ maupun $\det(B)$ tidak dapat bernilai nol.
Sebaliknya, jika $\det(A)$ tidak nol, maka $A$
dapat dibalik.
Kita nyatakan pengamatan ini sebagai teorema
karena sangat penting dalam konteks mata kuliah ini.

\begin{theorem}
Matriks $n \times n$ $A$ dapat dibalik jika dan hanya jika $\det (A) \neq 0$.
\end{theorem}

Sebenarnya, $\det(A^{-1}) \det(A) = 1$ menyatakan bahwa $\det(A^{-1}) =
\frac{1}{\det(A)}$.  Jadi kita bahkan mengetahui determinan $A^{-1}$
sebelum mengetahui cara menghitung $A^{-1}$.

Terdapat rumus sederhana untuk invers matriks $2 \times 2$
\begin{equation*}
\begin{bmatrix}
a & b \\
c & d
\end{bmatrix}^{-1}
=
\frac{1}{ad-bc}
\begin{bmatrix}
d & -b \\
-c & a
\end{bmatrix} .
\end{equation*}
Perhatikan determinan matriks
$[\begin{smallmatrix}a&b\\c&d\end{smallmatrix}]$
pada penyebut pecahan.
Rumus tersebut hanya berlaku jika determinannya taknol; jika tidak, kita
membagi dengan nol.

\subsection{Menyelesaikan sistem linear}

Salah satu penerapan matriks yang akan kita perlukan ialah menyelesaikan sistem
persamaan linear.  Hal ini paling baik ditunjukkan melalui contoh.
Andaikan kita memiliki sistem persamaan linear berikut
\begin{align*}
          2 x_1 +           2 x_2 +           2 x_3 & = 2 , \\
\phantom{9} x_1 + \phantom{9} x_2 +           3 x_3 & = 5 , \\
\phantom{9} x_1 +           4 x_2 + \phantom{9} x_3 & = 10 .
\end{align*}

Tanpa mengubah solusi,
kita dapat menukar persamaan dalam sistem ini,
mengalikan setiap persamaan dengan bilangan taknol, dan
menambahkan kelipatan suatu persamaan ke persamaan lain.
Ternyata operasi-operasi ini selalu cukup untuk menemukan solusi.

Lebih mudah menuliskan sistem tersebut sebagai persamaan matriks.
Sistem di atas dapat
dituliskan sebagai
\begin{equation*}
\begin{bmatrix}
2 & 2 & 2 \\
1 & 1 & 3 \\
1 & 4 & 1
\end{bmatrix}
\begin{bmatrix}
x_1 \\
x_2 \\
x_3
\end{bmatrix}
=
\begin{bmatrix}
2 \\
5 \\
10
\end{bmatrix} .
\end{equation*}
Untuk menyelesaikan sistem tersebut,
kita gabungkan matriks koefisien (matriks pada
ruas kiri persamaan) dengan vektor pada ruas kanan
dan memperoleh apa yang disebut
\emph{\myindex{matriks diperbesar}}:
\begin{equation*}
\left[
\begin{array}{ccc|c}
2 & 2 & 2 & 2 \\
1 & 1 & 3 & 5 \\
1 & 4 & 1 & 10
\end{array}
\right] .
%\qquad
%\text{atau cukup}
%\qquad
%\begin{bmatrix}
%2 & 2 & 2 & 2 \\
%1 & 1 & 3 & 5 \\
%1 & 4 & 1 & 10
%\end{bmatrix} .
\end{equation*}

\pagebreak[2]
Kita menerapkan rangkaian tiga
\emph{\myindex{operasi baris elementer}}\index{operasi baris} berikut.
\begin{enumerate}[(i)]
\item Tukarkan dua baris.
\item Kalikan satu baris dengan bilangan taknol.
\item Tambahkan kelipatan suatu baris ke baris lain.
\end{enumerate}
Kita terus melakukan operasi-operasi ini hingga mencapai keadaan yang
memudahkan kita membaca jawabannya, atau hingga memperoleh kontradiksi yang menunjukkan
tidak adanya solusi, misalnya jika kita memperoleh persamaan seperti $0=1$.

Mari kita kerjakan contoh tersebut.  Pertama, kalikan baris pertama dengan
$\nicefrac{1}{2}$ untuk memperoleh
\begin{equation*}
\left[
\begin{array}{ccc|c}
1 & 1 & 1 & 1 \\
1 & 1 & 3 & 5 \\
1 & 4 & 1 & 10
\end{array}
\right] .
\end{equation*}
Sekarang kurangkan baris pertama dari baris kedua dan ketiga:
\begin{equation*}
\left[
\begin{array}{ccc|c}
1 & 1 & 1 & 1 \\
0 & 0 & 2 & 4 \\
0 & 3 & 0 & 9
\end{array}
\right]
\end{equation*}
Kalikan baris terakhir dengan $\nicefrac{1}{3}$ dan baris kedua dengan $\nicefrac{1}{2}$:
\begin{equation*}
\left[
\begin{array}{ccc|c}
1 & 1 & 1 & 1 \\
0 & 0 & 1 & 2 \\
0 & 1 & 0 & 3
\end{array}
\right]
\end{equation*}
Tukarkan baris 2 dan 3:
\begin{equation*}
\left[
\begin{array}{ccc|c}
1 & 1 & 1 & 1 \\
0 & 1 & 0 & 3 \\
0 & 0 & 1 & 2
\end{array}
\right]
\end{equation*}
Kurangkan baris terakhir dari baris pertama, lalu kurangkan baris kedua
dari baris pertama:
\begin{equation*}
\left[
\begin{array}{ccc|c}
1 & 0 & 0 & -4 \\
0 & 1 & 0 & 3 \\
0 & 0 & 1 & 2
\end{array}
\right]
\end{equation*}
Terakhir, kita pikirkan persamaan yang dinyatakan oleh matriks diperbesar ini:
$x_1 = -4$, $x_2 = 3$, dan $x_3 = 2$.  Kita uji solusi ini dalam sistem
semula dan, voil\`a, solusi tersebut berhasil!

\begin{exercise}
Periksa bahwa solusi di atas benar-benar menyelesaikan persamaan yang diberikan.
\end{exercise}

Kita menuliskan persamaan ini dalam notasi matriks sebagai
\begin{equation*}
A \vec{x} = \vec{b} ,
\end{equation*}
dengan $A$ berupa matriks
$\left[ \begin{smallmatrix}
2 & 2 & 2 \\
1 & 1 & 3 \\
1 & 4 & 1
\end{smallmatrix} \right]$ dan $\vec{b}$ berupa vektor
$\left[ \begin{smallmatrix}
2 \\
5 \\
10
\end{smallmatrix} \right]$.  Solusi juga dapat dihitung melalui
invers,
\begin{equation*}
\vec{x} = A^{-1} A \vec{x} = A^{-1} \vec{b} .
\end{equation*}

\medskip

Ada kemungkinan
solusinya tidak tunggal, atau tidak ada solusi.
Mudah untuk mengetahui jika solusi tidak ada.  Jika selama reduksi
baris Anda memperoleh baris yang semua entrinya kecuali entri terakhir
bernilai nol (entri terakhir dalam suatu baris bersesuaian dengan ruas kanan
persamaan), maka sistem tersebut \emph{tak konsisten\index{sistem tak konsisten}} dan
tidak memiliki solusi.  Sebagai
contoh, untuk sistem 3 persamaan dan 3 variabel tak diketahui, jika Anda menemukan baris
seperti $[\,0 \quad 0 \quad 0 ~\,|\,~ 1\,]$ dalam matriks diperbesar,
Anda mengetahui bahwa sistem tersebut tak konsisten.  Baris itu bersesuaian dengan $0=1$.

\medskip

Pada umumnya, Anda mencoba menggunakan operasi baris hingga syarat-syarat berikut
dipenuhi.  Entri taknol pertama (dari kiri) pada setiap baris disebut
\emph{\myindex{entri utama}}.
\begin{enumerate}[(i)]
\item Entri utama pada setiap baris terletak tepat di sebelah kanan
entri utama baris di atasnya.
\item Setiap baris nol berada di bawah semua baris taknol.
\item Semua entri utama bernilai 1.
\item Semua entri di atas dan di bawah suatu entri utama bernilai nol.
\end{enumerate}
Matriks semacam itu dikatakan berada dalam
\emph{\myindex{bentuk eselon baris tereduksi}}.  Variabel yang
bersesuaian dengan kolom tanpa entri utama disebut
\emph{variabel bebas\index{variabel bebas}}.
Variabel bebas berarti bahwa kita dapat memilih variabel tersebut
sesuka hati, lalu menyelesaikan variabel tak diketahui lainnya.

\begin{example}
Matriks diperbesar berikut berada dalam bentuk eselon baris tereduksi.
\begin{equation*}
\left[
\begin{array}{ccc|c}
1 & 2 & 0 & 3 \\
0 & 0 & 1 & 1 \\
0 & 0 & 0 & 0
\end{array}
\right]
\end{equation*}
Andaikan
variabelnya adalah $x_1$, $x_2$, dan $x_3$.  Maka $x_2$ adalah
variabel bebas, $x_1 = 3 - 2x_2$, dan $x_3 = 1$.

\medskip

Sebaliknya, jika selama proses reduksi baris Anda memperoleh
matriks
\begin{equation*}
\left[
\begin{array}{ccc|c}
1 & 2 & 13 & 3 \\
0 & 0 & 1 & 1 \\
0 & 0 & 0 & 3
\end{array}
\right]
,
\end{equation*}
tidak perlu melangkah lebih jauh.  Baris terakhir bersesuaian dengan
persamaan $0 x_1 + 0 x_2 + 0 x_3 = 3$, yang mustahil.  Jadi,
tidak ada solusi.
\end{example}

\subsection{Menghitung invers}

Jika matriks $A$ persegi dan terdapat solusi tunggal
$\vec{x}$ bagi $A \vec{x} = \vec{b}$ untuk setiap $\vec{b}$ (tidak ada variabel
bebas), maka $A$ dapat dibalik.
Dengan mengalikan kedua ruas dengan $A^{-1}$, Anda dapat melihat bahwa $\vec{x} =
A^{-1} \vec{b}$.  Jadi, menghitung invers berguna jika Anda ingin
menyelesaikan persamaan untuk banyak ruas kanan $\vec{b}$ yang berbeda.

Kita memiliki rumus untuk
invers $2 \times 2$, tetapi menghitung invers
matriks yang lebih besar juga tidak sulit.
Walaupun kita tidak akan terlalu sering menghitung invers secara manual untuk matriks
yang lebih besar daripada $2 \times 2$, mari kita singgung caranya.
Menemukan invers $A$ sebenarnya hanya berarti menyelesaikan sekumpulan persamaan
linear.  Jika kita dapat menyelesaikan $A \vec{x}_k = \vec{e}_k$ dengan $\vec{e}_k$
sebagai vektor yang semua entrinya nol kecuali 1 pada posisi $k^{\text{th}}$, maka
inversnya adalah matriks dengan kolom $\vec{x}_k$ untuk $k=1,2,\ldots,n$
(latihan: mengapa?).  Oleh karena itu, untuk menemukan invers kita tuliskan matriks diperbesar
$n \times 2n$ yang lebih besar, $[ \,A ~|~ I\, ]$, dengan $I$ sebagai matriks
identitas.
Kemudian kita melakukan reduksi baris.
Bentuk eselon baris tereduksi dari $[ \,A ~|~ I\, ]$
akan berbentuk $[ \,I ~|~ A^{-1}\, ]$ jika dan hanya jika
$A$ dapat dibalik.  Kemudian kita cukup membaca invers $A^{-1}$.

\subsection{Latihan}

\begin{exercise}
Selesaikan
$\left[ \begin{smallmatrix}
1 & 2 \\
3 & 4
\end{smallmatrix} \right] \vec{x} =
\left[ \begin{smallmatrix}
5 \\
6
\end{smallmatrix} \right]$ dengan menggunakan invers matriks.
\end{exercise}

\begin{exercise}
Hitung determinan
$\left[ \begin{smallmatrix}
9 & -2 & -6 \\
-8 & 3 & 6 \\
10 & -2 & -6
\end{smallmatrix} \right]$.
\end{exercise}

\begin{exercise}
Hitung determinan
$\left[ \begin{smallmatrix}
1 & 2 & 3 & 1 \\
4 & 0 & 5 & 0 \\
6 & 0 & 7 & 0 \\
8 & 0 & 10 & 1
\end{smallmatrix} \right]$.  Petunjuk: Lakukan ekspansi sepanjang baris atau kolom yang tepat untuk
menyederhanakan perhitungan.
\end{exercise}

\begin{exercise}
Hitung invers
$\left[ \begin{smallmatrix}
1 & 2 & 3 \\
1 & 1 & 1 \\
0 & 1 & 0
\end{smallmatrix} \right]$.
\end{exercise}

\begin{exercise}
Untuk nilai $h$ mana
$\left[ \begin{smallmatrix}
1 & 2 & 3 \\
4 & 5 & 6 \\
7 & 8 & h
\end{smallmatrix} \right]$
tidak dapat dibalik?  Apakah hanya ada satu nilai $h$ semacam itu?  Apakah ada beberapa?  Tak
berhingga banyaknya?
\end{exercise}

\begin{exercise}
Untuk nilai $h$ mana
$\left[ \begin{smallmatrix}
h & 1 & 1 \\
0 & h & 0 \\
1 & 1 & h
\end{smallmatrix} \right]$
tidak dapat dibalik?  Temukan semua $h$ semacam itu.
\end{exercise}

\begin{exercise}
Selesaikan
$\left[ \begin{smallmatrix}
9 & -2 & -6 \\
-8 & 3 & 6 \\
10 & -2 & -6
\end{smallmatrix} \right] \vec{x} =
\left[ \begin{smallmatrix}
1 \\
2 \\
3
\end{smallmatrix} \right]$ .
\end{exercise}

\begin{exercise}
Selesaikan
$\left[ \begin{smallmatrix}
5 & 3 & 7 \\
8 & 4 & 4 \\
6 & 3 & 3
\end{smallmatrix} \right] \vec{x} =
\left[ \begin{smallmatrix}
2 \\
0 \\
0
\end{smallmatrix} \right]$.
\end{exercise}

\begin{exercise}
Selesaikan
$\left[ \begin{smallmatrix}
3 & 2 & 3 & 0 \\
3 & 3 & 3 & 3 \\
0 & 2 & 4 & 2 \\
2 & 3 & 4 & 3
\end{smallmatrix} \right] \vec{x} =
\left[ \begin{smallmatrix}
2 \\
0 \\
4 \\
1
\end{smallmatrix} \right]$.
\end{exercise}

\begin{exercise}
Carilah 3 matriks $2 \times 2$ taknol, $A$, $B$, dan $C$, sedemikian sehingga
$AB = AC$ tetapi $B \neq C$.
\end{exercise}

\setcounter{exercise}{100}

\begin{exercise}
Hitung determinan
$\left[ \begin{smallmatrix}
1 & 1 & 1 \\
2 & 3 & -5 \\
1 & -1 & 0
\end{smallmatrix}\right]$
\end{exercise}
\exsol{%
$-15$
}

\begin{exercise}
Carilah $t$ sedemikian sehingga
$\left[ \begin{smallmatrix}
1 & t \\
-1 & 2
\end{smallmatrix}\right]$
tidak dapat dibalik.
\end{exercise}
\exsol{%
$-2$
}

\begin{exercise}
Selesaikan
$\left[ \begin{smallmatrix}
1 & 1 \\
1 & -1
\end{smallmatrix}\right] \vec{x} =
\left[ \begin{smallmatrix}
10 \\ 20
\end{smallmatrix}\right]$.
\end{exercise}
\exsol{%
$\vec{x} =
\left[ \begin{smallmatrix}
15 \\ -5
\end{smallmatrix}\right]$
}

\begin{exercise}
Andaikan $a, b, c$ adalah bilangan taknol.
Misalkan
$M=\left[ \begin{smallmatrix}
a & 0 \\
0 & b
\end{smallmatrix}\right]$,
$N=\left[ \begin{smallmatrix}
a & 0 & 0 \\
0 & b & 0 \\
0 & 0 & c
\end{smallmatrix}\right]$.
\begin{tasks}(2)
\task Hitung $M^{-1}$.
\task Hitung $N^{-1}$.
\end{tasks}
\end{exercise}
\exsol{%
a) $\left[ \begin{smallmatrix}
\nicefrac{1}{a} & 0 \\
0 & \nicefrac{1}{b}
\end{smallmatrix}\right]$
\quad
b)
$\left[ \begin{smallmatrix}
\nicefrac{1}{a} & 0 & 0 \\
0 & \nicefrac{1}{b} & 0 \\
0 & 0 & \nicefrac{1}{c}
\end{smallmatrix}\right]$
}

%%%%%%%%%%%%%%%%%%%%%%%%%%%%%%%%%%%%%%%%%%%%%%%%%%%%%%%%%%%%%%%%%%%%%%%%%%%%%%

\sectionnewpage
\section{Sistem linear PDB}
\label{linsystems:section}

\sectionnotes{kurang dari 1 perkuliahan\EPref{,
bagian kedua \S5.1 dalam \cite{EP}}\BDref{,
\S7.4 dalam \cite{BD}}}

Fungsi bernilai matriks atau bernilai vektor
hanyalah matriks atau vektor yang entrinya bergantung pada suatu variabel.
Jika $t$ adalah variabel independen, kita menuliskan
\emph{\myindex{fungsi bernilai vektor}}
$\vec{x}(t)$ sebagai
\begin{equation*}
\vec{x}(t) = \begin{bmatrix}
x_1(t) \\
x_2(t) \\
\vdots \\
x_n(t)
\end{bmatrix} .
\end{equation*}
Demikian pula, suatu \emph{\myindex{fungsi bernilai matriks}} $A(t)$ adalah
\begin{equation*}
A(t) =
\begin{bmatrix}
a_{11}(t) & a_{12}(t) & \cdots & a_{1n}(t) \\
a_{21}(t) & a_{22}(t) & \cdots & a_{2n}(t) \\
\vdots & \vdots & \ddots & \vdots \\
a_{n1}(t) & a_{n2}(t) & \cdots & a_{nn}(t)
\end{bmatrix} .
\end{equation*}
Turunan $A'(t)$ atau $\frac{dA}{dt}$
hanyalah fungsi bernilai matriks
yang entri $ij^{\text{th}}$-nya adalah $a_{ij}'(t)$.

Aturan diferensiasi fungsi bernilai matriks
serupa dengan aturan untuk fungsi
biasa.  Misalkan $A(t)$ dan $B(t)$ adalah fungsi bernilai
matriks.  Misalkan $c$ suatu skalar dan $C$ suatu matriks konstan.
Maka
\begin{equation*}
\begin{aligned}
\bigl(A(t)+B(t)\bigr)' & = A'(t) + B'(t), \\
\bigl(A(t)B(t)\bigr)' & = A'(t)B(t) + A(t)B'(t), \\
\bigl(cA(t)\bigr)' & = cA'(t), \\
\bigl(CA(t)\bigr)' & = CA'(t), \\
\bigl(A(t)\,C\bigr)' & = A'(t)\,C .
\end{aligned}
\end{equation*}
Perhatikan urutan perkalian dalam dua ekspresi terakhir.

Suatu \emph{\myindex{sistem linear PDB orde pertama}} adalah sistem yang dapat
dituliskan sebagai persamaan vektor
\begin{equation*}
{\vec{x}}'(t) = P(t)\vec{x}(t) + \vec{f}(t),
\end{equation*}
dengan $P(t)$ fungsi bernilai matriks,
serta $\vec{x}(t)$ dan $\vec{f}(t)$ fungsi bernilai vektor.
Kita akan sering menghilangkan ketergantungan pada $t$ dan
hanya menuliskan ${\vec{x}}' = P\vec{x} + \vec{f}$.  Solusi
sistem tersebut adalah fungsi bernilai vektor
$\vec{x}$ yang memenuhi persamaan vektor itu.

Sebagai contoh, persamaan
\begin{equation*}
\begin{aligned}
x_1' &= 2t x_1 + e^t x_2 + t^2 , \\
x_2' &= \frac{x_1}{t} -x_2 + e^t ,
\end{aligned}
\end{equation*}
dapat dituliskan sebagai
\begin{equation*}
{\vec{x}}' =
\begin{bmatrix}
2t & e^t \\
\nicefrac{1}{t} & -1
\end{bmatrix}
\vec{x}
+
\begin{bmatrix}
t^2 \\
e^t
\end{bmatrix} .
\end{equation*}

Kita terutama akan berfokus pada persamaan yang bukan sekadar linear, tetapi
bahkan merupakan persamaan \emph{\myindex{koefisien konstan}}.  Artinya, matriks $P$
akan konstan; matriks tersebut tidak bergantung pada $t$.

\medskip

Ketika $\vec{f} = \vec{0}$ (vektor nol), kita mengatakan sistem tersebut
\emph{homogen\index{sistem homogen}}.
Untuk sistem linear homogen, kita memiliki
prinsip superposisi, sama seperti untuk persamaan linear homogen tunggal.

\begin{theorem}[Superposisi]\index{superposisi}
Misalkan
${\vec{x}}' = P\vec{x}$ suatu sistem linear homogen PDB.  Andaikan
$\vec{x}_1,\vec{x}_2,\ldots,\vec{x}_n$ adalah $n$ solusi persamaan tersebut
dan $c_1,c_2,\ldots,c_n$ adalah sebarang konstanta, maka
\begin{equation} \label{syshom:eq1}
\vec{x} = c_1 \vec{x}_1 + c_2 \vec{x}_2 + \cdots + c_n \vec{x}_n ,
\end{equation}
juga merupakan solusi.
Selanjutnya, jika ini adalah sistem $n$ persamaan ($P$ berukuran $n\times n$), dan
$\vec{x}_1,\vec{x}_2,\ldots,\vec{x}_n$ bebas linear, maka setiap
solusi $\vec{x}$ dapat dituliskan seperti \eqref{syshom:eq1}.
\end{theorem}

Kebebasan linear untuk fungsi bernilai vektor merupakan gagasan yang sama seperti
untuk fungsi biasa.
Fungsi-fungsi bernilai vektor
$\vec{x}_1,\vec{x}_2,\ldots,\vec{x}_n$
bersifat bebas linear
\index{bebas linear!untuk fungsi bernilai vektor}
apabila
\begin{equation*}
c_1 \vec{x}_1 + c_2 \vec{x}_2 + \cdots + c_n \vec{x}_n  = \vec{0}
\end{equation*}
hanya memiliki solusi $c_1 = c_2 = \cdots = c_n = 0$, dengan persamaan tersebut
harus berlaku untuk semua $t$.

\begin{example}
$\vec{x}_1 = \Bigl[ \begin{smallmatrix} t^2 \\ t \end{smallmatrix} \Bigr]$,
$\vec{x}_2 = \Bigl[ \begin{smallmatrix} 0 \\ 1+t \end{smallmatrix} \Bigr]$,
$\vec{x}_3 = \Bigl[ \begin{smallmatrix} -t^2 \\ 1 \end{smallmatrix} \Bigr]$
bergantung linear karena
$\vec{x}_1 + \vec{x}_3 = \vec{x}_2$, dan hal ini berlaku untuk semua $t$.  Jadi $c_1 =
1$, $c_2 = -1$, dan $c_3 = 1$ di atas akan memenuhi.

Sebaliknya, jika kita sedikit mengubah contoh menjadi
$\vec{x}_1 = \Bigl[ \begin{smallmatrix} t^2 \\ t \end{smallmatrix} \Bigr]$,
$\vec{x}_2 = \Bigl[ \begin{smallmatrix} 0 \\ t \end{smallmatrix} \Bigr]$,
$\vec{x}_3 = \Bigl[ \begin{smallmatrix} -t^2 \\ 1 \end{smallmatrix}
\Bigr]$,
maka fungsi-fungsi tersebut bebas linear.
Pertama, tuliskan
$c_1 \vec{x}_1 + c_2 \vec{x}_2 + c_3 \vec{x}_3  = \vec{0}$ dan perhatikan bahwa persamaan itu
harus berlaku untuk semua $t$.  Kita peroleh
\begin{equation*}
c_1 \vec{x}_1 + c_2 \vec{x}_2 + c_3 \vec{x}_3
=
\begin{bmatrix}
c_1 t^2 - c_3 t^2
\\
c_1 t + c_2 t + c_3
\end{bmatrix}
=
\begin{bmatrix}
0
\\
0
\end{bmatrix} .
\end{equation*}
Dengan kata lain,
$c_1 t^2 - c_3 t^2 = 0$ dan
$c_1 t + c_2 t + c_3 = 0$.
Jika kita menetapkan $t = 0$, persamaan kedua menjadi $c_3 = 0$.  Namun kemudian
persamaan pertama menjadi
$c_1 t^2 = 0$ untuk semua $t$, sehingga $c_1 = 0$.  Dengan demikian persamaan kedua
hanya menjadi $c_2 t = 0$, yang berarti $c_2 = 0$.  Jadi $c_1 = c_2 = c_3 = 0$
adalah satu-satunya solusi dan $\vec{x}_1$,
$\vec{x}_2$, serta $\vec{x}_3$ bebas linear.
\end{example}

Kombinasi linear $c_1 \vec{x}_1 + c_2 \vec{x}_2 + \cdots + c_n
\vec{x}_n$ selalu dapat dituliskan sebagai
\begin{equation*}
X(t)\,\vec{c} ,
\end{equation*}
dengan $X(t)$ matriks yang kolom-kolomnya adalah $\vec{x}_1, \vec{x}_2, \ldots, \vec{x}_n$,
dan $\vec{c}$ vektor kolom dengan entri $c_1, c_2, \ldots, c_n$.
Dengan mengandaikan bahwa $\vec{x}_1,\vec{x}_2,\ldots,\vec{x}_n$ bebas
linear,
fungsi bernilai matriks $X(t)$ disebut \emph{\myindex{matriks fundamental}},
atau \emph{\myindex{solusi matriks fundamental}}.
Perhatikan bahwa matriks fundamental tidak tunggal; jika Anda mengambil sebarang matriks
konstan $n\times n$ yang dapat dibalik, $C$, maka $X(t) C$ adalah matriks fundamental lain.

\medskip

Untuk menyelesaikan sistem linear orde pertama takhomogen, kita menggunakan
teknik yang sama seperti yang kita terapkan untuk menyelesaikan persamaan linear takhomogen tunggal.

\begin{theorem}
Misalkan
${\vec{x}}' = P\vec{x} + \vec{f}$ suatu sistem linear PDB.
Andaikan $\vec{x}_p$ suatu solusi khusus.  Maka setiap solusi
dapat dituliskan sebagai
\begin{equation*}
\vec{x} = \vec{x}_c + \vec{x}_p ,
\end{equation*}
dengan $\vec{x}_c$ suatu solusi bagi
\myindex{persamaan homogen terkait}
(${\vec{x}}' = P\vec{x}$).
\end{theorem}

Prosedur untuk sistem sama seperti untuk persamaan tunggal.
Kita menemukan solusi khusus
bagi persamaan takhomogen, lalu menemukan solusi umum bagi
persamaan homogen terkait, dan akhirnya menjumlahkan keduanya.

\medskip

Baiklah, andaikan Anda telah menemukan solusi umum
${\vec{x}}' = P\vec{x} + \vec{f}$.  Selanjutnya, andaikan Anda diberi
syarat awal berbentuk
\begin{equation*}
\vec{x}(t_0) = \vec{b}
\end{equation*}
untuk suatu $t_0$ tetap dan vektor konstan
$\vec{b}$.  Misalkan $X(t)$ suatu solusi matriks fundamental
bagi persamaan homogen terkait
(yakni,\ kolom-kolom $X(t)$ merupakan solusi).  Solusi umum dapat
dituliskan sebagai
\begin{equation*}
\vec{x}(t) = X(t)\,\vec{c} + \vec{x}_p(t).
\end{equation*}
Kita mencari vektor
$\vec{c}$ sedemikian sehingga
\begin{equation*}
\vec{b} = \vec{x}(t_0) = X(t_0)\,\vec{c} + \vec{x}_p(t_0).
\end{equation*}
Dengan kata lain, untuk $\vec{c}$ kita menyelesaikan sistem persamaan linear takhomogen
\begin{equation*}
X(t_0)\,\vec{c} = \vec{b} - \vec{x}_p(t_0) .
\end{equation*}

\begin{example}
Dalam \examplevref{example:simplesystem}, kita menyelesaikan sistem
\begin{align*}
x_1' & = x_1 , \\
x_2' & = x_1 - x_2 ,
\end{align*}
dengan syarat awal $x_1(0) = 1$, $x_2(0) = 2$.
Mari kita tinjau masalah ini dalam bahasa bagian ini.
Sistem tersebut homogen, sehingga $\vec{f}(t) = \vec{0}$.
Kita tuliskan sistem dan syarat awalnya sebagai
\begin{equation*}
{\vec{x}}'
=
\begin{bmatrix}
1 & 0 \\
1 & -1
\end{bmatrix}
\vec{x} ,
\qquad
\vec{x}(0) =
\begin{bmatrix}
1 \\
2
\end{bmatrix} .
\end{equation*}

Dengan mengabaikan syarat awal sejenak,
solusi umumnya adalah
$x_1 = c_1 e^t $ dan
$x_2 = \frac{c_1}{2}e^{t} + c_2e^{-t}$.
Dengan menetapkan $c_1=1$ dan $c_2=0$, kita memperoleh solusi
$\vec{x}(t) = \left[ \begin{smallmatrix} e^t \\ (1/2) e^t \end{smallmatrix} \right]$.
Dengan menetapkan $c_1=0$ dan $c_2=1$, kita memperoleh
$\vec{x}(t) = \left[ \begin{smallmatrix} 0 \\ e^{-t} \end{smallmatrix} \right]$.
Kedua solusi ini bebas linear,
seperti dapat dilihat dengan menetapkan
$t=0$, dan memperhatikan bahwa vektor konstan yang dihasilkan
bebas linear.
Oleh karena itu, dalam notasi matriks, suatu solusi matriks fundamental adalah
\begin{equation*}
X(t) =
\begin{bmatrix}
e^t & 0 \\
\frac{1}{2} e^t & e^{-t}
\end{bmatrix} .
\end{equation*}

Untuk menyelesaikan masalah nilai awal, kita selesaikan $\vec{c}$ dalam persamaan
\begin{equation*}
X(0)\,\vec{c} = \vec{b} ,
\end{equation*}
atau dengan kata lain,
\begin{equation*}
\begin{bmatrix}
1 & 0 \\
\frac{1}{2} & 1
\end{bmatrix}
\vec{c} =
\begin{bmatrix}
1 \\ 2
\end{bmatrix} .
\end{equation*}
Satu operasi baris elementer menunjukkan
$\vec{c} =
\left[ \begin{smallmatrix} 1 \\ 3/2 \end{smallmatrix} \right]$.
Solusi kita adalah
\begin{equation*}
\vec{x}(t) =
X(t)\,\vec{c} =
\begin{bmatrix}
e^t & 0 \\
\frac{1}{2} e^t & e^{-t}
\end{bmatrix}
\begin{bmatrix}
1 \\ \frac{3}{2}
\end{bmatrix} =
\begin{bmatrix}
e^t \\
\frac{1}{2} e^t + \frac{3}{2} e^{-t}
\end{bmatrix} .
\end{equation*}
Solusi baru ini sesuai dengan solusi sebelumnya dari \sectionref{sec:introtosys}.
\end{example}

\subsection{Latihan}

\begin{exercise}
Tuliskan sistem $x_1' = 2 x_1 - 3t x_2 + \sin t$,
$x_2' = e^t x_1 + 3 x_2 + \cos t$ dalam bentuk
${\vec{x}}' = P(t) \vec{x} + \vec{f}(t)$.
\end{exercise}

\begin{exercise}
\leavevmode
\begin{tasks}
\task
Verifikasi bahwa sistem ${\vec{x}}' =
\left[ \begin{smallmatrix}
1 & 3 \\ 3 & 1
\end{smallmatrix} \right] \vec{x}$ memiliki dua solusi
$\left[ \begin{smallmatrix}
1 \\ 1
\end{smallmatrix} \right] e^{4t}$ dan
$\left[ \begin{smallmatrix}
1 \\ -1
\end{smallmatrix} \right] e^{-2t}$.
\task Tuliskan solusi umumnya.
\task Tuliskan solusi umum dalam bentuk $x_1 = ?$, $x_2 = ?$
(yakni,\ tuliskan rumus untuk setiap elemen solusi).
\end{tasks}
\end{exercise}

\begin{exercise}
Verifikasi bahwa
$\left[ \begin{smallmatrix}
1 \\ 1
\end{smallmatrix} \right] e^{t}$ dan
$\left[ \begin{smallmatrix}
1 \\ -1
\end{smallmatrix} \right] e^{t}$
bebas linear.  Petunjuk: Cukup substitusikan $t=0$.
\end{exercise}

\begin{exercise}
Verifikasi bahwa
$\left[ \begin{smallmatrix}
1 \\ 1 \\ 0
\end{smallmatrix} \right] e^{t}$ dan
$\left[ \begin{smallmatrix}
1 \\ -1 \\ 1
\end{smallmatrix} \right] e^{t}$
serta
$\left[ \begin{smallmatrix}
1 \\ -1 \\ 1
\end{smallmatrix} \right] e^{2t}$
bebas linear.  Petunjuk: Anda harus sedikit lebih cerdik daripada dalam
latihan sebelumnya.
\end{exercise}

\begin{exercise}
Verifikasi bahwa
$\left[ \begin{smallmatrix}
t \\ t^2
\end{smallmatrix} \right]$ dan
$\left[ \begin{smallmatrix}
t^3 \\ t^4
\end{smallmatrix} \right]$
bebas linear.
\end{exercise}

\begin{exercise}
Ambil sistem $x_1' + x_2' = x_1$,
$x_1' - x_2' = x_2$.
\begin{tasks}
\task Tuliskan dalam bentuk
$A {\vec{x}}' = B \vec{x}$ untuk matriks $A$ dan $B$.
\task Hitung $A^{-1}$ dan gunakan hasilnya untuk menuliskan sistem dalam bentuk
${\vec{x}}' = P \vec{x}$.
\end{tasks}
\end{exercise}

\setcounter{exercise}{100}

\begin{exercise}
Apakah
$\left[ \begin{smallmatrix}
e^{2t} \\ e^t
\end{smallmatrix}\right]$
dan
$\left[ \begin{smallmatrix}
e^{t} \\ e^{2t}
\end{smallmatrix}\right]$
bebas linear?  Berikan alasan.
\end{exercise}
\exsol{%
Ya.
}

\begin{exercise}
Apakah
$\left[ \begin{smallmatrix}
\cosh(t) \\ 1
\end{smallmatrix}\right]$,
$\left[ \begin{smallmatrix}
e^{t} \\ 1
\end{smallmatrix}\right]$,
dan
$\left[ \begin{smallmatrix}
e^{-t} \\ 1
\end{smallmatrix}\right]$
bebas linear?  Berikan alasan.
\end{exercise}
\exsol{%
Tidak.
$2 \left[ \begin{smallmatrix}
\cosh(t) \\ 1
\end{smallmatrix}\right] -
\left[ \begin{smallmatrix}
e^{t} \\ 1
\end{smallmatrix}\right]
-
\left[ \begin{smallmatrix}
e^{-t} \\ 1
\end{smallmatrix}\right] = \vec{0}$
}

\begin{exercise}
Tuliskan $x'=3x-y+e^t$, $y'=tx$ dalam notasi matriks.
\end{exercise}
\exsol{%
$\left[ \begin{smallmatrix}
x \\ y
\end{smallmatrix}\right] '
=
\left[ \begin{smallmatrix}
3 & -1 \\
t & 0
\end{smallmatrix}\right]
\left[ \begin{smallmatrix}
x \\ y
\end{smallmatrix}\right]
+
\left[ \begin{smallmatrix}
e^{t} \\ 0
\end{smallmatrix}\right]$
}

\begin{exercise}
\leavevmode
\begin{tasks}
\task Tuliskan $x_1'=2tx_2$, $x_2'=2tx_2$ dalam notasi matriks.
\task Selesaikan dan tuliskan
solusinya dalam notasi matriks.
\end{tasks}
\end{exercise}
\exsol{%
a)
$\vec{x}\,'
=
\left[ \begin{smallmatrix}
0 & 2t \\
0 & 2t
\end{smallmatrix}\right]
\vec{x}$
\quad
b)
$\vec{x}
=
\left[ \begin{smallmatrix}
C_2 e^{t^2} + C_1 \\
C_2 e^{t^2}
\end{smallmatrix}\right]$
}

%%%%%%%%%%%%%%%%%%%%%%%%%%%%%%%%%%%%%%%%%%%%%%%%%%%%%%%%%%%%%%%%%%%%%%%%%%%%%%

\sectionnewpage
\section{Metode nilai eigen}
\label{eigenmethod:section}

%mbxINTROSUBSECTION

\sectionnotes{2 perkuliahan\EPref{, \S5.2 dalam \cite{EP}}\BDref{,
bagian dari \S7.3, \S7.5, dan \S7.6 dalam \cite{BD}}}

Dalam bagian ini, kita akan mempelajari cara menyelesaikan sistem linear homogen
PDB berkoefisien konstan dengan metode nilai eigen.
Andaikan kita memiliki sistem semacam itu
\begin{equation*}
{\vec{x}}' = P\vec{x} ,
\end{equation*}
dengan
$P$ suatu
matriks persegi konstan.
Kita ingin
mengadaptasi metode untuk persamaan tunggal
berkoefisien konstan dengan mencoba fungsi $e^{\lambda t}$
untuk suatu konstanta tak diketahui $\lambda$.
Namun, $\vec{x}$ adalah suatu vektor.  Jadi kita coba $\vec{x} = \vec{v} e^{\lambda t}$, dengan
$\vec{v}$ suatu vektor konstan tak diketahui.  Kita substitusikan $\vec{x}$ ini ke dalam persamaan untuk memperoleh
\begin{equation*}
\underbrace{\lambda \vec{v} e^{\lambda t}}_{{\vec{x}}'} =
\underbrace{P\vec{v} e^{\lambda t}}_{P\vec{x}} .
\end{equation*}
Kita bagi dengan $e^{\lambda t}$ dan perhatikan bahwa kita mencari skalar $\lambda$
dan vektor $\vec{v}$ yang memenuhi persamaan
\begin{equation*}
\lambda \vec{v} = P\vec{v} .
\end{equation*}

Untuk menyelesaikan persamaan ini, kita memerlukan sedikit lagi aljabar linear, yang sekarang
kita tinjau.

\subsection{Nilai eigen dan vektor eigen suatu matriks}

Misalkan $A$ suatu matriks persegi konstan.  Andaikan
$\lambda$ suatu skalar dan
$\vec{v}$ suatu vektor taknol sedemikian sehingga
\begin{equation*}
A \vec{v} = \lambda \vec{v}.
\end{equation*}
Kita menyebut $\lambda$ semacam itu sebagai \emph{\myindex{nilai eigen}} dari $A$,
dan kita
menyebut $\vec{v}$
sebagai \emph{\myindex{vektor eigen}} yang bersesuaian.

\begin{example}
Matriks $\left[ \begin{smallmatrix}
2 & 1 \\
0 & 1
\end{smallmatrix} \right]$ memiliki nilai eigen $\lambda = 2$ dengan
vektor eigen
yang bersesuaian $\left[ \begin{smallmatrix}
1 \\ 0
\end{smallmatrix} \right]$, karena
\begin{equation*}
\begin{bmatrix}
2 & 1 \\
0 & 1
\end{bmatrix}
\begin{bmatrix}
1 \\ 0
\end{bmatrix}
=
\begin{bmatrix}
2 \\
0
\end{bmatrix}
=
2
\begin{bmatrix}
1 \\ 0
\end{bmatrix} .
\end{equation*}
\end{example}

Mari kita hitung nilai eigen untuk sebarang matriks.
Tuliskan kembali persamaan nilai eigen sebagai $A\vec{v}-\lambda \vec{v}=\vec{0}$
dan dengan demikian sebagai
\begin{equation*}
(A - \lambda I)\vec{v} = \vec{0} .
\end{equation*}
Persamaan ini memiliki solusi taknol $\vec{v}$ jika dan hanya jika
$A - \lambda I$ tidak dapat dibalik.
Seandainya $A - \lambda I$ dapat dibalik,
kita dapat menuliskan
${(A - \lambda I)}^{-1}(A - \lambda I)\vec{v} = {(A-\lambda I)}^{-1}\vec{0}$,
yang mengakibatkan $\vec{v} = \vec{0}$.  Oleh karena itu,
$A$ memiliki
nilai eigen $\lambda$ jika dan hanya jika $\lambda$ menyelesaikan persamaan
\begin{equation*}
\det (A-\lambda I) = 0 .
\end{equation*}
Akibatnya, kita dapat menemukan nilai eigen $A$ tanpa
menemukan vektor eigen yang bersesuaian pada saat yang sama.  Vektor eigen harus
ditemukan kemudian, setelah $\lambda$ diketahui.

\begin{example}
Carilah semua nilai eigen dari
$\left[ \begin{smallmatrix}
2 & 1 & 1 \\
1 & 2 & 0 \\
0 & 0 & 2
\end{smallmatrix} \right]$.

Kita tuliskan
\begin{multline*}
\det \left(
\begin{bmatrix}
2 & 1 & 1 \\
1 & 2 & 0 \\
0 & 0 & 2
\end{bmatrix}
- \lambda
\begin{bmatrix}
1 & 0 & 0 \\
0 & 1 & 0 \\
0 & 0 & 1
\end{bmatrix}
\right)
=
\det \left(
\begin{bmatrix}
2-\lambda & 1 & 1 \\
1 & 2-\lambda & 0 \\
0 & 0 & 2-\lambda
\end{bmatrix}
\right)
= \\
=
(2-\lambda) \bigl({(2-\lambda)}^2 - 1\bigr)
=
-(\lambda -1)(\lambda -2)(\lambda-3) .
\end{multline*}
Dengan menetapkan bentuk ini sama dengan nol, kita menemukan bahwa
nilai eigennya adalah $\lambda = 1$, $\lambda = 2$, dan
$\lambda = 3$.
\end{example}

Untuk matriks $n \times n$, polinomial yang kita peroleh dengan
menghitung $\det(A - \lambda I)$ berderajat $n$, dan oleh karena itu,
secara umum kita memiliki $n$ nilai eigen.  Beberapa mungkin berulang, beberapa mungkin
kompleks.

\medskip

Untuk menemukan vektor eigen yang bersesuaian dengan nilai eigen $\lambda$, kita tuliskan
\begin{equation*}
(A-\lambda I) \vec{v} = \vec{0} ,
\end{equation*}
dan selesaikan untuk suatu vektor taktrivial (taknol) $\vec{v}$.
Jika $\lambda$ suatu nilai eigen, akan ada setidaknya satu
variabel bebas, sehingga
untuk
setiap nilai eigen berbeda $\lambda$, kita selalu dapat menemukan vektor eigen.

\begin{example}
Carilah vektor eigen dari
$\left[ \begin{smallmatrix}
2 & 1 & 1 \\
1 & 2 & 0 \\
0 & 0 & 2
\end{smallmatrix} \right]$ yang bersesuaian dengan nilai eigen $\lambda = 3$.

Kita tuliskan
\begin{equation*}
(A-\lambda I) \vec{v} =
\left(
\begin{bmatrix}
2 & 1 & 1 \\
1 & 2 & 0 \\
0 & 0 & 2
\end{bmatrix}
- 3
\begin{bmatrix}
1 & 0 & 0 \\
0 & 1 & 0 \\
0 & 0 & 1
\end{bmatrix}
\right)
\begin{bmatrix}
v_1 \\ v_2 \\ v_3
\end{bmatrix}
=
\begin{bmatrix}
-1 & 1 & 1 \\
1 & -1 & 0 \\
0 & 0 & -1
\end{bmatrix}
\begin{bmatrix}
v_1 \\ v_2 \\ v_3
\end{bmatrix}
=
\vec{0} .
\end{equation*}
Untuk menyelesaikan
sistem persamaan linear ini,
kita tuliskan
matriks diperbesar
\begin{equation*}
\left[
\begin{array}{ccc|c}
-1 & 1 & 1 & 0 \\
1 & -1 & 0 & 0 \\
0 & 0 & -1 & 0
\end{array}
\right] ,
\end{equation*}
dan kita melakukan operasi baris (latihan: operasi mana?) hingga memperoleh:
\begin{equation*}
\left[
\begin{array}{ccc|c}
1 & -1 & 0 & 0 \\
0 & 0 & 1 & 0 \\
0 & 0 & 0 & 0
\end{array}
\right] .
\end{equation*}
Entri $\vec{v}$ harus memenuhi persamaan
$v_1 - v_2 = 0$, $v_3 = 0$, dan $v_2$ merupakan variabel bebas.
Kita pilih $v_2$ secara sebarang (tetapi taknol), kita tetapkan
$v_1 = v_2$, dan tentu saja $v_3 = 0$.
Sebagai contoh, jika kita memilih $v_2 = 1$, maka
$\vec{v} =
\left[ \begin{smallmatrix} 1 \\ 1 \\ 0 \end{smallmatrix} \right]$.
Mari kita verifikasi bahwa $\vec{v}$ benar-benar merupakan vektor eigen yang bersesuaian dengan $\lambda = 3$:
\begin{equation*}
\begin{bmatrix}
2 & 1 & 1 \\
1 & 2 & 0 \\
0 & 0 & 2 \\
\end{bmatrix}
\begin{bmatrix}
1 \\
1 \\
0
\end{bmatrix}
=
\begin{bmatrix}
3 \\
3 \\
0
\end{bmatrix}
=
3
\begin{bmatrix}
1 \\
1 \\
0
\end{bmatrix} .
\end{equation*}
Berhasil!
\end{example}

\begin{exercise}[mudah]
Apakah vektor eigen tunggal?  Dapatkah Anda menemukan vektor eigen lain untuk
$\lambda = 3$ dalam contoh di atas?  Bagaimana hubungan kedua vektor eigen tersebut?
\end{exercise}

\begin{exercise}
Ketika matriks berukuran $2 \times 2$, Anda tidak perlu
%menuliskan matriks diperbesar dan
melakukan operasi baris saat menghitung vektor eigen;
Anda dapat langsung membacanya dari $A-\lambda I$
(jika Anda telah menghitung nilai eigen dengan benar).
Dapatkah Anda melihat alasannya?  Jelaskan.  Cobalah untuk matriks
$\left[ \begin{smallmatrix} 2 & 1 \\ 1 & 2 \end{smallmatrix} \right]$.
\end{exercise}

\subsection{Metode nilai eigen dengan nilai eigen real yang berbeda}

Baiklah\@, kembali ke sistem homogen PDB berkoefisien konstan.  Kita memiliki sistem
\begin{equation*}
{\vec{x}}' = P\vec{x} .
\end{equation*}
Kita mencari nilai eigen $\lambda_1$, $\lambda_2$, \ldots, $\lambda_n$
dari matriks $P$, dan vektor eigen yang bersesuaian
$\vec{v}_1$, $\vec{v}_2$, \ldots, $\vec{v}_n$.
Fungsi-fungsi
$\vec{v}_1 e^{\lambda_1 t}$,
$\vec{v}_2 e^{\lambda_2 t}$, \ldots,
$\vec{v}_n e^{\lambda_n t}$ merupakan solusi sistem persamaan tersebut dan karena itu
$
\vec{x} = c_1 \vec{v}_1 e^{\lambda_1 t} +
c_2 \vec{v}_2 e^{\lambda_2 t} + \cdots +
c_n \vec{v}_n e^{\lambda_n t}
$
merupakan solusi.

\begin{theorem}
Ambil ${\vec{x}}' = P\vec{x}$.  Jika $P$ adalah matriks konstan $n \times n$
yang
memiliki $n$ nilai eigen real berbeda $\lambda_1$, $\lambda_2$, \ldots, $\lambda_n$,
maka terdapat $n$ vektor eigen bersesuaian yang bebas linear,
$\vec{v}_1$, $\vec{v}_2$, \ldots, $\vec{v}_n$, dan solusi umum
${\vec{x}}' = P\vec{x}$
dapat dituliskan sebagai
\begin{equation*}
\mybxbg{~~
\vec{x} = c_1 \vec{v}_1 e^{\lambda_1 t} +
c_2 \vec{v}_2 e^{\lambda_2 t} + \cdots +
c_n \vec{v}_n e^{\lambda_n t} .
~~}
\end{equation*}
\end{theorem}

Solusi matriks fundamental yang bersesuaian adalah
\begin{equation*}
X(t) = \bigl[\, \vec{v}_1 e^{\lambda_1 t} \quad \vec{v}_2 e^{\lambda_2 t}
\quad \cdots \quad \vec{v}_n e^{\lambda_n t} \,\bigr].
\end{equation*}
Artinya, $X(t)$
adalah matriks yang kolom $j^{\text{th}}$-nya
$\vec{v}_j e^{\lambda_j t}$.

\begin{example}
Tinjau sistem
\begin{equation*}
{\vec{x}}'
=
\begin{bmatrix}
2 & 1 & 1 \\
1 & 2 & 0 \\
0 & 0 & 2
\end{bmatrix}
\vec{x} .
\end{equation*}
Carilah solusi umumnya.

Sebelumnya, kita menemukan bahwa nilai eigennya adalah $1,2,3$.  Kita menemukan vektor eigen
$\left[ \begin{smallmatrix} 1 \\ 1 \\ 0 \end{smallmatrix} \right]$
untuk nilai eigen 3.  Demikian pula,
kita menemukan vektor eigen
$\left[ \begin{smallmatrix} 1 \\ -1 \\ 0 \end{smallmatrix} \right]$
untuk nilai eigen 1 dan
$\left[ \begin{smallmatrix} 0 \\ 1 \\ -1 \end{smallmatrix} \right]$
untuk nilai eigen 2 (latihan: periksa).
Solusi umumnya adalah
\begin{equation*}
\vec{x} =
c_1
\begin{bmatrix}
1 \\ -1 \\ 0
\end{bmatrix}
e^t
+
c_2
\begin{bmatrix}
0 \\ 1 \\ -1
\end{bmatrix}
e^{2t}
+
c_3
\begin{bmatrix}
1 \\ 1 \\ 0
\end{bmatrix}
e^{3t}
=
\begin{bmatrix}
c_1 e^t+c_3 e^{3t} \\ -c_1 e^t + c_2 e^{2t} + c_3 e^{3t} \\ - c_2 e^{2t}
\end{bmatrix} .
\end{equation*}
Dalam bentuk solusi matriks fundamental,
\begin{equation*}
\vec{x} = X(t)\, \vec{c}
=
\begin{bmatrix}
e^t & 0 & e^{3t} \\
-e^t & e^{2t} & e^{3t} \\
0 & -e^{2t} & 0
\end{bmatrix}
\begin{bmatrix}
c_1 \\ c_2 \\ c_3
\end{bmatrix} .
\end{equation*}
\end{example}

\begin{exercise}
Periksa bahwa $\vec{x}$ ini benar-benar menyelesaikan sistem tersebut.
\end{exercise}

Catatan: Jika kita menuliskan satu persamaan linear homogen berkoefisien konstan
orde $n^{\text{th}}$
sebagai sistem orde pertama (seperti yang kita lakukan dalam \sectionref{sec:introtosys}),
maka persamaan nilai eigen
\begin{equation*}
\det(P - \lambda I) = 0
\end{equation*}
pada dasarnya sama dengan persamaan karakteristik yang kita peroleh dalam
\sectionref{sec:ccsol} dan \sectionref{sec:hol}.

\subsection{Nilai eigen kompleks}

Suatu matriks sangat mungkin memiliki nilai eigen kompleks meskipun semua entrinya
real.  Ambil, sebagai contoh,
\begin{equation*}
{\vec{x}}' =
\begin{bmatrix}
1 & 1 \\
-1 & 1
\end{bmatrix}
\vec{x} .
\end{equation*}
Mari kita hitung nilai eigen
matriks $P = \left[ \begin{smallmatrix} 1 & 1 \\ -1 & 1 \end{smallmatrix}
\right]$.
\begin{equation*}
\det(P - \lambda I) =
\det\left(
\begin{bmatrix}
1-\lambda & 1 \\
-1 & 1-\lambda
\end{bmatrix}
\right)
= {(1-\lambda)}^2 + 1
= \lambda^2 - 2 \lambda + 2 = 0 .
\end{equation*}
Jadi $\lambda = 1 \pm i$.
Vektor eigen yang bersesuaian juga kompleks.
Mulailah dengan $\lambda = 1-i$.
\begin{align*}
\bigl(P-(1-i) I\bigr) \vec{v} & = \vec{0} , \\
\begin{bmatrix}
i & 1 \\
-1 & i
\end{bmatrix}
\vec{v} = \vec{0}.
\end{align*}
Persamaan $i v_1 + v_2 = 0$ dan $-v_1 + iv_2 = 0$
merupakan kelipatan satu sama lain.  Jadi kita hanya perlu meninjau salah satunya.
Setelah memilih $v_2 = 1$, misalnya, kita memiliki
vektor eigen
$\vec{v} = \left[ \begin{smallmatrix} i \\ 1 \end{smallmatrix} \right]$.
Dengan cara serupa, kita menemukan bahwa
$\left[ \begin{smallmatrix} -i \\ 1 \end{smallmatrix} \right]$
adalah vektor eigen yang bersesuaian dengan nilai eigen $1+i$.

Kita dapat menuliskan solusi sebagai
\begin{equation*}
\vec{x} =
c_1 \begin{bmatrix} i \\ 1 \end{bmatrix} e^{(1-i)t} +
c_2 \begin{bmatrix} -i \\ 1 \end{bmatrix} e^{(1+i)t}
=
\begin{bmatrix}
c_1 i e^{(1-i)t} - c_2 i e^{(1+i)t} \\
c_1 e^{(1-i)t} + c_2 e^{(1+i)t}
\end{bmatrix} .
\end{equation*}
Kemudian kita perlu mencari nilai kompleks $c_1$ dan $c_2$ untuk memenuhi
setiap syarat awal.  Mungkin belum sepenuhnya jelas
bahwa kita memperoleh solusi real.  Setelah menyelesaikan $c_1$ dan $c_2$,
kita dapat menggunakan
\hyperref[eulersformula]{rumus Euler} dan melakukan seluruh rangkaian
yang telah kita lakukan sebelumnya, tetapi kita tidak akan melakukannya.  Pertama-tama kita akan menerapkan
rumus tersebut dengan cara yang lebih cerdas untuk menemukan solusi real yang bebas.

\medskip

Kita menyatakan bahwa sebenarnya kita tidak perlu mencari vektor eigen kedua
(ataupun nilai eigen kedua).  Semua nilai eigen kompleks muncul berpasangan
(karena matriks $P$ real).

Pertama, suatu penyimpangan kecil.  Bagian real dari
bilangan kompleks $z$ dapat dihitung sebagai $\frac{z + \bar{z}}{2}$, dengan
garis di atas $z$ berarti $\overline{a+ib} = a -ib$.  $z$ bergaris atas disebut
\emph{\myindex{konjugat kompleks}} dari $z$.
Jika $a$ suatu bilangan real, maka $\bar{a} = a$.
Demikian pula,
kita memberi garis atas pada seluruh vektor atau matriks dengan mengambil konjugat kompleks
dari setiap entri.  Andaikan suatu matriks $P$ bersifat riil.  Maka
$\overline{P} = P$, sehingga $\overline{P\vec{x}} = \overline{P} \,
\overline{\vec{x}} = P \overline{\vec{x}}$.
Selain itu, konjugat kompleks dari 0 tetaplah 0,
oleh karena itu,
\begin{equation*}
\vec{0} = \overline{\vec{0}} = 
\overline{(P-\lambda I)\vec{v}}
=
(P-\bar{\lambda} I)\overline{\vec{v}} .
\end{equation*}
Dengan kata lain, jika $\lambda = a+ib$ adalah suatu nilai eigen, maka $\bar{\lambda} = a-ib$ juga demikian.
Dan jika $\vec{v}$ adalah suatu vektor eigen yang bersesuaian dengan nilai eigen
$\lambda$, maka $\overline{\vec{v}}$ adalah suatu vektor eigen yang bersesuaian
dengan nilai eigen~$\bar{\lambda}$.  

Andaikan $a + ib$ adalah suatu nilai eigen kompleks dari $P$, dan $\vec{v}$
adalah suatu vektor eigen yang bersesuaian.  Maka
\begin{equation*}
\vec{x}_1 = \vec{v} e^{(a+ib)t}
\end{equation*}
adalah suatu solusi (bernilai kompleks) dari
${\vec{x}}' = P \vec{x}$.  \hyperref[eulersformula]{Rumus Euler}
menunjukkan bahwa $\overline{e^{a+ib}} =
e^{a-ib}$,
sehingga
\begin{equation*}
\vec{x}_2 = \overline{\vec{x}_1} = \overline{\vec{v}} e^{(a-ib)t}
\end{equation*}
juga merupakan suatu solusi.
Karena $\vec{x}_1$ dan $\vec{x}_2$ merupakan solusi, fungsi
\begin{equation*}
\vec{x}_3 =
\operatorname{Re} \vec{x}_1 =
\operatorname{Re} \vec{v} e^{(a+ib)t} =
\frac{\vec{x}_1 + \overline{\vec{x}_1}}{2}  =
\frac{\vec{x}_1 + \vec{x}_2}{2} 
=
\frac{1}{2} \vec{x}_1 + \frac{1}{2}\vec{x}_2
\end{equation*}
juga merupakan suatu solusi.  Dan $\vec{x}_3$ bernilai riil!  Demikian pula, karena
$\operatorname{Im} z = \frac{z-\bar{z}}{2i}$ adalah bagian imajiner, kita memperoleh
bahwa
\begin{equation*}
\vec{x}_4 =
\operatorname{Im} \vec{x}_1 =
\frac{\vec{x}_1 - \overline{\vec{x}_1}}{2i}  =
\frac{\vec{x}_1 - \vec{x}_2}{2i} .
\end{equation*}
juga merupakan suatu solusi bernilai riil.  Ternyata $\vec{x}_3$ dan
$\vec{x}_4$ bebas linear.
Kita akan menggunakan \hyperref[eulersformula]{rumus Euler}
untuk memisahkan bagian riil dan bagian imajiner.

\medskip

Kembali ke contoh soal kita,
\begin{equation*}
\begin{split}
\vec{x}_1 & =
\begin{bmatrix} i \\ 1 \end{bmatrix} e^{(1-i)t}
=
\begin{bmatrix} i \\ 1 \end{bmatrix} \left( e^t \cos t - i e^t \sin t \right)
\\
& =
\begin{bmatrix}
i e^t \cos t + e^t \sin t  \\
e^t \cos t - i e^t \sin t
\end{bmatrix}
=
\begin{bmatrix}
e^t \sin t  \\
e^t \cos t
\end{bmatrix}
+ i
\begin{bmatrix}
e^t \cos t  \\
- e^t \sin t
\end{bmatrix}
.
\end{split}
\end{equation*}
Dua solusi bernilai riil dan bebas linear yang kita cari adalah
\begin{equation*}
\operatorname{Re} \vec{x}_1 = 
\begin{bmatrix}
e^t \sin t  \\
e^t \cos t
\end{bmatrix}
\qquad \text{dan} \qquad
\operatorname{Im} \vec{x}_1 = 
\begin{bmatrix}
e^t \cos t \\
- e^t \sin t
\end{bmatrix}
.
\end{equation*}

\begin{exercise}
Periksa bahwa keduanya benar-benar merupakan solusi.
\end{exercise}

Solusi umumnya adalah
\begin{equation*}
\vec{x}
=
c_1
\begin{bmatrix}
e^t \sin t  \\
e^t \cos t
\end{bmatrix} 
+ c_2
\begin{bmatrix}
e^t \cos t \\
-e^t \sin t
\end{bmatrix} 
=
\begin{bmatrix}
c_1 e^t \sin t + c_2 e^t \cos t \\
c_1 e^t \cos t - c_2 e^t \sin t
\end{bmatrix} .
\end{equation*}
Solusi ini bernilai riil untuk $c_1$ dan $c_2$ riil.  Pada tahap ini, kita akan menyelesaikan
sebarang syarat awal yang mungkin kita miliki untuk menemukan $c_1$ dan $c_2$.

\medskip

Kita rangkum pembahasan tersebut sebagai sebuah teorema.

\begin{theorem}
Misalkan $P$ suatu matriks konstan bernilai riil.
Jika $P$ memiliki nilai eigen kompleks $a+ib$ dan suatu vektor eigen yang bersesuaian
$\vec{v}$, maka $P$ juga memiliki nilai eigen kompleks $a-ib$ dengan
vektor eigen yang bersesuaian $\overline{\vec{v}}$.
Selanjutnya,
${\vec{x}}' = P\vec{x}$ memiliki
dua solusi bernilai riil dan bebas linear
\begin{equation*}
\vec{x}_1 = \operatorname{Re} \vec{v} e^{(a+ib)t}
\qquad
\text{dan}
\qquad
\vec{x}_2 = \operatorname{Im} \vec{v} e^{(a+ib)t} .
\end{equation*}
\end{theorem}

Untuk setiap pasangan nilai eigen kompleks $a+ib$ dan $a-ib$,
kita memperoleh dua solusi bernilai riil yang
bebas linear.
Kemudian kita beralih ke
nilai eigen berikutnya, yang berupa nilai eigen riil atau pasangan nilai eigen
kompleks lainnya.  Jika kita memiliki $n$ nilai eigen berbeda (riil atau kompleks), maka pada akhirnya kita memperoleh $n$ solusi bebas linear.
Jika kita hanya memiliki dua persamaan ($n=2$) seperti pada contoh di atas,
maka setelah menemukan dua solusi, pekerjaan kita
selesai, dan solusi umum kita adalah
\begin{equation*}
\vec{x} =
c_1 \vec{x}_1 + c_2 \vec{x}_2
= 
c_1 \bigl( \operatorname{Re} \vec{v} e^{(a+ib)t} \bigr) +
c_2 \bigl( \operatorname{Im} \vec{v} e^{(a+ib)t} \bigr)
.
\end{equation*}

Sekarang kita dapat menemukan solusi umum bernilai riil untuk sebarang sistem
homogen yang matriksnya memiliki nilai-nilai eigen berbeda.  Ketika kita memiliki nilai eigen
berulang, persoalannya menjadi sedikit lebih rumit dan kita akan membahas
situasi tersebut dalam \sectionref{sec:multeigen}.

\subsection{Latihan}

\begin{exercise}[mudah]
Misalkan $A$ suatu matriks $3 \times 3$ dengan nilai eigen 3 dan suatu
vektor eigen yang bersesuaian $\vec{v} =
\left[ \begin{smallmatrix} 1 \\ -1 \\ 3 \end{smallmatrix} \right]$.
Temukan $A \vec{v}$.
\end{exercise}

\begin{exercise}
\leavevmode
\begin{tasks}
\task
Temukan solusi umum dari $x_1' = 2 x_1$, $x_2' = 3 x_2$ menggunakan
metode nilai eigen (mula-mula tuliskan sistem dalam bentuk
${\vec{x}}' = A \vec{x}$).
\task
Selesaikan sistem dengan menyelesaikan setiap persamaan
secara terpisah dan verifikasi bahwa Anda memperoleh solusi umum yang sama.
\end{tasks}
\end{exercise}

\begin{exercise}
Temukan solusi umum dari $x_1' = 3 x_1 + x_2$,
$x_2' = 2 x_1 + 4 x_2$ menggunakan metode nilai eigen.
\end{exercise}

\begin{exercise}
Temukan solusi umum dari $x_1' = x_1 -2 x_2$,
$x_2' = 2 x_1 + x_2$ menggunakan metode nilai eigen.
Jangan gunakan eksponensial kompleks dalam solusi Anda.
\end{exercise}

\begin{exercise}
\leavevmode
\begin{tasks}
\task
Hitung nilai eigen dan vektor eigen dari
$A = \left[ \begin{smallmatrix}
9 & -2 & -6 \\
-8 & 3 & 6 \\
10 & -2 & -6
\end{smallmatrix} \right]$.
\task
Temukan solusi umum dari ${\vec{x}}' = A \vec{x}$.
\end{tasks}
\end{exercise}

%Nilai eigen berulang, jangan meminta penyelesaian persamaan diferensial dulu
\begin{exercise}
Hitung nilai eigen dan vektor eigen dari
$\left[ \begin{smallmatrix}
-2 & -1 & -1 \\
3 & 2 & 1 \\
-3 & -1 & 0 \\
\end{smallmatrix} \right]$.
\end{exercise}

\begin{exercise}
Misalkan $a,b,c,d,e,f$ bilangan-bilangan.  Temukan nilai eigen dari
$\left[ \begin{smallmatrix}
a & b & c \\
0 & d & e \\
0 & 0 & f \\
\end{smallmatrix} \right]$.
\end{exercise}

\setcounter{exercise}{100}

\begin{exercise}
\leavevmode
\begin{tasks}
\task
Hitung nilai eigen dan vektor eigen dari
$A= \left[ \begin{smallmatrix}
1 & 0 & 3 \\
-1 & 0 & 1 \\
2 & 0 & 2
\end{smallmatrix}\right]$.
\task
Selesaikan sistem
$\vec{x}\,' = A \vec{x}$.
\end{tasks}
\end{exercise}
\exsol{%
a)
Nilai eigen: $4,0,-1$
\quad
Vektor eigen:
$\left[ \begin{smallmatrix}
1 \\ 0 \\ 1
\end{smallmatrix}\right]$,
$\left[ \begin{smallmatrix}
0 \\ 1 \\ 0
\end{smallmatrix}\right]$,
$\left[ \begin{smallmatrix}
3 \\ 5 \\ -2
\end{smallmatrix}\right]$
\\
b)
$\vec{x} = 
C_1
\left[ \begin{smallmatrix}
1 \\ 0 \\ 1
\end{smallmatrix}\right] e^{4t}
+
C_2
\left[ \begin{smallmatrix}
0 \\ 1 \\ 0
\end{smallmatrix}\right] +
C_3
\left[ \begin{smallmatrix}
3 \\ 5 \\ -2
\end{smallmatrix}\right] e^{-t}$
}

\begin{exercise}
\leavevmode
\begin{tasks}
\task
Hitung nilai eigen dan vektor eigen dari
$A=\left[ \begin{smallmatrix}
1 & 1 \\
-1 & 0 
\end{smallmatrix}\right]$.
\task
Selesaikan sistem $\vec{x}\,' = A\vec{x}$.
\end{tasks}
\end{exercise}
\exsol{%
a)
Nilai eigen: $\frac{1+\sqrt{3}i}{2},
\frac{1-\sqrt{3}i}{2}$,
\quad
Vektor eigen:
$\left[ \begin{smallmatrix}
-2 \\ 1-\sqrt{3}i
\end{smallmatrix}\right]$,
$\left[ \begin{smallmatrix}
-2 \\ 1+\sqrt{3}i
\end{smallmatrix}\right]$
\\
b)
$\vec{x} = C_1
e^{t/2}
\left[ \begin{smallmatrix}
-2\cos\bigl(\frac{\sqrt{3}t}{2}\bigr)
\\
\cos\bigl(\frac{\sqrt{3}t}{2}\bigr) + \sqrt{3}\sin\bigl(\frac{\sqrt{3}t}{2}\bigr)
\end{smallmatrix}\right]
+
C_2
e^{t/2}
\left[ \begin{smallmatrix}
- 2\sin\bigl(\frac{\sqrt{3}t}{2}\bigr)
\\
\sin\bigl(\frac{\sqrt{3}t}{2}\bigr) -\sqrt{3}\cos\bigl(\frac{\sqrt{3}t}{2}\bigr)
\end{smallmatrix}\right]$
}

\begin{exercise}
Selesaikan $x_1' = x_2$, $x_2' = x_1$ menggunakan metode nilai eigen.
\end{exercise}
\exsol{%
$\vec{x} = C_1 \left[ \begin{smallmatrix}
1 \\ 1
\end{smallmatrix}\right] e^{t}
+
C_2 \left[ \begin{smallmatrix}
1 \\ -1
\end{smallmatrix}\right] e^{-t}$
}

\begin{exercise}
Selesaikan $x_1' = x_2$, $x_2' = -x_1$ menggunakan metode nilai eigen.
\end{exercise}
\exsol{%
$\vec{x} =
C_1 \left[ \begin{smallmatrix}
\cos(t) \\ -\sin(t)
\end{smallmatrix}\right]
+
C_2 \left[ \begin{smallmatrix}
\sin(t) \\ \cos(t)
\end{smallmatrix}\right]$
}

%%%%%%%%%%%%%%%%%%%%%%%%%%%%%%%%%%%%%%%%%%%%%%%%%%%%%%%%%%%%%%%%%%%%%%%%%%%%%%

\sectionnewpage
\section{Sistem dua dimensi dan medan vektornya}
\label{sec:twodimaut}

\sectionnotes{1 kuliah\EPref{, bagian dari \S6.2 dalam \cite{EP}}\BDref{,
bagian dari \S7.5 dan \S7.6 dalam \cite{BD}}}

Kita berhenti sejenak untuk membahas sistem homogen linear berkoefisien
konstan di bidang.
Banyak intuisi dapat diperoleh dengan mempelajari kasus sederhana ini.
Kita menggunakan koordinat $(x,y)$ untuk bidang seperti biasa,
dan andaikan
$P = \left[ \begin{smallmatrix} a & b \\ c & d \end{smallmatrix} \right]$ 
adalah suatu matriks $2 \times 2$.  Perhatikan sistem
\begin{equation} \label{pln:eq}
\begin{bmatrix} x \\ y \end{bmatrix} ' =
P \begin{bmatrix} x \\ y \end{bmatrix} 
\qquad 
\text{atau}
\qquad
\begin{bmatrix} x \\ y \end{bmatrix} ' =
\begin{bmatrix} a & b \\ c & d \end{bmatrix} 
\begin{bmatrix} x \\ y \end{bmatrix} 
.
\end{equation}
Sistem tersebut otonom (bandingkan bagian ini
dengan \sectionref{auteq:section}),
sehingga kita menggambar suatu medan vektor (lihat akhir
\sectionref{sec:introtosys}).
Kita akan dapat memahami secara visual medan vektor ini
dan solusi-solusi PDB dalam kaitannya dengan
nilai eigen dan vektor eigen dari matriks $P$.
Untuk bagian ini,
kita andaikan $P$ memiliki dua nilai eigen berbeda dan dua vektor eigen yang
bersesuaian.
Kita juga akan mengandaikan bahwa $P$ nonsingular, yaitu,
tidak satu pun nilai eigennya nol.

\medskip

\emph{Kasus 1.}  Andaikan nilai-nilai eigen $P$ riil dan positif.
Kita menemukan dua vektor eigen yang bersesuaian dan menggambarnya di bidang.  Sebagai
contoh, ambil
matriks $\left[ \begin{smallmatrix} 1 & 1 \\ 0 & 2 \end{smallmatrix}
\right]$.
Nilai eigennya adalah 1 dan 2, dan vektor eigen yang bersesuaian adalah
$\left[ \begin{smallmatrix} 1 \\ 0 \end{smallmatrix} \right]$ dan
$\left[ \begin{smallmatrix} 1 \\ 1 \end{smallmatrix} \right]$.  Lihat
\figurevref{pln:source-eigfig}.

\begin{mywrapfig}[14]{3.25in}
\capstart
\diffyincludegraphics{width=3in}{width=4.5in}{pln-source-eig}{%
Diagram bidang xy dengan dua garis yang merepresentasikan vektor (1,0)
dan (1,1), dengan menampilkan dua garis dari titik asal ke masing-masing
titik tersebut.}
\caption{Vektor eigen dari $P$.\label{pln:source-eigfig}}
\end{mywrapfig}

Misalkan $(x,y)$ suatu titik pada garis yang ditentukan oleh vektor eigen
$\vec{v}$ untuk suatu nilai eigen $\lambda$.
Artinya,
$\left[ \begin{smallmatrix} x \\ y \end{smallmatrix} \right] = \alpha \vec{v}$
untuk suatu skalar $\alpha$.
Maka 
\begin{equation*}
\begin{bmatrix} x \\ y \end{bmatrix} '
=
P \begin{bmatrix} x \\ y \end{bmatrix}
=
P ( \alpha \vec{v} ) =  \alpha ( P \vec{v} )
= \alpha \lambda \vec{v} .
\end{equation*}
Turunannya merupakan kelipatan dari $\vec{v}$ dan karenanya mengarah sepanjang
garis yang ditentukan oleh $\vec{v}$.  Karena $\lambda > 0$, turunan itu mengarah ke
arah $\vec{v}$ ketika $\alpha$ positif dan ke arah yang berlawanan
ketika $\alpha$ negatif.  Kita menggambar garis-garis yang ditentukan oleh
vektor eigen dan
panah-panah pada garis untuk menunjukkan arahnya.
Lihat \figurevref{pln:source-eig-arrfig}.

Kita mengisi panah-panah lainnya untuk medan vektor,
dan menggambar beberapa solusi.  Lihat
\figurevref{pln:source-fullfig}.
Gambar tersebut tampak seperti suatu sumber
dengan panah-panah yang keluar dari titik asal.
Karena itu, kita menyebut jenis gambar ini suatu
\emph{\myindex{sumber}} atau terkadang suatu \emph{\myindex{simpul takstabil}}.

\begin{myfig}
\parbox[t]{3.0in}{
 \capstart
 \diffyincludegraphics{width=3in}{width=4.5in}{pln-source-eig-arr}{%
Diagram bidang xy dengan dua panah merah yang merepresentasikan 
dua vektor (1,0) dan (1,1), sekarang dengan panah yang menjauhi
titik asal dan juga garis-garis hitam sepanjang kedua panah; salah satunya
adalah garis x sama dengan y dan yang lain adalah sumbu x.}
 \caption{Vektor eigen dari $P$ beserta arahnya.\label{pln:source-eig-arrfig}}
}
\quad
\parbox[t]{3.0in}{
 \capstart
 \diffyincludegraphics{width=3in}{width=4.5in}{pln-source-full}{%
Diagram seperti gambar sebelumnya, tetapi dengan panah-panah medan arah
ditambahkan dan beberapa lintasan solusi digambarkan.  Semua panah medan arah
menjauhi titik asal, dan panah menjadi semakin panjang ketika
semakin jauh dari titik asal.  Lintasan solusi yang berada di sepanjang kedua
garis yang diberikan oleh vektor eigen kembali berupa garis lurus.  Lintasan
solusi dalam arah lain dari titik asal sedikit melengkung alih-alih
sekadar mengikuti garis lurus, seolah-olah akhirnya ingin mengarah ke
arah yang sama dengan garis diagonal.}
 \caption{Contoh medan vektor sumber dengan vektor eigen dan
 solusi.\label{pln:source-fullfig}}
}
\end{myfig}

\medskip

\emph{Kasus 2.} Andaikan kedua nilai eigen negatif.  Sebagai contoh, ambil
negatif dari
matriks dalam kasus 1,
$\left[ \begin{smallmatrix} -1 & -1 \\ 0 & -2 \end{smallmatrix} \right]$.
Nilai eigennya adalah $-1$ dan $-2$, dan vektor eigen yang bersesuaian
tetap sama,
$\left[ \begin{smallmatrix} 1 \\ 0 \end{smallmatrix} \right]$ dan
$\left[ \begin{smallmatrix} 1 \\ 1 \end{smallmatrix} \right]$.  Perhitungan
dan gambarnya hampir sama.  Perbedaannya adalah
nilai eigennya negatif dan panah-panahnya dibalik.  Kita memperoleh
gambar dalam \figurevref{pln:sink-fullfig}.  Kita menyebut jenis gambar ini suatu
\emph{\myindex{serapan}} atau suatu \emph{\myindex{simpul stabil}}.

\begin{myfig}
\parbox[t]{3.0in}{
 \capstart
 \diffyincludegraphics{width=3in}{width=4.5in}{pln-sink-full}{%
Diagram yang sangat mirip dengan gambar terakhir, tetapi sekarang semua panah
mengarah ke titik asal alih-alih menjauhinya.}
 \caption{Contoh medan vektor serapan dengan vektor eigen dan
 solusi.\label{pln:sink-fullfig}}
}
\quad
\parbox[t]{3.0in}{
 \capstart
 \diffyincludegraphics{width=3in}{width=4.5in}{pln-saddle-full}{%
Diagram suatu pelana; di sini salah satu garis hitam masih merupakan sumbu x,
tetapi garis lainnya adalah garis melalui titik asal dengan kemiringan minus 3.
Panah pada
sumbu x yang merepresentasikan vektor eigen (1,0) mengarah ke kanan,
sedangkan panah yang merepresentasikan vektor eigen (1, minus 3)
sekarang mengarah ke
titik asal karena nilai eigennya negatif.  Panah-panah medan arah
di masing-masing dari 4 sektor di antara kedua garis mula-mula bergerak sepanjang garis dengan
kemiringan minus 3 menuju titik asal, dan setelah panah cukup dekat dengan
sumbu x, panah mulai menjauhi titik asal dalam arah sumbu x.
Beberapa lintasan yang mengikuti panah-panah ini ditampilkan.}
 \caption{Contoh medan vektor pelana dengan vektor eigen dan
 solusi.\label{pln:saddle-fullfig}}
}
\end{myfig}

\medskip

\emph{Kasus 3.} Andaikan satu nilai eigen positif dan satu negatif.
Sebagai contoh, perhatikan
$P = \left[ \begin{smallmatrix} 1 & 1 \\ 0 & -2 \end{smallmatrix} \right]$.
Nilai eigennya adalah 1 dan $-2$, dan vektor eigen yang bersesuaian adalah
$\left[ \begin{smallmatrix} 1 \\ 0 \end{smallmatrix} \right]$ dan
$\left[ \begin{smallmatrix} 1 \\ -3 \end{smallmatrix} \right]$.  Kita membalik
panah-panah pada garis yang bersesuaian dengan nilai eigen negatif dan
memperoleh gambar dalam \figurevref{pln:saddle-fullfig}.  Kita menyebut gambar ini suatu
\emph{\myindex{titik pelana}}.

\medskip

\pagebreak[0]
Untuk tiga kasus berikutnya, kita akan mengandaikan nilai-nilai eigennya kompleks.  Dalam
kasus ini, vektor eigennya juga kompleks dan kita tidak dapat sekadar menggambarnya di
bidang.

\medskip

\pagebreak[0]
\emph{Kasus 4.} Andaikan nilai-nilai eigennya imajiner murni, yaitu,
$\pm ib$.  Sebagai contoh,
misalkan $P = 
\left[ \begin{smallmatrix} 0 & 1 \\ -4 & 0 \end{smallmatrix} \right]$.
Nilai eigennya adalah $\pm 2i$, dan vektor eigen yang bersesuaian adalah
$\left[ \begin{smallmatrix} 1 \\ 2i \end{smallmatrix} \right]$ dan
$\left[ \begin{smallmatrix} 1 \\ -2i \end{smallmatrix} \right]$.  Perhatikan
nilai eigen $2i$ dan vektor eigennya
$\left[ \begin{smallmatrix} 1 \\ 2i \end{smallmatrix} \right]$.
Bagian riil dan imajiner
dari $\vec{v} e^{2it}$ adalah
\begin{equation*}
\operatorname{Re}
\begin{bmatrix} 1 \\ 2i \end{bmatrix} e^{2it} =
\begin{bmatrix} \cos (2t) \\ -2 \sin (2t)  \end{bmatrix} ,
\qquad
\operatorname{Im}
\begin{bmatrix} 1 \\ 2i \end{bmatrix} e^{2it} =
\begin{bmatrix} \sin (2t) \\ 2 \cos (2t) \end{bmatrix} .
\end{equation*}
Kita dapat mengambil sebarang kombinasi linear dari keduanya untuk memperoleh solusi lain;
mana yang kita ambil
bergantung pada syarat
awal.  Sekarang perhatikan bahwa bagian riil merupakan
persamaan parametrik untuk suatu elips.  Demikian pula bagian imajiner
dan, pada kenyataannya, sebarang kombinasi linear dari keduanya.
Inilah yang terjadi secara umum ketika nilai eigennya imajiner murni.
Jadi ketika nilai eigennya imajiner murni, kita memperoleh
\emph{elips\index{elips (medan vektor)}} untuk
solusi-solusinya.  Jenis gambar ini terkadang disebut suatu
\emph{\myindex{pusat}}.  Lihat \figurevref{pln:ellipsesfig}.

\begin{myfig}
\parbox[t]{3.0in}{
 \capstart
 \diffyincludegraphics{width=3in}{width=4.5in}{pln-ellipses}{%
Ditampilkan suatu medan arah dengan panah yang mengarah sedemikian rupa sehingga
mengikutinya akan membawa kita sepanjang suatu elips yang lebih tinggi daripada lebarnya
dan dalam arah jarum jam.
Panah-panah semakin panjang ketika kita semakin jauh dari titik asal.
Beberapa elips semacam itu yang merepresentasikan lintasan ditampilkan.}
 \caption{Contoh medan vektor pusat.\label{pln:ellipsesfig}}
}
\quad
\parbox[t]{3.0in}{
 \capstart
 \diffyincludegraphics{width=3in}{width=4.5in}{pln-spiral-source}{%
Ditampilkan suatu medan arah dengan panah yang berpilin keluar dari titik asal dalam
arah jarum jam.  Panah-panah semakin panjang ketika kita semakin jauh dari titik asal.
Ditampilkan beberapa lintasan yang berpilin keluar dari titik asal.}
 \caption{Contoh medan vektor sumber spiral.\label{pln:spiral-sourcefig}}
}
\end{myfig}

\medskip

\emph{Kasus 5.} Sekarang andaikan nilai eigen kompleks memiliki bagian riil
positif.  Artinya, andaikan nilai eigennya adalah $a \pm ib$ untuk suatu $a > 0$.
Sebagai contoh, misalkan $P = 
\left[ \begin{smallmatrix} 1 & 1 \\ -4 & 1 \end{smallmatrix} \right]$.
Nilai eigennya ternyata $1\pm 2i$, dan vektor eigennya adalah
$\left[ \begin{smallmatrix} 1 \\ 2i \end{smallmatrix} \right]$ dan
$\left[ \begin{smallmatrix} 1 \\ -2i \end{smallmatrix} \right]$.  Kita ambil
$1 + 2i$ dan vektor eigennya
$\left[ \begin{smallmatrix} 1 \\ 2i \end{smallmatrix} \right]$, lalu menemukan
bagian riil dan imajiner dari
$\vec{v} e^{(1+2i)t}$ adalah
\begin{equation*}
\operatorname{Re}
\begin{bmatrix} 1 \\ 2i \end{bmatrix} e^{(1+2i)t} =
e^t
\begin{bmatrix} \cos (2t) \\ -2 \sin (2t)  \end{bmatrix} ,
\qquad
\operatorname{Im}
\begin{bmatrix} 1 \\ 2i \end{bmatrix} e^{(1+2i)t} =
e^t
\begin{bmatrix} \sin (2t) \\ 2 \cos (2t) \end{bmatrix} .
\end{equation*}
Perhatikan $e^t$ di depan solusi-solusi tersebut.  Solusi-solusinya
bertambah besar sambil berputar mengelilingi titik asal.  Karena itu, kita memperoleh
suatu \emph{\myindex{sumber spiral}}.
Lihat \figurevref{pln:spiral-sourcefig}.

\medskip

\emph{Kasus 6.} Terakhir, andaikan nilai eigen kompleks memiliki bagian riil
negatif.  Artinya, andaikan nilai eigennya adalah $-a \pm ib$ untuk suatu $a > 0$.
Sebagai contoh, misalkan $P = 
\left[ \begin{smallmatrix} -1 & -1 \\ 4 & -1 \end{smallmatrix} \right]$.
Nilai eigennya ternyata $-1\pm 2i$, dan vektor eigennya adalah
$\left[ \begin{smallmatrix} 1 \\ -2i \end{smallmatrix} \right]$ dan
$\left[ \begin{smallmatrix} 1 \\ 2i \end{smallmatrix} \right]$.  Kita ambil
$-1 - 2i$ dan vektor eigennya
$\left[ \begin{smallmatrix} 1 \\ 2i \end{smallmatrix} \right]$, lalu menemukan
bagian riil dan imajiner dari
$\vec{v} e^{(-1-2i)t}$ adalah
\begin{equation*}
\operatorname{Re}
\begin{bmatrix} 1 \\ 2i \end{bmatrix} e^{(-1-2i)t} =
e^{-t}
\begin{bmatrix} \cos (2t) \\ 2 \sin (2t)  \end{bmatrix} ,
\qquad
\operatorname{Im}
\begin{bmatrix} 1 \\ 2i \end{bmatrix} e^{(-1-2i)t} =
e^{-t}
\begin{bmatrix} -\sin (2t) \\ 2 \cos (2t) \end{bmatrix} .
\end{equation*}
Perhatikan $e^{-t}$ di depan solusi-solusi tersebut.  Solusi-solusinya
mengecil sambil berputar mengelilingi titik asal.  Karena itu, kita memperoleh
suatu \emph{\myindex{serapan spiral}}.
Lihat \figurevref{pln:spiral-sinkfig}.

\begin{myfig}
\capstart
\diffyincludegraphics{width=3in}{width=4.5in}{pln-spiral-sink}{%
Diagram suatu medan arah yang sangat mirip dengan gambar sebelumnya, tetapi
dibalik.  Panah-panah berpilin berlawanan arah jarum jam
menuju titik asal.  Beberapa lintasan spiral semacam itu ditampilkan.}
\caption{Contoh medan vektor serapan spiral.\label{pln:spiral-sinkfig}}
\end{myfig}

\medskip

Kita rangkum perilaku sistem linear homogen dua dimensi
yang diberikan oleh suatu matriks nonsingular
dalam \tableref{pln:behtab}.
Sistem yang salah satu nilai eigennya nol (matriksnya singular)
muncul dalam praktik dari waktu ke waktu; lihat
\examplevref{sintro:closedbrine-example}, dan gambarnya agak
berbeda (dalam suatu arti, lebih sederhana).  Lihat latihan-latihan.
Ketika nilai eigennya riil dan berulang, 
analisisnya sedikit lebih sulit dan kita belum
membahas cara menyelesaikan sistem seperti itu, tetapi
gagasannya kurang lebih sama---nilai eigen positif yang berulang
berarti suatu sumber dan nilai eigen negatif yang berulang
berarti suatu serapan.


\begin{table}[h!t]
\mybeginframe
\capstart
\begin{center}
\begin{tabular}{@{}ll@{}}
\toprule
Nilai eigen & Perilaku \\
\midrule
riil dan keduanya positif & sumber / simpul takstabil \\
riil dan keduanya negatif & serapan / simpul stabil \\
riil dan berlawanan tanda & pelana \\
imajiner murni & titik pusat / elips \\
kompleks dengan bagian riil positif & sumber spiral \\
kompleks dengan bagian riil negatif & serapan spiral \\
\bottomrule
\end{tabular}
\end{center}
\caption{Ringkasan perilaku sistem linear homogen dua dimensi.\label{pln:behtab}}
\myendframe
\end{table}

\subsection{Latihan}

\begin{exercise}
Ambil persamaan $m x'' + c x' + kx = 0$, dengan $m > 0$, $c \geq 0$, $k > 0$
untuk sistem massa-pegas.
\begin{tasks}
\task Ubah ini menjadi suatu sistem persamaan orde pertama.
\task Klasifikasikan perilaku yang diperoleh untuk berbagai $m, c, k$.
\task Jelaskan berdasarkan intuisi fisik mengapa Anda tidak memperoleh semua
jenis perilaku yang berbeda di sini.
\end{tasks}
\end{exercise}

\begin{exercise}
\pagebreak[2]
Apa yang terjadi dalam kasus ketika $P = 
\left[ \begin{smallmatrix} 1 & 1 \\ 0 & 1 \end{smallmatrix} \right]$?  Dalam
kasus ini, nilai eigennya berulang dan hanya ada satu vektor eigen bebas.
Gambarkan dan jelaskan medan vektornya.
\end{exercise}

\begin{exercise}
Apa yang terjadi dalam kasus ketika $P = 
\left[ \begin{smallmatrix} 1 & 1 \\ 1 & 1 \end{smallmatrix} \right]$?
Gambarkan medan vektornya.
Apakah ini tampak seperti salah satu gambar yang telah kita buat?
\end{exercise}

\begin{exercise}
Perilaku apa saja yang mungkin jika $P$ diagonal, yaitu,
$P = \left[ \begin{smallmatrix} a & 0 \\ 0 & b \end{smallmatrix} \right]$?
Anda boleh mengandaikan bahwa $a$ dan $b$ tidak nol.
\end{exercise}

\begin{exercise}
Ambil sistem dari \examplevref{sintro:closedbrine-example},
$x_1'=\frac{r}{V}(x_2-x_1)$,
$x_2'=\frac{r}{V}(x_1-x_2)$.
Seperti telah kita katakan, salah satu nilai eigennya nol.  Apakah nilai eigen lainnya, bagaimana
bentuk gambarnya, dan apa yang terjadi ketika $t$ menuju tak hingga?
\end{exercise}

\setcounter{exercise}{100}

\begin{exercise}
Jelaskan perilaku sistem-sistem berikut tanpa menyelesaikannya:
\begin{tasks}(2)
\task $x' = x + y$, \quad $y' = x-y$.
\task $x_1' = x_1 + x_2$, \quad $x_2' = 2 x_2$.
\task $x_1' = -2x_2$, \quad $x_2' = 2 x_1$.
\task $x' = x + 3y$, \quad $y' = -2x-4y$.
\task $x' = x - 4y$, \quad $y' = -4x+y$.
\end{tasks}
\end{exercise}
\exsol{%
a) Dua nilai eigen: $\pm \sqrt{2}$, sehingga perilakunya berupa pelana.
\quad
b) Dua nilai eigen: $1$ dan $2$, sehingga perilakunya berupa sumber.
\quad
c) Dua nilai eigen: $\pm 2i$, sehingga perilakunya berupa pusat (elips).
\quad
d) Dua nilai eigen: $-1$ dan $-2$, sehingga perilakunya berupa serapan.
\quad
e) Dua nilai eigen: $5$ dan $-3$, sehingga perilakunya berupa pelana.
}

\begin{exercise}
Andaikan bahwa $\vec{x}\,' = P \vec{x}$ dengan $P$ suatu matriks 2 kali 2
yang memiliki nilai eigen $2\pm i$.  Jelaskan perilakunya.
\end{exercise}
\exsol{%
Sumber spiral.
}

\begin{exercise}
Ambil
$\left[ \begin{smallmatrix}
x \\ y
\end{smallmatrix}\right] '
=
\left[ \begin{smallmatrix}
0 & 1 \\ 0 & 0
\end{smallmatrix}\right]
\left[ \begin{smallmatrix}
x \\ y
\end{smallmatrix}\right]$.
Gambarkan medan vektornya dan jelaskan perilakunya.  Apakah ini salah satu
perilaku yang telah kita lihat sebelumnya?
\end{exercise}
\exsol{%
%mbxSTARTIGNORE
\vtop{\vskip-1ex \hbox{%
%mbxENDIGNORE
\diffyincludegraphics{width=2in}{width=2in}{0100vectorfield}{%
Diagram medan arah dengan semua panah mendatar.
Panah-panah di atas sumbu x mengarah ke kanan dan panah-panah di bawah sumbu x mengarah ke kiri.
Panah-panah semakin panjang ketika semakin jauh dari sumbu x.}
%mbxSTARTIGNORE
}}
%mbxENDIGNORE
\\[6pt]
Solusi tidak bergerak ke mana pun jika $y = 0$.  Ketika $y$ positif,
solusi bergerak (dengan kecepatan konstan)
dalam arah $x$ positif.  Ketika $y$
negatif, solusi bergerak (dengan kecepatan konstan) dalam arah
$x$ negatif.  Ini bukan salah satu perilaku yang telah kita lihat.
\\
Perhatikan bahwa matriks tersebut memiliki nilai eigen ganda 0 dan solusi umumnya adalah
$x = C_1 t + C_2$ dan $y = C_1$, yang sesuai dengan
deskripsi di atas.
}

%%%%%%%%%%%%%%%%%%%%%%%%%%%%%%%%%%%%%%%%%%%%%%%%%%%%%%%%%%%%%%%%%%%%%%%%%%%%%%

\sectionnewpage
\section{Sistem orde kedua dan penerapannya}
\label{sol:section}

%mbxINTROSUBSECTION

\sectionnotes{lebih dari 2 kuliah\EPref{, \S5.4 dalam \cite{EP}}\BDref{,
tidak ada dalam \cite{BD}}}

\subsection{Sistem massa-pegas takteredam}

Meskipun kita memang mengatakan bahwa biasanya kita hanya dapat mempelajari sistem orde pertama, terkadang
lebih mudah menganalisis sistem dengan cara kemunculannya secara alami.
Sebagai contoh, perhatikan 3 massa yang dihubungkan oleh pegas di antara dua
dinding.  Kita dapat memilih sebarang jumlah yang lebih besar, dan matematikanya pada dasarnya
sama, tetapi demi kesederhanaan, sekarang kita memilih 3.  Kita juga mengandaikan tidak ada
gesekan, yaitu sistemnya takteredam\index{gerak takteredam!sistem}.
Massa-massanya adalah $m_1$, $m_2$, dan
$m_3$, dan konstanta pegasnya adalah $k_1$, $k_2$, $k_3$, dan $k_4$.
Misalkan $x_1$ adalah simpangan massa pertama dari posisi diam, dan
$x_2$ serta $x_3$ adalah simpangan massa kedua dan ketiga.  Seperti
biasa, kita menetapkan nilai positif mengarah ke kanan (ketika $x_1$ membesar, massa pertama
bergerak ke kanan).
Lihat \figurevref{sosa:threecartsfig}.

\begin{myfig}
\capstart
\inputpdft{threecartsfig}{%
Diagram tiga kereta beroda dalam satu baris yang diberi label, dari kiri ke kanan,
m sub 1, m sub 2, dan m sub 3 yang menandakan massanya.
Terdapat suatu pegas berlabel k sub 1 yang terpasang pada dinding di sebelah kiri
dan kereta pertama.
Terdapat suatu pegas berlabel k sub 2 yang menghubungkan kereta 1 dan kereta 2.
Terdapat suatu pegas berlabel k sub 3 yang menghubungkan kereta 2 dan kereta 3.
Terdapat suatu pegas berlabel k sub 4 yang menghubungkan kereta 3 ke dinding di sebelah kanan.}
\caption{Sistem massa dan pegas.\label{sosa:threecartsfig}}
\end{myfig}

Sistem sederhana ini muncul di tempat-tempat yang tidak terduga.  Sebagai contoh, 
dunia kita sesungguhnya terdiri dari banyak partikel kecil materi yang saling berinteraksi.
Ketika kita mencoba sistem di atas dengan jauh lebih banyak massa, kita memperoleh aproksimasi yang baik terhadap
perilaku suatu bahan elastis.  Dengan mengambil limit dari
jumlah massa yang menuju tak hingga dengan suatu cara, kita memperoleh persamaan gelombang kontinu
satu dimensi (yang kita pelajari dalam \sectionref{we:section}).
Namun, kita telah menyimpang.

Mari kita susun persamaan untuk \myindex{sistem tiga massa}.
Menurut \myindex{hukum Hooke}, gaya yang bekerja pada massa sama dengan 
kompresi pegas dikalikan konstanta pegas.  Menurut
\myindex{hukum kedua Newton},
gaya adalah massa dikalikan percepatan.  Jadi, jika kita menjumlahkan gaya-gaya yang bekerja
pada setiap massa, menempatkan tanda yang tepat di depan setiap suku bergantung pada
arah kerjanya, dan menyamakannya dengan massa
dikalikan percepatan, kita memperoleh sistem
persamaan yang diinginkan.
\begin{equation*}
\begin{aligned}
m_1 x_1'' &= -k_1 x_1 + k_2 (x_2-x_1)
& & = -(k_1+k_2) x_1 + k_2 x_2 , \\
m_2 x_2'' &= -k_2 (x_2-x_1) + k_3 (x_3-x_2)
& & = k_2 x_1 -(k_2+k_3) x_2 + k_3 x_3 , \\
m_3 x_3'' &= -k_3 (x_3-x_2) - k_4 x_3
& & = k_3 x_2 - (k_3+k_4) x_3 . 
\end{aligned}
\end{equation*}
Kita definisikan matriks-matriks
\begin{equation*}
M =
\begin{bmatrix}
m_1 & 0 & 0 \\
0 & m_2 & 0 \\
0 & 0 & m_3
\end{bmatrix}
\qquad
\text{dan}
\qquad
K =
\begin{bmatrix}
-(k_1+k_2) & k_2 & 0 \\
k_2 & -(k_2+k_3) & k_3 \\
0 & k_3 & -(k_3+k_4)
\end{bmatrix} .
\end{equation*}
Kita tuliskan persamaannya secara ringkas sebagai
\begin{equation*}
M {\vec{x}}'' = K \vec{x} .
\end{equation*}
Pada tahap ini, kita dapat memperkenalkan 3 variabel baru dan menuliskan suatu sistem
6 persamaan orde pertama.  Kita mengklaim susunan sederhana ini lebih mudah ditangani sebagai
suatu sistem orde kedua.
Kita menyebut $\vec{x}$ \emph{\myindex{vektor simpangan}}, $M$
\emph{\myindex{matriks massa}}, dan $K$ \emph{\myindex{matriks kekakuan}}.

\begin{exercise}
Ulangi penyusunan ini untuk 4 massa (temukan matriks $M$ dan $K$).
Lakukan untuk 5 massa.  Dapatkah Anda
menemukan suatu aturan untuk melakukannya bagi $n$ massa?
\end{exercise}

Seperti pada suatu persamaan tunggal, kita ingin \myquote{membagi dengan $M$.}  Ini berarti
menghitung invers dari $M$.
Semua massa taknol dan $M$ adalah suatu
matriks diagonal\index{matriks diagonal}, sehingga menghitung inversnya
mudah:
\begin{equation*}
M^{-1} =
\begin{bmatrix}
\frac{1}{m_1} & 0 & 0 \\
0 & \frac{1}{m_2} & 0 \\
0 & 0 & \frac{1}{m_3}
\end{bmatrix} .
\end{equation*}
Fakta ini langsung mengikuti
dari cara kita mengalikan matriks diagonal.  Sebagai latihan, Anda harus memverifikasi bahwa
$M M^{-1} = M^{-1} M = I$.

\medskip

Misalkan $A = M^{-1}K$.  Kita perhatikan sistem
${\vec{x}}'' = M^{-1}K \vec{x}$, atau
\begin{equation*}
{\vec{x}}'' = A \vec{x} .
\end{equation*}
Banyak sistem dunia nyata dapat dimodelkan oleh persamaan ini.  Demi kesederhanaan,
kita hanya akan membahas
soal massa-dan-pegas yang diberikan.  Kita coba suatu solusi berbentuk
\begin{equation*}
\vec{x} = \vec{v} e^{\alpha t} .
\end{equation*}
Kita menghitung bahwa untuk tebakan ini,
${\vec{x}}'' = \alpha^2 \vec{v} e^{\alpha t}$.
Kita substitusikan tebakan ke persamaan dan memperoleh
\begin{equation*}
\alpha^2 \vec{v} e^{\alpha t} = A\vec{v} e^{\alpha t} .
\end{equation*}
Kita bagi dengan $e^{\alpha t}$ sehingga sampai pada
$\alpha^2 \vec{v} = A\vec{v}$.  Jadi, jika $\alpha^2$ adalah suatu nilai eigen
dari $A$ dan $\vec{v}$ adalah vektor eigen yang bersesuaian, kita telah menemukan
suatu solusi.

Dalam contoh kita, dan dalam penerapan umum lainnya,
$A$ hanya memiliki nilai eigen riil negatif
(dan mungkin suatu nilai eigen nol).  Jadi, kita hanya mempelajari
kasus ini.  Ketika suatu nilai eigen $\lambda$ negatif, itu berarti
$\alpha^2 = \lambda$ negatif.  Oleh karena itu, terdapat suatu bilangan riil $\omega$
sedemikian sehingga $-\omega^2 = \lambda$.  Maka $\alpha = \pm i \omega$.
Untuk $\alpha = i \omega$, solusi yang kita tebak adalah $\vec{x} = \vec{v} e^{i \omega t}$ atau
\begin{equation*}
\vec{x} = \vec{v} \, \bigl(\cos (\omega t) + i \sin (\omega t) \bigr) .
\end{equation*}
Dengan mengambil bagian riil dan imajiner (perhatikan bahwa $\vec{v}$ riil), kita
menemukan bahwa 
$\vec{v} \cos (\omega t)$ dan
$\vec{v} \sin (\omega t)$ adalah solusi-solusi bebas linear.

Jika suatu nilai eigen nol, ternyata $\vec{v}$ dan $\vec{v} t$ keduanya merupakan
solusi, dengan $\vec{v}$ suatu vektor eigen yang bersesuaian dengan nilai eigen
0.

\begin{exercise}
Tunjukkan bahwa jika $A$ memiliki nilai eigen nol dan $\vec{v}$ adalah suatu
vektor eigen yang bersesuaian, maka $\vec{x} = \vec{v} (a + bt)$ adalah suatu solusi
dari ${\vec{x}}'' = A \vec{x}$ untuk
konstanta sebarang $a$ dan $b$.
\end{exercise}

\begin{theorem}
Misalkan $A$ suatu matriks riil $n \times n$
dengan $n$ nilai eigen riil negatif (atau nol) berbeda yang kita
nyatakan dengan $-\omega_1^2 > -\omega_2^2 > \cdots > -\omega_n^2$, dan
vektor eigen yang bersesuaian
dengan
$\vec{v}_1$, $\vec{v}_2$, \ldots, $\vec{v}_n$.  Jika $A$ invertibel
(yaitu, jika $\omega_1 > 0$), maka
\begin{equation*}
\mybxbg{~~
\vec{x}(t)
= \sum_{i=1}^n \vec{v}_i \bigl(a_i \cos (\omega_i t) + b_i \sin (\omega_i t) \bigr) 
~~}
\end{equation*}
adalah solusi umum dari
\begin{equation*}
{\vec{x}}'' = A \vec{x},
\end{equation*}
untuk beberapa konstanta sebarang $a_i$ dan $b_i$.
Jika $A$ memiliki nilai eigen nol, yaitu $\omega_1 = 0$,
dan semua nilai eigen lainnya berbeda dan
negatif, maka solusi umum dapat dituliskan sebagai
\begin{equation*}
\mybxbg{~~
\vec{x}(t) = \vec{v}_1 (a_1 + b_1 t) +
\sum_{i=2}^n \vec{v}_i \bigl(a_i \cos (\omega_i t) + b_i \sin (\omega_i t) \bigr) .
~~}
\end{equation*}
\end{theorem}

Kita menggunakan solusi ini dan susunan dari pengantar
bagian ini bahkan ketika beberapa massa dan pegas tidak ada.  
Sebagai contoh, ketika hanya ada 2 massa dan hanya 2 pegas, cukup ambil
hanya persamaan untuk
kedua massa dan tetapkan semua konstanta pegas untuk pegas yang
tidak ada menjadi nol.

\subsection{Contoh}

\begin{example}
Perhatikan susunan dalam \figurevref{sosa:twocartswallfig}, dengan
$m_1 = \unit[2]{kg}$, $m_2 = \unit[1]{kg}$, $k_1 = \unitfrac[4]{N}{m}$, dan
$k_2 = \unitfrac[2]{N}{m}$.

\begin{myfig}
\capstart
\inputpdft{twocartswallfig}{%
Diagram dua kereta beroda dalam satu baris yang diberi label, dari kiri ke kanan,
m sub 1 dan m sub 2 yang menandakan massanya.
Terdapat suatu pegas berlabel k sub 1 yang terpasang pada dinding di sebelah kiri
dan kereta pertama.
Terdapat suatu pegas berlabel k sub 2 yang menghubungkan kereta 1 dan kereta 2.}
\caption{Sistem massa dan pegas.\label{sosa:twocartswallfig}}
\end{myfig}

Persamaan yang kita tuliskan adalah
\begin{equation*}
\begin{bmatrix}
2 & 0 \\
0 & 1
\end{bmatrix}
{\vec{x}}'' =
\begin{bmatrix}
-(4+2) & 2 \\
2 & -2
\end{bmatrix}
\vec{x} ,
\end{equation*}
atau
\begin{equation*}
{\vec{x}}'' =
\begin{bmatrix}
-3 & 1 \\
2 & -2
\end{bmatrix}
\vec{x} .
\end{equation*}

Kita menemukan nilai eigen $A$ berupa $\lambda = -1, -4$ (latihan).
Kita menemukan vektor eigen yang bersesuaian berupa
$\left[ \begin{smallmatrix} 1 \\ 2 \end{smallmatrix} \right]$ dan
$\left[ \begin{smallmatrix} 1 \\ -1 \end{smallmatrix} \right]$ berturut-turut
(latihan).
Kita periksa teorema tersebut dan perhatikan bahwa $\omega_1 = 1$ dan $\omega_2 = 2$.
Oleh karena itu, solusi umumnya adalah
\begin{equation*}
\vec{x} = 
\begin{bmatrix} 1 \\ 2 \end{bmatrix}
\bigl( a_1 \cos (t) + b_1 \sin (t) \bigr)
+
\begin{bmatrix} 1 \\ -1 \end{bmatrix}
\bigl( a_2 \cos (2t) + b_2 \sin (2t) \bigr) .
\end{equation*}

Dua suku dalam solusi merepresentasikan dua
\emph{modus alami\index{modus alami osilasi}}
atau \emph{modus normal osilasi\index{modus normal osilasi}}.
Dan kedua frekuensi (sudut) tersebut adalah
\emph{frekuensi alami\index{frekuensi alami}}.
Frekuensi alami pertama adalah 1, dan yang kedua adalah 2.
Kedua modus digambarkan dalam \figurevref{sosa:modesfig}.

\begin{myfig}
\capstart
%berkas asli sosa-mode1 sosa-mode2
\diffyincludegraphics{width=6.24in}{width=9in}{sosa-modes-1-2}{%
Diberikan dua grafik, keduanya untuk t dari 0 sampai 10.  Pada grafik sebelah kiri,
dua fungsi sinusoidal (kosinus) bergerak serempak, satu dengan
amplitudo 1 dan satu dengan amplitudo 2; keduanya hanyalah grafik kosinus t
dan 2 kosinus t.  Pada grafik sebelah kanan, dua fungsi kosinus bergerak dalam
arah berlawanan dengan frekuensi dua kali lipat; keduanya hanyalah grafik
kosinus 2 t dan minus kosinus 2 t.}
\caption{Dua modus sistem massa-pegas.  Pada grafik sebelah kiri,
massa-massa bergerak serempak, dan pada grafik sebelah kanan massa-massa bergerak dalam
arah berlawanan.\label{sosa:modesfig}}
\end{myfig}

Mari kita tuliskan solusinya sebagai
\begin{equation*}
\vec{x} = 
\begin{bmatrix} 1 \\ 2 \end{bmatrix}
c_1 \cos (t - \alpha_1 )
+
\begin{bmatrix} 1 \\ -1 \end{bmatrix}
c_2 \cos (2t - \alpha_2 ) .
\end{equation*}
Suku pertama,
\begin{equation*}
\begin{bmatrix} 1 \\ 2 \end{bmatrix}
c_1 \cos (t - \alpha_1 )
=
\begin{bmatrix}
c_1 \cos (t - \alpha_1 ) \\
2c_1 \cos (t - \alpha_1 )
\end{bmatrix} ,
\end{equation*}
bersesuaian dengan modus ketika massa-massa bergerak sinkron
dalam arah yang sama.

Suku kedua,
\begin{equation*}
\begin{bmatrix} 1 \\ -1 \end{bmatrix}
c_2 \cos (2t - \alpha_2 )
=
\begin{bmatrix}
c_2 \cos (2t - \alpha_2 ) \\
- c_2 \cos (2t - \alpha_2 )
\end{bmatrix} ,
\end{equation*}
bersesuaian dengan modus ketika massa-massa bergerak sinkron
tetapi dalam arah berlawanan.

Solusi umum merupakan kombinasi dari kedua modus.  Artinya,
syarat awal menentukan amplitudo dan pergeseran fase setiap modus.
Sebagai contoh, andaikan kita memiliki syarat awal
\begin{equation*}
\vec{x}(0) = 
\begin{bmatrix} 1 \\ -1 \end{bmatrix}
, \qquad
\vec{x}'(0) = 
\begin{bmatrix} 0 \\ 6 \end{bmatrix} .
\end{equation*}
Kita gunakan konstanta $a_j, b_j$ untuk menyelesaikan syarat awal.  Pertama,
\begin{equation*}
\begin{bmatrix} 1 \\ -1 \end{bmatrix}
=
\vec{x}(0) = 
\begin{bmatrix} 1 \\ 2 \end{bmatrix}
a_1
+
\begin{bmatrix} 1 \\ -1 \end{bmatrix}
a_2 
=
\begin{bmatrix} a_1+a_2 \\2a_1 - a_2 \end{bmatrix} .
\end{equation*}
Kita selesaikan (latihan) untuk memperoleh $a_1 = 0$, $a_2 = 1$.
Untuk menemukan $b_1$ dan $b_2$, mula-mula kita diferensialkan:
\begin{equation*}
{\vec{x}}' = 
\begin{bmatrix} 1 \\ 2 \end{bmatrix}
\bigl( - a_1 \sin (t) + b_1 \cos (t) \bigr)
+
\begin{bmatrix} 1 \\ -1 \end{bmatrix}
\bigl( - 2a_2 \sin (2t) + 2 b_2 \cos (2t) \bigr) .
\end{equation*}
Sekarang kita selesaikan:
\begin{equation*}
\begin{bmatrix} 0 \\ 6 \end{bmatrix}
=
{\vec{x}}'(0) = 
\begin{bmatrix} 1 \\ 2 \end{bmatrix}
b_1
+
\begin{bmatrix} 1 \\ -1 \end{bmatrix}
2 b_2
=
\begin{bmatrix} b_1+2b_2 \\ 2b_1-2b_2 \end{bmatrix} .
\end{equation*}
Sekali lagi, selesaikan (latihan) untuk memperoleh $b_1 = 2$, $b_2 = -1$.  Jadi solusi kita adalah
\begin{equation*}
\vec{x} = 
\begin{bmatrix} 1 \\ 2 \end{bmatrix}
2 \sin (t)
+
\begin{bmatrix} 1 \\ -1 \end{bmatrix}
\bigl( \cos (2t) - \sin (2t) \bigr)
=
\begin{bmatrix}
2 \sin (t) + \cos(2t)- \sin(2t) \\
4 \sin (t) - \cos(2t) + \sin(2t)
\end{bmatrix} .
\end{equation*}
Grafik kedua simpangan, $x_1$ dan $x_2$, dari kedua kereta terdapat dalam
\figurevref{sosa:superposfig}.
\begin{myfig}
\capstart
\diffyincludegraphics{width=3in}{width=4.5in}{sosa-superpos}{%
Dua fungsi digambarkan untuk t sama dengan 0 hingga t sama dengan 10.
Geraknya agak
rumit tetapi tetap periodik.}
\caption{Superposisi kedua modus berdasarkan syarat
awal yang diberikan.\label{sosa:superposfig}}
\end{myfig}
\end{example}

\begin{example} \label{sosa:railcarexample}
Kita memiliki dua kereta rel
mainan.  Kereta 1 bermassa \unit[2]{kg} bergerak dengan kecepatan \unitfrac[3]{m}{s}
menuju kereta rel kedua
yang bermassa \unit[1]{kg}.  Terdapat suatu penyangga pada kereta rel kedua yang
mulai bekerja saat kedua kereta bertabrakan (penyangga itu menghubungkan kedua kereta)
dan tidak terlepas.
Penyangga tersebut bekerja seperti pegas dengan konstanta pegas $k=\unitfrac[2]{N}{m}$.
Kereta kedua
berjarak 10 meter dari suatu dinding.  Lihat \figurevref{sosa:railcarscrashfig}.

\begin{myfig}
\capstart
\inputpdft{railcarscrash}{%
Diagram dua kereta beroda; di sebelah kiri terdapat kereta berlabel
m1 dengan panah mengarah ke kanan yang menunjukkan bahwa kereta ini bergerak ke kanan.
Sedikit lebih jauh terdapat kereta kedua berlabel m2 dengan pegas di sisi
kiri berlabel k, yang tidak terhubung dengan apa pun pada sisi kirinya.
Terdapat dinding di sebelah kanan dan suatu label menunjukkan bahwa dinding itu berjarak 10 meter
di sebelah kanan kereta kedua.}
\caption{Tumbukan dua kereta rel.\label{sosa:railcarscrashfig}}
\end{myfig}

Kita ingin mengajukan beberapa pertanyaan.  Pada waktu berapa setelah kereta-kereta terhubung
tumbukan dengan dinding terjadi?  Berapakah kecepatan kereta 2 saat menabrak
dinding?

Baiklah\@, pertama-tama mari kita susun sistemnya.  Misalkan $t=0$ adalah waktu
ketika kedua kereta terhubung.  Misalkan $x_1$ adalah simpangan kereta pertama
dari posisi pada $t=0$, dan misalkan $x_2$ adalah
simpangan kereta kedua dari lokasi asalnya.  Maka
waktu ketika $x_2(t) = 10$ tepat merupakan waktu terjadinya tumbukan dengan dinding.
Untuk $t$ ini, $x_2'(t)$ adalah kecepatan saat tumbukan.  Sistem ini bekerja persis seperti
sistem dalam contoh sebelumnya tetapi tanpa $k_1$.  Jadi persamaannya adalah
\begin{equation*}
\begin{bmatrix}
2 & 0 \\
0 & 1
\end{bmatrix}
{\vec{x}}'' =
\begin{bmatrix}
-2 & 2 \\
2 & -2
\end{bmatrix}
\vec{x} ,
\end{equation*}
atau
\begin{equation*}
{\vec{x}}'' =
\begin{bmatrix}
-1 & 1 \\
2 & -2
\end{bmatrix}
\vec{x} .
\end{equation*}

Kita hitung nilai eigen dari $A$.  Tidak sulit melihat bahwa nilai eigennya
adalah 0 dan $-3$ (latihan).  Selanjutnya,
vektor eigennya adalah
$\left[ \begin{smallmatrix} 1 \\ 1 \end{smallmatrix} \right]$ dan
$\left[ \begin{smallmatrix} 1 \\ -2 \end{smallmatrix} \right]$ berturut-turut
(latihan). 
Maka
$\omega_1 = 0$,
$\omega_2 = \sqrt{3}$, dan menurut bagian kedua teorema,
solusi umumnya adalah
\begin{equation*}
\begin{split}
\vec{x} & = 
\begin{bmatrix} 1 \\ 1 \end{bmatrix}
\left( a_1 + b_1 t \right) 
+
\begin{bmatrix} 1 \\ -2 \end{bmatrix}
\left( a_2 \cos ( \sqrt{3} \, t) + b_2 \sin ( \sqrt{3} \, t ) \right)
\\
& =
\begin{bmatrix}
a_1 + b_1 t + a_2 \cos ( \sqrt{3} \, t ) + b_2 \sin ( \sqrt{3} \, t
) \\
a_1 + b_1 t - 2 a_2 \cos ( \sqrt{3} \, t ) - 2 b_2 \sin ( \sqrt{3}
\, t ) 
\end{bmatrix} .
\end{split}
\end{equation*}

Sekarang kita terapkan syarat awal.  Mula-mula kedua kereta mulai pada posisi 0,
sehingga $x_1 (0) = 0$ dan $x_2(0) = 0$.  Kereta pertama bergerak dengan kecepatan
\unitfrac[3]{m}{s},
sehingga $x_1'(0) = 3$, dan kereta kedua mulai dari keadaan diam, sehingga $x_2'(0) = 0$.
Syarat pertama menyatakan
\begin{equation*}
\vec{0} = \vec{x}(0) = 
\begin{bmatrix}
a_1 + a_2 \\
a_1 - 2 a_2 
\end{bmatrix} .
\end{equation*}
Tidak sulit melihat bahwa $a_1 = a_2 = 0$.  Kita tetapkan $a_1=0$
dan $a_2=0$ dalam $\vec{x}(t)$,
lalu diferensialkan untuk memperoleh
\begin{equation*}
{\vec{x}}'(t)
=
\begin{bmatrix}
b_1 + \sqrt{3} \, b_2 \cos ( \sqrt{3} \, t ) \\
b_1 - 2 \sqrt{3} \, b_2 \cos ( \sqrt{3} \, t )
\end{bmatrix} .
\end{equation*}
Jadi,
\begin{equation*}
\begin{bmatrix} 3 \\ 0 \end{bmatrix} = 
{\vec{x}}'(0)
=
\begin{bmatrix}
b_1 + \sqrt{3} \, b_2 \\
b_1 - 2 \sqrt{3} \, b_2 
\end{bmatrix} .
\end{equation*}
Dengan menyelesaikan kedua persamaan ini, kita
memperoleh $b_1 = 2$ dan $b_2 =
\frac{1}{\sqrt{3}}$.  Oleh karena itu, posisi kereta-kereta kita adalah
(hingga tumbukan dengan dinding)
\begin{equation*}
\vec{x} = 
\begin{bmatrix}
2 t + \frac{1}{\sqrt{3}} \sin ( \sqrt{3} \, t ) \\
2 t - \frac{2}{\sqrt{3}} \sin ( \sqrt{3} \, t )
\end{bmatrix} .
\end{equation*}
Perhatikan bagaimana adanya nilai eigen nol menghasilkan suatu suku yang memuat $t$.
Ini berarti bahwa kereta-kereta akan bergerak ke arah positif ketika
waktu bertambah, seperti yang kita harapkan.


Yang sebenarnya kita minati adalah ekspresi kedua, yaitu ekspresi untuk $x_2$.
Kita memiliki $x_2(t) = 
2 t - \frac{2}{\sqrt{3}} \sin ( \sqrt{3} \, t)$.  Lihat \figurevref{sosa:railcarfig}
untuk grafik $x_2$ terhadap waktu.

Hanya dari grafik, kita dapat melihat bahwa waktu tumbukan akan sedikit
lebih dari 5 detik sejak waktu nol.  Untuk itu, kita harus menyelesaikan
persamaan $10 = x_2(t) = 2 t - \frac{2}{\sqrt{3}} \sin ( \sqrt{3} \,
t)$.
Dengan menggunakan komputer (atau bahkan kalkulator grafik),
kita memperoleh $t_{\text{impact}} \approx 5.22$ detik.

\begin{mywrapfig}{3.25in}
\capstart
\diffyincludegraphics{width=3in}{width=4.5in}{sosa-railcar}{%
Grafik suatu fungsi untuk t sama dengan 0 hingga t sama dengan 6.  Fungsi dimulai dari
nol, mulai bergerak mendatar, lalu melengkung naik
dan melengkung kembali ke suatu dataran mendatar antara
t sama dengan 3 dan t sama dengan 4 pada nilai
sedikit kurang dari 7,5.  Fungsi tersebut 
kemudian kembali melengkung naik dan mencapai nilai 10 pada t yang sedikit
lebih besar dari 5.}
\caption{Posisi kereta kedua terhadap waktu (dengan mengabaikan dinding).\label{sosa:railcarfig}}
\end{mywrapfig}


Kecepatan kereta kedua adalah $x_2' = 2 - 2 \cos ( \sqrt{3} \, t)$.
Pada waktu tumbukan (5,22 detik sejak $t=0$), kita memperoleh 
$x_2'(t_{\text{impact}}) \approx 3.85$.
%
Kecepatan maksimum adalah maksimum dari $2 - 2 \cos ( \sqrt{3} \, t )$, yaitu 4.
Kita bergerak dengan kecepatan yang hampir maksimum saat menabrak dinding.

\medskip

Andaikan Bob adalah orang mungil yang duduk di kereta 2.  Bob memegang Martini dan
ingin agar minumannya tidak tumpah.  Mari kita andaikan Bob tidak akan menumpahkan
Martini-nya
ketika kereta pertama terhubung dengan kereta 2, tetapi jika kereta 2 menabrak dinding dengan
kecepatan lebih besar dari nol, Bob akan menumpahkan minumannya.  Andaikan Bob
dapat memindahkan kereta 2
beberapa meter mendekati atau menjauhi dinding (ia tidak dapat bergerak sampai ke
dinding, dan juga tidak dapat menyingkir dari jalur kereta pertama).  Adakah suatu
jarak \myquote{aman}
baginya?  Jarak sedemikian sehingga tumbukan dengan
dinding terjadi pada kecepatan nol?

Jawabannya ya.  Pada \figurevref{sosa:railcarfig},
perhatikan \myquote{dataran} antara $t=3$ dan $t=4$.  Terdapat suatu titik tempat
kecepatannya nol.  Untuk menemukannya, kita selesaikan $x_2'(t) = 0$.  Ini terjadi ketika
$\cos ( \sqrt{3} \, t) = 1$, atau dengan kata lain ketika $t = \frac{2 \pi}{\sqrt{3}}, 
\frac{4 \pi}{\sqrt{3}},\ldots$, dan seterusnya.  Kita substitusikan nilai pertama untuk memperoleh
$x_2\left(\frac{2 \pi}{\sqrt{3}}\right) = 
\frac{4 \pi}{\sqrt{3}} \approx 7.26$.  Jadi, suatu jarak \myquote{aman} adalah sekitar 7 seperempat
meter dari dinding.

Sebagai alternatif, Bob dapat bergerak menjauhi dinding
ke arah kereta 2 yang datang, dengan jarak aman lainnya adalah
$x_2 \left( \frac{4 \pi}{\sqrt{3}} \right) = \frac{8 \pi}{\sqrt{3}}
\approx 14.51$, dan seterusnya.  Kita dapat menggunakan semua nilai
$t$ yang berbeda sedemikian sehingga $x_2'(t) = 0$.  Tentu saja, $t=0$ juga suatu solusi,
yang bersesuaian dengan $x_2 = 0$, tetapi
itu berarti berdiri tepat di dinding.
\end{example}

\subsection{Osilasi paksa}

Akhirnya, kita beralih ke osilasi paksa\index{gerak paksa!sistem}.
Andaikan sekarang sistem kita adalah
\begin{equation} \label{sosa:forcedeq}
{\vec{x}}'' = A \vec{x} + \vec{F} \cos ( \omega t) .
\end{equation}
Artinya, kita menambahkan pemaksaan periodik pada sistem dalam arah
vektor $\vec{F}$.

Seperti sebelumnya, sistem ini hanya mengharuskan kita menemukan satu solusi khusus
$\vec{x}_p$, menambahkannya ke solusi umum dari sistem homogen terkait
$\vec{x}_c$, dan kita akan memperoleh solusi umum untuk \eqref{sosa:forcedeq}.
Andaikan $\omega$ bukan salah satu frekuensi alami dari
${\vec{x}}'' = A \vec{x}$, maka kita dapat menebak
\begin{equation*}
\vec{x}_p = \vec{c} \cos (\omega t) ,
\end{equation*}
di mana $\vec{c}$ adalah suatu vektor konstan yang tidak diketahui.  Perhatikan bahwa kita tidak perlu
menggunakan sinus karena hanya ada turunan kedua.  Kita selesaikan $\vec{c}$ untuk
menemukan $\vec{x}_p$.  Ini sesungguhnya hanyalah metode
\emph{koefisien tak tentu\index{koefisien tak tentu!untuk sistem orde kedua}}
untuk sistem.  Mari kita diferensialkan $\vec{x}_p$ dua kali untuk memperoleh
\begin{equation*}
{\vec{x}_p}'' = -\omega^2 \vec{c} \cos (\omega t) .
\end{equation*}
Substitusikan $\vec{x}_p$ dan ${\vec{x}_p}''$ ke dalam persamaan \eqref{sosa:forcedeq}:
\begin{equation*}
\overbrace{
-\omega^2 \vec{c} \cos (\omega t)
}^{{\vec{x}_p}''}
=
\overbrace{
A \vec{c} \cos (\omega t) 
}^{A \vec{x}_p}
+ \vec{F} \cos (\omega t) .
\end{equation*}
Kita coret kosinus dan susun ulang persamaan untuk memperoleh
\begin{equation*}
(A +\omega^2 I) \vec{c}
=
- \vec{F} .
\end{equation*}
Jadi,
\begin{equation*}
\vec{c}
=
{(A +\omega^2 I)}^{-1}
(-\vec{F} ).
\end{equation*}
Tentu saja, ini hanya mungkin jika
$(A+ \omega^2 I) = \bigl(A- (-\omega^2) I\bigr)$
invertibel.  Matriks itu invertibel jika dan hanya jika
$-\omega^2$ bukan suatu nilai eigen dari $A$.  Hal itu benar jika dan hanya jika $\omega$
bukan suatu frekuensi alami dari sistem.

Kita sedikit menyederhanakan persoalan.  Jika kita ingin
suku pemaksa berada dalam satuan gaya, misalnya Newton, maka kita harus menulis
\begin{equation*}
M \vec{x}'' = K \vec{x} + \vec{G} \cos(\omega t) .
\end{equation*}
Jika kemudian kita menuliskan segala sesuatu dalam $A = M^{-1} K$, kita memperoleh
\begin{equation*}
\vec{x}'' = M^{-1}K \vec{x} + M^{-1} \vec{G} \cos(\omega t) 
\qquad \text{atau} \qquad
\vec{x}'' = A \vec{x} + \vec{F} \cos(\omega t) ,
\end{equation*}
di mana $\vec{F} = M^{-1} \vec{G}$.

\begin{example}
Mari kita ambil contoh dalam \figurevref{sosa:twocartswallfig} dengan
parameter yang sama seperti sebelumnya:
$m_1 = 2$, $m_2 = 1$, $k_1 = 4$, dan $k_2 = 2$.  Sekarang andaikan
terdapat gaya $2 \cos (3t)$ yang bekerja pada kereta kedua.

Persamaannya adalah
\begin{equation*}
\begin{bmatrix}
2 & 0 \\
0 & 1
\end{bmatrix}
{\vec{x}}'' =
\begin{bmatrix}
-(4+2) & 2 \\
2 & -2
\end{bmatrix}
\vec{x} 
+ 
\begin{bmatrix}
0 \\ 2
\end{bmatrix}
\cos (3 t)
%mbxSTARTIGNORE
\qquad
%mbxENDIGNORE
%mbxlatex \quad
\text{atau}
%mbxSTARTIGNORE
\qquad
%mbxENDIGNORE
%mbxlatex \quad
{\vec{x}}'' =
\begin{bmatrix}
-3 & 1 \\
2 & -2
\end{bmatrix}
\vec{x} 
+ 
\begin{bmatrix}
0 \\ 2
\end{bmatrix}
\cos (3 t) .
\end{equation*}
Kita telah menyelesaikan persamaan homogen terkait sebelumnya dan memperoleh
solusi komplementer
\begin{equation*}
\vec{x}_c =
\begin{bmatrix} 1 \\ 2 \end{bmatrix}
\bigl( a_1 \cos (t) + b_1 \sin (t) \bigr)
+
\begin{bmatrix} 1 \\ -1 \end{bmatrix}
\bigl( a_2 \cos (2t) + b_2 \sin (2t) \bigr) .
\end{equation*}

Frekuensi alaminya adalah 1 dan 2.  Karena 3 bukan suatu
frekuensi alami, kita coba $\vec{c} \cos (3t)$.
Kita inverskan $(A+3^2 I)$:
\begin{equation*}
{\left( \begin{bmatrix}
-3 & 1 \\
\noalign{\smallskip}
2 & -2
\end{bmatrix}
+3^2 I\right)}^{-1}
=
{\begin{bmatrix}
6 & 1 \\
\noalign{\smallskip}
2 & 7
\end{bmatrix}}^{-1}
=
\begin{bmatrix}
\frac{7}{40} & \frac{-1}{40} \\
\noalign{\smallskip}
\frac{-1}{20} & \frac{3}{20}
\end{bmatrix} .
\end{equation*}
Oleh karena itu,
\begin{equation*}
\vec{c} = 
{(A +\omega^2 I)}^{-1}
(-\vec{F} ) = 
\begin{bmatrix}
\frac{7}{40} & \frac{-1}{40} \\
\noalign{\smallskip}
\frac{-1}{20} & \frac{3}{20}
\end{bmatrix}
\begin{bmatrix}
0 \\
\noalign{\smallskip}
-2
\end{bmatrix}
=
\begin{bmatrix}
\frac{1}{20} \\
\noalign{\smallskip}
\frac{-3}{10}
\end{bmatrix} .
\end{equation*}

Dengan menggabungkannya dengan solusi umum dari masalah
homogen terkait, kita memperoleh bahwa solusi umum dari
${\vec{x}}'' = A \vec{x} + \vec{F} \cos (\omega t)$ adalah
\begin{equation*}
%mbxlatex \begin{split}
\vec{x}
%mbxlatex &
= \vec{x}_c + \vec{x}_p
%mbxlatex \\
%mbxlatex &
=
\begin{bmatrix} 1 \\
\noalign{\smallskip}
2 \end{bmatrix}
\bigl( a_1 \cos (t) + b_1 \sin (t) \bigr)
+
\begin{bmatrix} 1 \\
\noalign{\smallskip}
-1 \end{bmatrix}
\bigl( a_2 \cos (2t) + b_2 \sin (2t) \bigr)
+
\begin{bmatrix}
\frac{1}{20} \\
\noalign{\smallskip}
\frac{-3}{10}
\end{bmatrix}
\cos (3t) .
%mbxlatex \end{split}
\end{equation*}
Kemudian kita akan menyelesaikan konstanta
$a_1$, $a_2$, $b_1$, dan $b_2$ menggunakan
sebarang syarat awal yang diberikan.
\end{example}

Perhatikan bahwa untuk gaya $\vec{f}$ yang diberikan, kita menuliskan 
persamaan sebagai $M {\vec{x}}'' = K \vec{x} + \vec{f}$ agar satuannya
benar. Kemudian kita tuliskan
${\vec{x}}'' = M^{-1}K \vec{x} + M^{-1}\vec{f}$.  Suku 
$\vec{g} = M^{-1} \vec{f}$ dalam ${\vec{x}}'' = A \vec{x} + \vec{g}$ memiliki satuan gaya
per satuan massa.

Jika $\omega$ adalah suatu frekuensi alami sistem,
\emph{\myindex{resonansi}} dapat terjadi,
karena kita harus mencoba suatu solusi khusus berbentuk
\begin{equation*}
\vec{x}_p =
\vec{c} \, t \sin (\omega t) +
\vec{d} \, \cos (\omega t) .
\end{equation*}
Ini dengan mengandaikan bahwa nilai eigen matriks koefisien berbeda-beda.
Selanjutnya, perhatikan bahwa amplitudo solusi ini bertambah tanpa batas ketika $t$ bertambah.

\subsection{Latihan}

\begin{exercise}
Temukan suatu solusi khusus untuk
\begin{equation*}
{\vec{x}}'' =
\begin{bmatrix}
-3 & 1 \\
2 & -2
\end{bmatrix}
\vec{x} 
+ 
\begin{bmatrix}
0 \\ 2
\end{bmatrix}
\cos (2 t) .
\end{equation*}
\end{exercise}

\begin{exercise}[menantang]
Perhatikan contoh dalam \figurevref{sosa:twocartswallfig} dengan
parameter yang sama seperti sebelumnya:
$m_1 = 2$, $k_1 = 4$, dan $k_2 = 2$, kecuali $m_2$, yang tidak diketahui.
Andaikan terdapat suatu gaya $\cos (5 t)$ yang bekerja pada massa pertama.
Temukan suatu $m_2$ sedemikian sehingga terdapat suatu solusi khusus dengan
massa pertama tidak bergerak.

Catatan: Gagasan ini disebut \emph{\myindex{peredaman dinamis}}.
Dalam praktik, akan ada sejumlah kecil
peredaman sehingga sebarang solusi transien akan menghilang dan setelah waktu yang cukup
lama, massa pertama akan selalu berhenti.
\end{exercise}

\begin{exercise}
Mari kita ambil \examplevref{sosa:railcarexample}, tetapi pada
waktu tumbukan, kereta 2 bergerak ke kiri dengan kecepatan
\unitfrac[3]{m}{s}.
\begin{tasks}
\task
Temukan perilaku sistem setelah terhubung.
\task
Akankah kereta kedua menabrak
dinding, atau akan bergerak menjauhi dinding seiring berjalannya waktu?
\task
Dengan kecepatan
berapa kereta pertama harus bergerak agar sistem
pada dasarnya tetap di tempat setelah terhubung?
\end{tasks}
\end{exercise}

\begin{exercise}
Mari kita ambil contoh dalam \figurevref{sosa:twocartswallfig} dengan
parameter
$m_1 = m_2 = 1$, $k_1 = k_2 = 1$.  Adakah suatu himpunan syarat
awal yang membuat kereta pertama bergerak tetapi kereta kedua tidak?
Jika ada, temukan syarat tersebut.  Jika tidak, jelaskan alasannya.
\end{exercise}

\setcounter{exercise}{100}

\begin{exercise}
Temukan solusi umum dari
$\left[ \begin{smallmatrix}
1 & 0 & 0\\
0 & 2 & 0\\
0 & 0 & 3
\end{smallmatrix}\right]
\vec{x}\,''
=
\left[ \begin{smallmatrix}
-3 & 0 & 0 \\
2 & -4 & 0 \\
0 & 6 & -3
\end{smallmatrix}\right]
\vec{x}
+ 
\left[ \begin{smallmatrix}
\cos(2t) \\ 0 \\ 0
\end{smallmatrix}\right]$.
\end{exercise}
\exsol{%
$\vec{x}
=
\left[ \begin{smallmatrix}
1 \\ -1 \\ 1
\end{smallmatrix}\right]
\bigl( a_1 \cos (\sqrt{3}\, t)  + b_1 \sin (\sqrt{3}\, t) \bigr)
+
\left[ \begin{smallmatrix}
0 \\ 1 \\ -2
\end{smallmatrix}\right]
\bigl( a_2 \cos (\sqrt{2}\, t)  + b_2 \sin (\sqrt{2}\, t) \bigr)
+
$
\\
$
\left[ \begin{smallmatrix}
0 \\ 0 \\ 1
\end{smallmatrix}\right]
\bigl( a_3 \cos (t)  + b_3 \sin (t) \bigr)
+
\left[ \begin{smallmatrix}
-1 \\ \nicefrac{1}{2} \\ \nicefrac{-1}{3}
\end{smallmatrix}\right]
\cos (2t)$
}

\begin{exercise}
Andaikan terdapat tiga kereta bermassa sama $m$ dan dihubungkan oleh dua pegas dengan
konstanta $k$ (dan tidak terhubung ke dinding).  Susun sistemnya dan temukan
solusi umumnya.
\end{exercise}
\exsol{%
$\left[ \begin{smallmatrix}
m & 0 & 0\\
0 & m & 0\\
0 & 0 & m
\end{smallmatrix}\right]
\vec{x}\,''
=
\left[ \begin{smallmatrix}
-k & k & 0 \\
k & -2k & k \\
0 & k & -k
\end{smallmatrix}\right]
\vec{x}$.
Solusi:
$\vec{x} =
\left[ \begin{smallmatrix}
1 \\ -2 \\ 1
\end{smallmatrix}\right]
\bigl( a_1 \cos (\sqrt{\nicefrac{3k}{m}}\, t)  + b_1 \sin
(\sqrt{\nicefrac{3k}{m}}\, t) \bigr)
\allowbreak
+
\left[ \begin{smallmatrix}
1 \\ 0 \\ -1
\end{smallmatrix}\right]
\bigl( a_2 \cos (\sqrt{\nicefrac{k}{m}}\, t)  + b_2 \sin
(\sqrt{\nicefrac{k}{m}}\, t) \bigr)
+
\left[ \begin{smallmatrix}
1 \\ 1 \\ 1
\end{smallmatrix}\right]
\bigl( a_3 t  + b_3 \bigr).$
}

\begin{exercise}
Andaikan suatu kereta bermassa \unit[2]{kg} terpasang melalui suatu pegas dengan
konstanta $k=1$ pada suatu kereta bermassa \unit[3]{kg}, yang
terpasang ke dinding melalui pegas lain yang juga berkonstanta $k=1$.
Andaikan posisi awal kereta pertama adalah 1 meter dalam
arah positif dari posisi diam, dan kereta kedua mulai pada
posisi diam.  Kereta-kereta tidak bergerak dan kemudian dilepaskan.  Temukan
posisi kereta kedua sebagai fungsi waktu.
\end{exercise}
\exsol{%
$x_2 =
( \nicefrac{2}{5} )
\cos (\sqrt{\nicefrac{1}{6}}\, t)
-
( \nicefrac{2}{5} )
\cos (t)$
%sol
%$\vec{x} =
%\left[ \begin{smallmatrix}
%3 \\ 2
%\end{smallmatrix}\right]
%\bigl( a_1 \cos (\sqrt{\nicefrac{1}{6}}\, t)  + b_1 \sin
%(\sqrt{\nicefrac{1}{6}}\, t) \bigr)
%+
%\left[ \begin{smallmatrix}
%1 \\ -1
%\end{smallmatrix}\right]
%\bigl( a_2 \cos (t)  + b_2 \sin (t) \bigr)$.
}

%%%%%%%%%%%%%%%%%%%%%%%%%%%%%%%%%%%%%%%%%%%%%%%%%%%%%%%%%%%%%%%%%%%%%%%%%%%%%%

\sectionnewpage
\section{Nilai eigen berulang} \label{sec:multeigen}

%mbxINTROSUBSECTION

\sectionnotes{1 atau 1,5 kuliah\EPref{, \S5.5 dalam \cite{EP}}\BDref{,
\S7.8 dalam \cite{BD}}}

Mungkin saja suatu matriks $A$ memiliki beberapa nilai eigen \myquote{berulang}.
Artinya, persamaan karakteristik $\det(A-\lambda I) = 0$ mungkin memiliki
akar berulang.  Hal ini sebenarnya tidak mungkin terjadi
untuk suatu matriks acak.  Jika kita mengambil suatu perturbasi kecil dari $A$ (kita mengubah
entri-entri $A$ sedikit), kita memperoleh suatu matriks dengan nilai eigen berbeda.  Karena
sebarang sistem yang ingin kita selesaikan dalam praktik merupakan aproksimasi terhadap kenyataan
bagaimanapun juga, sehingga mengetahui cara menyelesaikan
kasus-kasus khusus ini tidak mutlak diperlukan.  
Di sisi lain, kasus-kasus ini terkadang muncul dalam penerapan.
Selanjutnya, jika kita memiliki nilai eigen berbeda tetapi sangat berdekatan, perilakunya
serupa dengan nilai eigen berulang, sehingga memahami kasus tersebut
akan memberi kita wawasan mengenai apa yang terjadi.

\subsection{Multiplisitas geometrik}

Ambil matriks diagonal
\begin{equation*}
A =
\begin{bmatrix}
3 & 0 \\ 0 & 3
\end{bmatrix} .
\end{equation*}
$A$ memiliki nilai eigen 3 dengan multiplisitas 2.  Kita menyebut
multiplisitas nilai eigen\index{multiplisitas nilai eigen}
dalam persamaan karakteristik sebagai
\emph{\myindex{multiplisitas aljabar}}.  Dalam kasus ini, juga terdapat 2
vektor eigen yang
bebas linear,
$\left[ \begin{smallmatrix} 1 \\ 0 \end{smallmatrix} \right]$
dan
$\left[ \begin{smallmatrix} 0 \\ 1 \end{smallmatrix} \right]$ yang bersesuaian
dengan nilai eigen 3.  Ini berarti
bahwa apa yang disebut \emph{\myindex{multiplisitas geometrik}}
dari nilai eigen ini juga 2.

Dalam semua
teorema tempat kita mensyaratkan suatu matriks memiliki $n$ nilai eigen berbeda, kita sebenarnya hanya
memerlukan $n$ vektor eigen bebas linear.  Sebagai contoh,
${\vec{x}}' = A\vec{x}$ memiliki solusi umum
\begin{equation*}
\vec{x} = 
c_1 \begin{bmatrix} 1 \\ 0 \end{bmatrix} e^{3t}
+ c_2 \begin{bmatrix} 0 \\ 1 \end{bmatrix} e^{3t} .
\end{equation*}
Kita nyatakan kembali teorema mengenai nilai eigen riil.  Dalam teorema tersebut,
kita akan mengulang nilai eigen menurut multiplisitas (aljabar).  Jadi,
untuk matriks $A$ di atas, kita akan mengatakan bahwa nilai eigennya adalah 3 dan 3.

\begin{theorem}
Andaikan matriks $n \times n$ $A$ 
memiliki $n$ nilai eigen riil (tidak harus berbeda), $\lambda_1$,
$\lambda_2$, \ldots, $\lambda_n$,
dan terdapat $n$ vektor eigen bersesuaian yang bebas linear,
$\vec{v}_1$, $\vec{v}_2$, \ldots, $\vec{v}_n$.  Maka solusi umum untuk 
${\vec{x}}' = A\vec{x}$
dapat dituliskan sebagai
\begin{equation*}
\vec{x} = c_1 \vec{v}_1 e^{\lambda_1 t} +
c_2 \vec{v}_2 e^{\lambda_2 t} + \cdots +
c_n \vec{v}_n e^{\lambda_n t} .
\end{equation*}
\end{theorem}

\emph{Multiplisitas geometrik} dari suatu nilai eigen
sama dengan banyaknya vektor eigen bersesuaian yang bebas linear.
Multiplisitas geometrik selalu kurang dari atau
sama dengan multiplisitas aljabar.  Teorema tersebut menangani kasus
ketika kedua multiplisitas ini sama untuk semua nilai eigen.
Jika untuk suatu nilai eigen, multiplisitas geometrik sama
dengan multiplisitas aljabar, maka kita mengatakan bahwa nilai eigen tersebut
\emph{lengkap\index{nilai eigen lengkap}}.

Dengan kata lain, 
hipotesis teorema dapat dinyatakan dengan mengatakan bahwa jika semua
nilai eigen dari $A$ lengkap, maka terdapat $n$ vektor eigen bebas linear,
sehingga kita memiliki solusi umum yang diberikan.

\medskip

Jika multiplisitas geometrik suatu nilai eigen adalah 2 atau lebih,
maka himpunan vektor eigen bebas linear tidak unik hingga
kelipatan seperti sebelumnya.  Sebagai contoh, untuk matriks diagonal $A =
\left[ \begin{smallmatrix} 3 & 0 \\ 0 & 3 \end{smallmatrix} \right]$,
kita juga dapat memilih vektor eigen
$\left[ \begin{smallmatrix} 1 \\ 1 \end{smallmatrix} \right]$
dan
$\left[ \begin{smallmatrix} 1 \\ -1 \end{smallmatrix} \right]$, atau bahkan
sebarang pasangan dua vektor bebas linear.  Banyaknya vektor eigen
bebas linear yang bersesuaian dengan $\lambda$
adalah banyaknya variabel bebas yang kita peroleh ketika menyelesaikan $A\vec{v} =
\lambda \vec{v}$.  Kita memilih nilai tertentu untuk variabel bebas tersebut untuk
memperoleh vektor eigen.  Jika Anda memilih nilai yang berbeda, Anda mungkin memperoleh vektor eigen
yang berbeda.


\subsection{Nilai eigen defektif}

Jika suatu matriks $n \times n$ memiliki kurang dari $n$ vektor eigen
bebas linear, matriks itu disebut \emph{defisien\index{matriks defisien}}.
Maka terdapat setidaknya
satu nilai eigen dengan multiplisitas geometrik yang kurang dari
multiplisitas
aljabarnya.  Kita menyebut nilai eigen ini \emph{defektif\index{nilai eigen
defektif}},
dan selisih
antara kedua multiplisitas itu kita sebut \emph{\myindex{defek}}.
Dengan kata lain, defek adalah banyaknya vektor eigen yang
\myquote{hilang.}

\begin{example}
Matriks
\begin{equation*}
\begin{bmatrix}
3 & 1 \\ 0 & 3
\end{bmatrix}
\end{equation*}
memiliki nilai eigen 3 dengan multiplisitas aljabar 2.
Karena 
$A-\lambda I = \left[ \begin{smallmatrix} 3 & 1 \\ 0 & 3 \end{smallmatrix} \right]
- 3
\left[ \begin{smallmatrix} 1 & 0 \\ 0 & 1 \end{smallmatrix} \right]
=
\left[ \begin{smallmatrix} 0 & 1 \\ 0 & 0 \end{smallmatrix} \right]
$, untuk menghitung vektor eigen, kita harus menyelesaikan
\begin{equation*}
\begin{bmatrix}
0 & 1 \\ 0 & 0
\end{bmatrix}
\begin{bmatrix}
v_1 \\ v_2
\end{bmatrix}
= \vec{0} .
\end{equation*}
Jadi, $v_2 = 0$.  Semua vektor eigen berbentuk
$\left[ \begin{smallmatrix} v_1 \\ 0 \end{smallmatrix} \right]$.  Sebarang dua
vektor semacam itu bergantung linear, sehingga multiplisitas geometrik
dari nilai eigen tersebut adalah 1.  Oleh karena itu, defeknya adalah 1, dan kita tidak dapat lagi
menerapkan metode nilai eigen secara langsung pada suatu sistem PDB dengan
matriks koefisien seperti itu.

\medskip

Untuk menyelesaikan PDB seperti itu\@, kita memerlukan suatu gagasan baru.
Secara kasar, pengamatan kunci yang akan kita gunakan adalah bahwa
jika $\lambda$ adalah nilai eigen dari $A$ dengan multiplisitas aljabar $m$,
maka kita dapat menemukan $m$ vektor bebas linear tertentu
yang menyelesaikan ${(A-\lambda I)}^k \vec{v} = \vec{0}$ untuk berbagai pangkat
$k$.  Kita akan menyebutnya \emph{\myindex{vektor eigen tergeneralisasi}}.

\medskip

Kita lanjutkan dengan
contoh persamaan ${\vec{x}}' = A\vec{x}$
di mana
$A = \left[ \begin{smallmatrix}
3 & 1 \\ 0 & 3
\end{smallmatrix} \right]$.
Kita menemukan nilai eigen $\lambda=3$ dengan multiplisitas (aljabar) 2 dan defek 1.
Kita menemukan satu vektor eigen 
$\vec{v} = \left[ \begin{smallmatrix} 1 \\ 0 \end{smallmatrix} \right]$.
Kita memiliki satu solusi
\begin{equation*}
\vec{x}_1 = \vec{v} e^{3t} = \begin{bmatrix} 1 \\ 0 \end{bmatrix} e^{3t} .
\end{equation*}
Sekarang kita mengalami kebuntuan; kita tidak memperoleh solusi lain dari vektor eigen standar.  Namun,
kita memerlukan dua solusi bebas linear untuk menemukan solusi umum dari
persamaan tersebut.

Mengikuti semangat akar berulang dari
persamaan karakteristik untuk suatu persamaan tunggal, kita coba solusi lain berbentuk
\begin{equation*}
\vec{x}_2 = ( \vec{v}_2 +  \vec{v}_1 t )\, e^{3t} .
\end{equation*}
Kita diferensialkan untuk memperoleh
\begin{equation*}
{\vec{x}_2}' =
\vec{v}_1 e^{3t} +
3 ( \vec{v}_2 +  \vec{v}_1 t )\, e^{3t}
=
( 3 \vec{v}_2 + \vec{v}_1 )\, e^{3t} +  3 \vec{v}_1 t e^{3t} .
\end{equation*}
Karena kita mengandaikan bahwa $\vec{x}_2$ adalah suatu solusi, ${\vec{x}_2}'$ harus
sama dengan $A \vec{x}_2$. Jadi, kita hitung $A \vec{x}_2$:
\begin{equation*}
A \vec{x}_2 = 
A ( \vec{v}_2 +  \vec{v}_1 t )\, e^{3t}
=
A \vec{v}_2 e^{3t} +  A \vec{v}_1 t e^{3t} .
\end{equation*}
Dengan melihat koefisien $e^{3t}$ dan $t e^{3t}$, kita melihat bahwa
$3 \vec{v}_2 + \vec{v}_1 = A \vec{v}_2$ dan
$3 \vec{v}_1 = A \vec{v}_1$.
Ini berarti bahwa
\begin{equation*}
(A-3I)\vec{v}_2 = \vec{v}_1,
\qquad \text{dan} \qquad
(A-3I)\vec{v}_1 = \vec{0}.
\end{equation*}
Oleh karena itu, $\vec{x}_2$ adalah suatu solusi
jika kedua persamaan ini dipenuhi.
Persamaan kedua dipenuhi jika $\vec{v}_1$ adalah suatu
vektor eigen, dan kita telah menemukan vektor eigen di atas, jadi misalkan
$\vec{v}_1 = 
\left[ \begin{smallmatrix} 1 \\ 0 \end{smallmatrix} \right]$.
Jadi, jika kita dapat menemukan suatu $\vec{v}_2$ yang menyelesaikan
%${(A-3I)}^2\vec{v}_2 = \vec{0}$ dan sedemikian sehingga
$(A-3I)\vec{v}_2 = \vec{v}_1$, maka kita selesai.
Ini hanyalah sekumpulan persamaan linear
yang harus diselesaikan, dan sekarang kita sudah sangat mahir melakukannya.
%
Mari kita selesaikan
$(A-3I)\vec{v}_2 = \vec{v}_1$.  Tuliskan
\begin{equation*}
\begin{bmatrix}
0 & 1 \\ 0 & 0
\end{bmatrix}
\begin{bmatrix}
a \\ b
\end{bmatrix}
=
\begin{bmatrix}
1 \\ 0
\end{bmatrix} .
\end{equation*}
Dengan inspeksi, kita melihat bahwa mengambil $a=0$ (sebenarnya $a$ dapat berupa apa saja) dan
$b=1$ memenuhi tujuan.  Oleh karena itu, kita dapat mengambil $\vec{v}_2 = 
\left[ \begin{smallmatrix} 0 \\ 1 \end{smallmatrix} \right]$.  Solusi umum
kita untuk
${\vec{x}}' = A\vec{x}$ adalah
\begin{equation*}
\vec{x} =
c_1 
\begin{bmatrix}
1 \\ 0
\end{bmatrix}
e^{3t}
+
c_2
\left(
\begin{bmatrix}
0 \\ 1
\end{bmatrix}
+
\begin{bmatrix}
1 \\ 0
\end{bmatrix}
t
\right)
\,
e^{3t}
=
\begin{bmatrix}
c_1 e^{3t}+c_2 te^{3t} \\
c_2 e^{3t}
\end{bmatrix} .
\end{equation*}
Mari kita periksa bahwa kita benar-benar memiliki solusi tersebut.  Pertama,
$x_1' = 
c_1 3 e^{3t}+c_2 e^{3t} + 3 c_2 te^{3t} = 3 x_1 + x_2$.  Benar.  Sekarang,
$x_2' = 3 c_2 e^{3t} = 3x_2$.  Benar.
\end{example}

Dalam contoh tersebut, jika kita substitusikan $(A-3I)\vec{v}_2 = \vec{v}_1$ ke dalam
$(A-3I)\vec{v}_1 = \vec{0}$, kita memperoleh
\begin{equation*}
(A-3I)(A-3I) \vec{v}_2 = \vec{0},
\qquad \text{atau} \qquad
{(A-3I)}^2\vec{v}_2 = \vec{0}.
\end{equation*}
Jika 
$(A-3I) \vec{w} \neq \vec{0}$
dan
${(A-3I)}^2 \vec{w} = \vec{0}$, maka 
$(A-3I) \vec{w}$ adalah suatu vektor eigen, suatu kelipatan dari $\vec{v}_1$.
Dalam kasus $2 \times 2$ ini, ${(A-3I)}^2$ hanyalah matriks nol (latihan).
Jadi, sebarang vektor $\vec{w}$ menyelesaikan
${(A-3I)}^2\vec{w} = \vec{0}$, dan kita hanya memerlukan suatu $\vec{w}$ sedemikian sehingga
$(A-3I)\vec{w} \neq \vec{0}$.  Kemudian kita dapat menggunakan
$\vec{w}$ untuk $\vec{v}_2$ dan $(A-3I)\vec{w}$ untuk $\vec{v}_1$.

Perhatikan bahwa sistem contoh ${\vec{x}}' = A \vec{x}$ ini memiliki solusi yang lebih sederhana karena
$A$ adalah apa yang disebut \emph{\myindex{matriks segitiga atas}}, yaitu,
setiap entri di bawah diagonal bernilai nol.
Khususnya, persamaan untuk $x_2$
tidak bergantung pada $x_1$.  Namun ingat, tidak setiap matriks defektif bersifat
segitiga.

\begin{exercise}
Selesaikan ${\vec{x}}' = \left[ \begin{smallmatrix}
3 & 1 \\ 0 & 3
\end{smallmatrix} \right] \vec{x}$ dengan mula-mula menyelesaikan $x_2$, kemudian
$x_1$ secara terpisah.  Periksa bahwa Anda memperoleh solusi yang sama seperti yang kita peroleh
di atas.
\end{exercise}

Mari kita jelaskan algoritma umumnya.  Andaikan $\lambda$ adalah suatu
nilai eigen dengan multiplisitas 2 dan defek 1.
Mula-mula temukan suatu vektor eigen $\vec{v}_1$ dari $\lambda$.  
Artinya, $\vec{v}_1$ menyelesaikan
$(A-\lambda I)\vec{v}_1  = \vec{0}$.
Kemudian, temukan suatu vektor $\vec{v}_2$ sedemikian sehingga
\begin{equation*}
%{(A-\lambda I)}^2\vec{v}_2 & = \vec{0} , \\
(A-\lambda I)\vec{v}_2 = \vec{v}_1 .
\end{equation*}
Ini memberi kita dua solusi bebas linear
\begin{align*}
\vec{x}_1 & = \vec{v}_1 e^{\lambda t} , \\
\vec{x}_2 & = \left( \vec{v}_2 + \vec{v}_1 t \right) e^{\lambda t} .
\end{align*}

\begin{example}
Perhatikan sistem
\begin{equation*}
\vec{x}' =
\begin{bmatrix}
2 & -5 & 0 \\
0 & 2 & 0 \\
-1 & 4 & 1
\end{bmatrix}
\vec{x} .
\end{equation*}
Hitung nilai eigennya,
\begin{equation*}
0 =
\det(A-\lambda I) = 
\det\left(
\begin{bmatrix}
2-\lambda & -5 & 0 \\
0 & 2-\lambda & 0 \\
-1 & 4 & 1-\lambda
\end{bmatrix}
\right)
= (2-\lambda)^2(1-\lambda) .
\end{equation*}
Nilai eigennya adalah 1 dan 2, dengan 2 memiliki multiplisitas 2.
Kita serahkan kepada pembaca untuk menemukan bahwa
$\left[ \begin{smallmatrix} 0 \\ 0 \\ 1 \end{smallmatrix} \right]$
adalah suatu vektor eigen
untuk nilai eigen $\lambda = 1$.

Kita berfokus pada $\lambda = 2$.  Kita hitung vektor eigennya:
\begin{equation*}
\vec{0} =
(A - 2 I) \vec{v}
=
\begin{bmatrix}
0 & -5 & 0 \\
0 & 0 & 0 \\
-1 & 4 & -1
\end{bmatrix}
\begin{bmatrix}
v_1 \\ v_2 \\ v_3
\end{bmatrix}
.
\end{equation*}
Persamaan pertama menyatakan bahwa $v_2 = 0$, sehingga persamaan terakhir
adalah $-v_1 -v_3 = 0$.  Misalkan $v_3$ menjadi variabel bebas untuk memperoleh
bahwa $v_1 = -v_3$.  Kita dapat mengambil $v_3 = -1$ untuk memperoleh vektor eigen
$\left[ \begin{smallmatrix} 1 \\ 0 \\ -1 \end{smallmatrix} \right]$.
Masalahnya, menetapkan $v_3$ ke nilai apa pun lainnya hanya menghasilkan kelipatan
vektor ini, sehingga kita memiliki defek~1.
Misalkan $\vec{v}_1$ adalah vektor eigen tersebut dan kita mencari
suatu vektor eigen tergeneralisasi $\vec{v}_2$:
\begin{equation*}
(A - 2 I) \vec{v}_2 = \vec{v}_1 , 
\end{equation*}
atau
\begin{equation*}
\begin{bmatrix}
0 & -5 & 0 \\
0 & 0 & 0 \\
-1 & 4 & -1
\end{bmatrix}
\begin{bmatrix}
a \\ b \\ c
\end{bmatrix}
=
\begin{bmatrix}
1 \\ 0 \\ -1
\end{bmatrix} ,
\end{equation*}
di mana kita menggunakan $a$, $b$, $c$ sebagai komponen $\vec{v}_2$ demi kesederhanaan.
Persamaan pertama menyatakan $-5b = 1$, sehingga $b = \nicefrac{-1}{5}$.  Persamaan
kedua tidak menyatakan apa pun.
Persamaan terakhir adalah $-a + 4b - c = -1$, atau
$a + \nicefrac{4}{5} + c = 1$, atau
$a + c = \nicefrac{1}{5}$.  Kita ambil $c$ sebagai variabel bebas, dan kita
memilih $c=0$.  Kita memperoleh
$\vec{v}_2 = \left[ \begin{smallmatrix} \nicefrac{1}{5} \\ \nicefrac{-1}{5}
\\ 0 \end{smallmatrix} \right]$.

Oleh karena itu, solusi umumnya adalah
\begin{equation*}
\vec{x} =
c_1
\begin{bmatrix} 0 \\ 0 \\ 1 \end{bmatrix}
e^t
+
c_2 
\begin{bmatrix} 1 \\ 0 \\ -1 \end{bmatrix}
e^{2t}
+
c_3
\left(
\begin{bmatrix} \nicefrac{1}{5} \\ \nicefrac{-1}{5} \\ 0 \end{bmatrix}
+
\begin{bmatrix} 1 \\ 0 \\ -1 \end{bmatrix}
t
\right)
e^{2t}
.
\end{equation*}
\end{example}

Perangkat ini juga dapat digeneralisasi ke multiplisitas yang lebih tinggi
dan defek yang lebih tinggi.
Kita tidak akan membahas metode ini secara terperinci, tetapi kita uraikan gagasannya.
Andaikan bahwa $A$
memiliki nilai eigen $\lambda$ dengan multiplisitas $m$.
Kita mencari vektor-vektor sedemikian sehingga
\begin{equation*}
{(A - \lambda I)}^k \vec{v} = \vec{0},
\qquad \text{tetapi} \qquad
{(A - \lambda I)}^{k-1} \vec{v} \neq \vec{0}.
\end{equation*}
Vektor semacam itu disebut \emph{\myindex{vektor eigen tergeneralisasi}} (kemudian
$\vec{v}_1 = {(A - \lambda I)}^{k-1} \vec{v}$ adalah suatu vektor eigen).
Untuk
vektor eigen $\vec{v}_1$, terdapat suatu rantai vektor eigen tergeneralisasi
$\vec{v}_2$ hingga $\vec{v}_k$ sedemikian sehingga:
\begin{align*}
(A - \lambda I) \vec{v}_1 & = \vec{0} , \\
(A - \lambda I) \vec{v}_2 & = \vec{v}_1 , \\
& ~~\vdots \\
(A - \lambda I) \vec{v}_k & = \vec{v}_{k-1} .
\end{align*}
Sesungguhnya, setelah Anda menemukan $\vec{v}_k$ sedemikian sehingga
${(A - \lambda I)}^k \vec{v}_k = \vec{0}$ tetapi
${(A - \lambda I)}^{k-1} \vec{v}_k \neq \vec{0}$, Anda menemukan seluruh
rantainya karena Anda dapat menghitung sisanya,
$\vec{v}_{k-1} = (A - \lambda I) \vec{v}_k$,
$\vec{v}_{k-2} = (A - \lambda I) \vec{v}_{k-1}$, dan seterusnya.
Kita bentuk solusi-solusi bebas linear
\begin{align*}
\vec{x}_1 & = \vec{v}_1 e^{\lambda t} , \\
\vec{x}_2 & = ( \vec{v}_2 + \vec{v}_1 t ) \, e^{\lambda t} , \\
& ~~\vdots \\
\vec{x}_k & = \left( \vec{v}_k + \vec{v}_{k-1} t +
\vec{v}_{k-2} \frac{t^2}{2} +
\cdots + \vec{v}_2 \frac{t^{k-2}}{(k-2)!} + \vec{v}_1 \frac{t^{k-1}}{(k-1)!}
\right) \, e^{\lambda t} .
\end{align*}
Ingat bahwa $k! = 1 \cdot 2 \cdot 3 \cdots (k-1) \cdot k$ adalah faktorial.
Jika Anda memiliki suatu nilai eigen dengan multiplisitas geometrik $\ell$,
Anda harus menemukan $\ell$ rantai semacam itu (beberapa di antaranya mungkin
pendek: hanya persamaan vektor eigen tunggal).
Kita lanjutkan hingga membentuk $m$ solusi bebas linear, dengan
$m$ adalah multiplisitas aljabar.
Kita tidak sepenuhnya mengetahui vektor eigen tertentu mana yang berpasangan dengan rantai mana, jadi
mulailah dengan menemukan $\vec{v}_k$ terlebih dahulu untuk rantai terpanjang yang mungkin dan
lanjutkan dari sana.

Sebagai contoh, jika $\lambda$ adalah suatu nilai eigen dari $A$
dengan multiplisitas aljabar 3 dan defek 2, maka selesaikan
\begin{equation*}
(A - \lambda I) \vec{v}_1 = \vec{0} , \qquad
(A - \lambda I) \vec{v}_2 = \vec{v}_1 , \qquad
(A - \lambda I) \vec{v}_3 = \vec{v}_2 .
\end{equation*}
Artinya, temukan $\vec{v}_3$ sedemikian sehingga 
${(A - \lambda I)}^3 \vec{v}_3 = \vec{0}$, tetapi
${(A - \lambda I)}^2 \vec{v}_3 \neq \vec{0}$.
Kemudian Anda selesai karena
$\vec{v}_2 = (A - \lambda I) \vec{v}_3$
dan 
$\vec{v}_1 = (A - \lambda I) \vec{v}_2$.
Ketiga solusi bebas linear tersebut adalah
\begin{equation*}
\vec{x}_1 = \vec{v}_1 e^{\lambda t} , \qquad
\vec{x}_2 = ( \vec{v}_2 + \vec{v}_1 t ) \, e^{\lambda t} , \qquad
\vec{x}_3 = \left( \vec{v}_3 + \vec{v}_2 t +
\vec{v}_{1} \frac{t^2}{2} \right) \, e^{\lambda t} .
\end{equation*}

Di sisi lain, jika $A$ memiliki nilai eigen $\lambda$
dengan multiplisitas aljabar 3 dan defek 1, maka 
selesaikan
\begin{equation*}
(A - \lambda I) \vec{v}_1 = \vec{0} , \qquad
(A - \lambda I) \vec{v}_2 = \vec{0} , \qquad
(A - \lambda I) \vec{v}_3 = \vec{v}_2 .
\end{equation*}
Di sini, $\vec{v}_1$ dan $\vec{v}_2$ benar-benar merupakan vektor eigen,
dan $\vec{v}_3$ adalah suatu vektor eigen tergeneralisasi.
Jadi, terdapat dua rantai.
Untuk menyelesaikannya, mula-mula temukan suatu 
$\vec{v}_3$ sedemikian sehingga 
${(A - \lambda I)}^2 \vec{v}_3 = \vec{0}$, tetapi
$(A - \lambda I) \vec{v}_3 \neq \vec{0}$.
Kemudian $\vec{v}_2 = (A - \lambda I) \vec{v}_3$ akan menjadi suatu vektor eigen.
Lalu selesaikan suatu vektor eigen $\vec{v}_1$ yang bebas linear 
dari $\vec{v}_2$.
Anda memperoleh 3 solusi bebas linear
\begin{equation*}
\vec{x}_1 = \vec{v}_1 e^{\lambda t} , \qquad
\vec{x}_2 = \vec{v}_2 e^{\lambda t} , \qquad
\vec{x}_3 = ( \vec{v}_3 + \vec{v}_2 t ) \, e^{\lambda t} .
\end{equation*}

\subsection{Latihan}

\begin{exercise}
Misalkan
$A = \left[ \begin{smallmatrix} 5 & -3 \\ 3 & -1 \end{smallmatrix} \right]$.
Temukan solusi umum dari ${\vec{x}}' = A \vec{x}$.
\end{exercise}

\begin{samepage}
\begin{exercise}
Misalkan
$A = \left[ \begin{smallmatrix}
5 & -4 & 4 \\
0 & 3 & 0 \\
-2 & 4 & -1
\end{smallmatrix} \right]$.
\begin{tasks}
\task Apakah nilai eigennya?
\task Apakah defek dari nilai eigennya?
\task Temukan solusi umum dari ${\vec{x}}' = A \vec{x}$.
\end{tasks}
\end{exercise}
\end{samepage}


\begin{exercise}
Misalkan
$A = \left[ \begin{smallmatrix} 2 & 1 & 0 \\ 0 & 2 & 0 \\ 0 & 0 & 2 \end{smallmatrix} \right]$.
\begin{tasks}
\task Apakah nilai eigennya?
\task Apakah defek dari nilai eigennya?
\task Temukan solusi umum dari ${\vec{x}}' = A \vec{x}$ dengan dua cara berbeda,
dan verifikasi bahwa Anda memperoleh jawaban yang sama.
\end{tasks}
\end{exercise}

\begin{exercise}
Misalkan
$A = \left[ \begin{smallmatrix}
0 & 1 & 2 \\
-1 & -2 & -2 \\
-4 & 4 & 7
\end{smallmatrix} \right]$.
\begin{tasks}
\task Apakah nilai eigennya?
\task Apakah defek dari nilai eigennya?
\task Temukan solusi umum dari ${\vec{x}}' = A \vec{x}$.
\end{tasks}
\end{exercise}

\begin{exercise}
\pagebreak[2]
Misalkan
$A = \left[ \begin{smallmatrix}
0 & 4 & -2 \\
-1 & -4 & 1 \\
0 & 0 & -2
\end{smallmatrix} \right]$.
\begin{tasks}
\task Apakah nilai eigennya?
\task Apakah defek dari nilai eigennya?
\task Temukan solusi umum dari ${\vec{x}}' = A \vec{x}$.
\end{tasks}
\end{exercise}

\begin{exercise}
Misalkan
$A = \left[ \begin{smallmatrix}
2 & 1 & -1 \\
-1 & 0 & 2 \\
-1 & -2 & 4
\end{smallmatrix} \right]$.
\begin{tasks}
\task Apakah nilai eigennya?
\task Apakah defek dari nilai eigennya?
\task Temukan solusi umum dari ${\vec{x}}' = A \vec{x}$.
\end{tasks}
\end{exercise}

\begin{exercise}
Andaikan $A$ adalah suatu matriks $2 \times 2$ dengan nilai eigen berulang
$\lambda$.
Andaikan terdapat dua vektor eigen bebas linear.  Tunjukkan bahwa
$A = \lambda I$.
\end{exercise}

\setcounter{exercise}{100}

\begin{exercise}
Misalkan $A =
\left[ \begin{smallmatrix}
1 & 1 & 1 \\
1 & 1 & 1 \\
1 & 1 & 1 
\end{smallmatrix}\right]$.  
\begin{tasks}
\task Apakah nilai eigennya?
\task Apakah defek dari nilai eigennya?
\task Temukan solusi umum dari $\vec{x}\,' = A\vec{x}$.
\end{tasks}
\end{exercise}
\exsol{%
a) $3,0,0$
\quad
b) Tidak ada defek.
\quad
c)
$\vec{x} =
C_1
\left[ \begin{smallmatrix}
1 \\ 1 \\ 1
\end{smallmatrix}\right]
e^{3t}
+
C_2
\left[ \begin{smallmatrix}
1 \\ 0 \\ -1
\end{smallmatrix}\right]
+
C_3
\left[ \begin{smallmatrix}
0 \\ 1 \\ -1
\end{smallmatrix}\right]$
}

\begin{exercise}
Misalkan $A =
\left[ \begin{smallmatrix}
1 & 3 & 3 \\
1 & 1 & 0 \\
-1 & 1 & 2 \\
\end{smallmatrix}\right]$.  
\begin{tasks}
\task Apakah nilai eigennya?
\task Apakah defek dari nilai eigennya?
\task Temukan solusi umum dari $\vec{x}\,' = A\vec{x}$.
\end{tasks}
\end{exercise}
\exsol{%
a) $1,1,2$
\\
b) Nilai eigen 1 memiliki defek 1
\\
c)
$\vec{x} =
C_1
\left[ \begin{smallmatrix}
0 \\ 1 \\ -1
\end{smallmatrix}\right]
e^{t}
+
C_2
\left(
\left[ \begin{smallmatrix}
1 \\ 0 \\ 0
\end{smallmatrix}\right]
+
t
\left[ \begin{smallmatrix}
0 \\ 1 \\ -1
\end{smallmatrix}\right]
\right)
e^{t}
+
C_3
\left[ \begin{smallmatrix}
3 \\ 3 \\ -2
\end{smallmatrix}\right]
e^{2t}$
}

\begin{exercise}
Misalkan $A =
\left[ \begin{smallmatrix}
2 & 0 & 0 \\
-1 & -1 & 9 \\
0 & -1 & 5
\end{smallmatrix}\right]$.  
\begin{tasks}
\task Apakah nilai eigennya?
\task Apakah defek dari nilai eigennya?
\task Temukan solusi umum dari $\vec{x}\,' = A\vec{x}$.
\end{tasks}
\end{exercise}
\exsol{%
a) $2,2,2$
\\
b) Nilai eigen 2 memiliki defek 2
\\
c)
$\vec{x} =
C_1
\left[ \begin{smallmatrix}
0 \\ 3 \\ 1
\end{smallmatrix}\right]
e^{2t}
+
C_2
\left(
\left[ \begin{smallmatrix}
0 \\ -1 \\ 0
\end{smallmatrix}\right]
+
t
\left[ \begin{smallmatrix}
0 \\ 3 \\ 1
\end{smallmatrix}\right]
\right)
e^{2t}
+
C_3
\left(
\left[ \begin{smallmatrix}
1 \\ 0 \\ 0
\end{smallmatrix}\right]
+
t
\left[ \begin{smallmatrix}
0 \\ -1 \\ 0
\end{smallmatrix}\right]
+
\frac{t^2}{2}
\left[ \begin{smallmatrix}
0 \\ 3 \\ 1
\end{smallmatrix}\right]
\right)
e^{2t}$
}

\begin{exercise}
Misalkan $A =
\left[ \begin{smallmatrix}
a & a \\
b & c
\end{smallmatrix}\right]$, dengan $a$, $b$, dan $c$ tidak diketahui.
Misalkan $5$ adalah nilai eigen ganda dengan defek 1, dan misalkan
$\left[ \begin{smallmatrix}
1 \\ 0
\end{smallmatrix}\right]$ adalah vektor eigen yang bersesuaian.  Carilah $A$ dan tunjukkan bahwa
hanya ada satu matriks $A$ demikian.
\end{exercise}
\exsol{%
$A = \left[ \begin{smallmatrix}
5 & 5 \\ 0 & 5
\end{smallmatrix}\right]$
}

%%%%%%%%%%%%%%%%%%%%%%%%%%%%%%%%%%%%%%%%%%%%%%%%%%%%%%%%%%%%%%%%%%%%%%%%%%%%%%

\sectionnewpage
\section{Eksponensial matriks} \label{sec:matexp}

%mbxINTROSUBSECTION

\sectionnotes{2 kuliah, mungkin perlu merujuk ke \sectionref{powerseries:section}\EPref{, \S5.6 dalam \cite{EP}}\BDref{, \S7.7 dalam \cite{BD}}}

\subsection{Definisi}

Ada cara lain untuk mencari solusi matriks fundamental
dari suatu sistem.
Tinjau persamaan berkoefisien konstan
\begin{equation*}
{\vec{x}}' = P \vec{x} .
\end{equation*}
Jika ini hanya satu persamaan (ketika $P$ adalah suatu bilangan atau matriks
$1 \times 1$), maka solusinya adalah
\begin{equation*}
\vec{x} = e^{Pt} .
\end{equation*}
Hal itu (belum) bermakna jika $P$ adalah matriks yang lebih besar, tetapi
pada dasarnya perhitungan yang sama yang menghasilkan bentuk di atas
berlaku untuk matriks apabila kita mendefinisikan
$e^{Pt}$ dengan tepat.  Pertama, kita tuliskan deret Taylor untuk $e^{at}$
bagi suatu bilangan $a$:
\begin{equation*}
e^{at} = 1 + at
+ \frac{{(at)}^2}{2}
+ \frac{{(at)}^3}{6}
+ \frac{{(at)}^4}{24}
+ \cdots
= \sum_{k=0}^\infty \frac{{(at)}^k}{k!} .
\end{equation*}
Ingat bahwa $k! = 1 \cdot 2 \cdot 3 \cdots k$ adalah faktorial, dan $0! = 1$.
Kita turunkan deret ini suku demi suku
\begin{equation*}
%mbxlatex \begin{split}
\frac{d}{dt} \left(e^{at} \right)
%mbxlatex &
= 
0
+ a
+ a^2 t
+ \frac{a^3t^2}{2}
+ \frac{a^4t^3}{6}
+ \cdots
%mbxlatex \\
%mbxlatex &
=
a \left(
1
+ a t
+ \frac{{(at)}^2}{2}
+ \frac{{(at)}^3}{6}
+ \cdots \right)
= a e^{at}.
%mbxlatex \end{split}
\end{equation*}
Mungkin kita dapat mencoba cara yang sama untuk matriks.  Untuk suatu matriks
berukuran $n \times n$, yaitu $A$, kita definisikan
\emph{\myindex{eksponensial matriks}\index{eksponensial matriks}} sebagai
\begin{equation*}
\mybxbg{~~
e^A \overset{\text{def}}{=} I + A + \frac{1}{2} A^2 + 
\frac{1}{6} A^3 + \cdots + \frac{1}{k!} A^k + \cdots
~~}
\end{equation*}
Kita tidak akan mempersoalkan konvergensi---deret tersebut memang
selalu konvergen.
Menurut konvensi, kita biasanya menuliskan $Pt$ sebagai $tP$ ketika $P$ adalah matriks.
Dengan perubahan kecil ini dan perhitungan yang persis sama
seperti di atas,
kita memperoleh
\begin{equation*}
\frac{d}{dt} \left(e^{tP} \right) = 
P e^{tP} .
\end{equation*}
Sekarang $P$, dan karena itu $e^{tP}$, adalah matriks $n \times n$.  Yang kita cari
adalah suatu vektor.  Dalam kasus $1 \times 1$, pada tahap ini kita
mengalikannya dengan konstanta sebarang untuk memperoleh solusi umum.  Dalam kasus
matriks, kita mengalikannya dengan vektor kolom $\vec{c}$.

\begin{theorem}
Misalkan $P$ matriks $n \times n$.  Maka solusi umum dari
${\vec{x}}' = P \vec{x}$ adalah
\begin{equation*}
\vec{x} = e^{tP} \vec{c} ,
\end{equation*}
dengan $\vec{c}$ suatu vektor konstan sebarang.  Bahkan, $\vec{x}(0) =
\vec{c}$.
\end{theorem}

Pemeriksaan:
\begin{equation*}
\frac{d}{dt}
\vec{x} =
\frac{d}{dt} \left( 
e^{tP} \vec{c}\, \right)
=
P e^{tP} \vec{c} =
P \vec{x}.
\end{equation*}

Jadi $e^{tP}$ merupakan suatu \myindex{solusi matriks fundamental}
dari sistem homogen.
Jika kita dapat menghitung eksponensial matriks,
kita memperoleh metode lain untuk menyelesaikan sistem homogen
berkoefisien konstan.  Metode ini juga memudahkan penyelesaian syarat awal.  Untuk menyelesaikan 
${\vec{x}}' = P \vec{x}$, $\vec{x}(0) = \vec{b}$, kita ambil solusi
\begin{equation*}
\vec{x} = e^{tP} \vec{b} .
\end{equation*}
Persamaan ini berlaku karena $e^{0P} = I$,
sehingga $\vec{x} (0) = e^{0P} \vec{b} = \vec{b}$.

\medskip

Kita sebutkan satu kelemahan eksponensial matriks.
Secara umum $e^{A+B} \neq e^A e^B$.  Masalahnya adalah matriks-matriks tidak
komutatif, yakni secara umum $AB \neq BA$.
Jika Anda mencoba menuliskan kembali $e^A e^B$ sebagai $e^{A+B}$ dengan deret Taylor,
Anda akan melihat mengapa ketidakkomutatifan menjadi masalah.
Memang benar bahwa jika
$AB = BA$, yakni jika $A$ dan $B$ komutatif, maka $e^{A+B} = e^Ae^B$.  Fakta ini akan
berguna bagi kita.  Kita nyatakan ulang sebagai teorema untuk menegaskan hal tersebut.

\begin{theorem}
Jika $AB = BA$, maka $e^{A+B} = e^Ae^B$.  Jika tidak,
secara umum $e^{A+B} \neq e^Ae^B$.
\end{theorem}

\subsection{Kasus-kasus sederhana}

Dalam beberapa keadaan, kita dapat langsung menyubstitusikan ke definisi deret.
Misalkan matriksnya diagonal\index{matriks diagonal!eksponensial matriks dari}.
Sebagai contoh,
$D = \left[ \begin{smallmatrix} a & 0 \\ 0 & b \end{smallmatrix} \right]$.
Maka 
\begin{equation*}
D^k = \begin{bmatrix} a^k & 0 \\ 0 & b^k \end{bmatrix} ,
\end{equation*}
dan
\begin{equation*}
\begin{split}
e^D & =
I + D + \frac{1}{2} D^2 + 
\frac{1}{6} D^3 + \cdots
\\
&=
\begin{bmatrix} 1 & 0 \\ 0 & 1 \end{bmatrix} +
\begin{bmatrix} a & 0 \\ 0 & b \end{bmatrix} +
\frac{1}{2}
\begin{bmatrix} a^2 & 0 \\ 0 & b^2 \end{bmatrix} +
\frac{1}{6}
\begin{bmatrix} a^3 & 0 \\ 0 & b^3 \end{bmatrix} + \cdots
=
\begin{bmatrix} e^a & 0 \\ 0 & e^b \end{bmatrix} .
\end{split}
\end{equation*}
Secara khusus,
\begin{equation*}
e^I = \begin{bmatrix} e & 0\\ 0 & e \end{bmatrix}
\qquad \text{dan} \qquad
e^{aI} = \begin{bmatrix} e^a & 0\\ 0 & e^a \end{bmatrix}.
\end{equation*}

Hal ini membuat eksponensial matriks-matriks tertentu lainnya mudah dihitung.  
Sebagai contoh, matriks
$A = \left[ \begin{smallmatrix} 5 & 4 \\ -1 & 1 \end{smallmatrix} \right]$
dapat dituliskan sebagai
$3I + B$ dengan
$B = \left[ \begin{smallmatrix} 2 & 4 \\ -1 & -2 \end{smallmatrix} \right]$.
Perhatikan bahwa $B^2 = 
\left[ \begin{smallmatrix} 0 & 0 \\ 0 & 0 \end{smallmatrix} \right]$.  Jadi
$B^k = 0$ untuk semua $k \geq 2$.  Oleh karena itu, $e^B = I + B$.  Misalkan kita
benar-benar ingin menghitung $e^{tA}$.  Matriks $3tI$ dan $tB$ komutatif
(latihan: periksalah hal ini)
dan $e^{tB} = I + tB$, karena ${(tB)}^2 = t^2 B^2 = 0$.
Kita tuliskan
\begin{multline*}
e^{tA} = 
e^{3tI + tB} = e^{3tI} e^{tB} = 
\begin{bmatrix} e^{3t} & 0 \\ 0 & e^{3t} \end{bmatrix}
\left(
I + tB
\right)
=
\\
=
\begin{bmatrix} e^{3t} & 0 \\ 0 & e^{3t} \end{bmatrix}
\begin{bmatrix} 1+2t & 4t \\ -t & 1-2t \end{bmatrix}
=
\begin{bmatrix} (1+2t)\,e^{3t} & 4te^{3t} \\ -te^{3t} & (1-2t)\,e^{3t} \end{bmatrix} .
\end{multline*}
Kita telah menemukan solusi matriks fundamental untuk
sistem ${\vec{x}}' = A \vec{x}$.  Perhatikan bahwa matriks ini memiliki nilai eigen
berulang dengan defek; hanya ada satu vektor eigen untuk nilai eigen
3.  Jadi kita menemukan cara yang mungkin lebih mudah untuk menangani kasus ini.  Faktanya, jika
matriks $A$ berukuran $2 \times 2$ dan memiliki nilai eigen $\lambda$ bermultiplisitas
2, maka
entah $A = \lambda I$, atau $A = \lambda I + B$ dengan $B^2 = 0$.  Ini merupakan
latihan yang baik.

\begin{exercise}%[menantang] Tidak lagi menantang dengan petunjuk ini?
Misalkan $A$ berukuran $2 \times 2$ dan $\lambda$
adalah satu-satunya nilai eigen.  Tunjukkan bahwa ${(A - \lambda I)}^2 = 0$, dan dengan demikian
kita dapat menuliskan $A = \lambda I + B$, dengan $B^2 = 0$ (dan mungkin $B=0$).
Petunjuk: Pertama,
tuliskan arti nilai eigen bermultiplisitas 2.
Anda akan memperoleh
suatu persamaan untuk entri-entrinya.  Sekarang hitung kuadrat $B$.
\end{exercise}

Matriks $B$ sedemikian sehingga $B^k = 0$ untuk suatu $k$ disebut
\emph{\myindex{nilpoten}}.  Perhitungan eksponensial matriks untuk
matriks nilpoten mudah dilakukan dengan hanya menuliskan $k$ suku pertama
deret Taylor.

\subsection{Matriks umum}

Secara umum, eksponensial tidak semudah di atas untuk dihitung.  Biasanya kita
tidak dapat menuliskan suatu matriks sebagai jumlah matriks-matriks komutatif yang masing-masing
memiliki eksponensial sederhana.  Namun jangan khawatir, persoalannya tetap tidak terlalu sulit asalkan
kita dapat menemukan cukup banyak vektor eigen.  Pertama, kita memerlukan hasil menarik berikut
tentang eksponensial matriks.  Untuk dua matriks persegi $A$ dan $B$, dengan
$B$ invertibel,
kita memiliki
\begin{equation*}
e^{BAB^{-1}} = B e^A B^{-1} .
\end{equation*}
Hal ini dapat dilihat dengan menuliskan deret Taylor.  Pertama, 
\begin{equation*}
{(BAB^{-1})}^2 =
BAB^{-1} BAB^{-1} =
BAIAB^{-1} =
BA^2B^{-1} .
\end{equation*}
Dengan penalaran yang sama, ${(BAB^{-1})}^k = B A^k B^{-1}$.  Sekarang tuliskan 
deret Taylor untuk
$e^{BAB^{-1}}$:
\begin{equation*}
\begin{split}
e^{BAB^{-1}} & =
I + {BAB^{-1}} + \frac{1}{2} {(BAB^{-1})}^2 + 
\frac{1}{6} {(BAB^{-1})}^3 + \cdots
\\
& =
BB^{-1} + {BAB^{-1}} + \frac{1}{2} BA^2B^{-1} + 
\frac{1}{6} BA^3B^{-1} + \cdots
\\
& =
B \bigl(
I + A + \frac{1}{2} A^2 + 
\frac{1}{6} A^3 + \cdots \bigr) B^{-1} \\
& = B e^A B^{-1} .
\end{split}
\end{equation*}

Untuk suatu matriks persegi $A$, 
biasanya kita dapat menuliskan $A = E D E^{-1}$, dengan $D$
diagonal dan $E$ invertibel.
Prosedur ini disebut \emph{\myindex{diagonalisasi}}.
Jika kita dapat melakukan
hal tersebut, perhitungan
eksponensial menjadi mudah karena $e^D$ diperoleh hanya dengan mengambil eksponensial
entri-entri pada diagonal.  Dengan menyertakan $t$, 
kita kemudian dapat menghitung eksponensial
\begin{equation*}
e^{tA} = E e^{tD} E^{-1} .
\end{equation*}

Untuk mendiagonalkan $A$,
kita memerlukan $n$ vektor eigen $A$ yang bebas linear.
Jika tidak, metode perhitungan eksponensial ini tidak berlaku dan kita memerlukan cara yang lebih cerdik, tetapi kita
tidak akan membahas rincian tersebut.
Misalkan $E$ matriks yang kolom-kolomnya adalah vektor eigen.  
Misalkan $\lambda_1$, $\lambda_2$, \ldots, $\lambda_n$ nilai-nilai eigen dan
$\vec{v}_1$, $\vec{v}_2$, \ldots, $\vec{v}_n$ vektor-vektor eigennya, maka
$E = [\, \vec{v}_1 \quad \vec{v}_2 \quad \cdots \quad \vec{v}_n \,]$.
Bentuklah matriks diagonal $D$ dengan nilai-nilai eigen pada diagonalnya:
\begin{equation*}
D =
\begin{bmatrix}
\lambda_1 & 0 & \cdots & 0 \\
0 & \lambda_2 & \cdots & 0 \\
\vdots & \vdots & \ddots & \vdots \\
0 & 0 & \cdots & \lambda_n
\end{bmatrix} .
\end{equation*}
Kita hitung
\begin{equation*}
\begin{split}
AE & = A 
[\, \vec{v}_1 \quad \vec{v}_2 \quad \cdots \quad \vec{v}_n \,]
\\
& =
[\, A\vec{v}_1 \quad A\vec{v}_2 \quad \cdots \quad A\vec{v}_n \,]
\\
& =
[\, \lambda_1 \vec{v}_1 \quad \lambda_2 \vec{v}_2 \quad \cdots \quad
\lambda_n \vec{v}_n \,]
\\
& =
[\, \vec{v}_1 \quad \vec{v}_2 \quad \cdots \quad \vec{v}_n \,] D
\\
& =
ED .
\end{split}
\end{equation*}
Kolom-kolom $E$ bebas linear karena merupakan
vektor-vektor eigen $A$ yang
bebas linear.  Jadi $E$ invertibel.
Karena $AE = ED$, kita kalikan di sebelah kanan dengan $E^{-1}$ dan memperoleh
\begin{equation*}
A = E D E^{-1}.
\end{equation*}
Ini berarti $e^A = E e^D E^{-1}$.  Dengan mengalikan matriks tersebut dengan $t$, kita memperoleh
\begin{equation} \label{matexp:diagfundsol}
\mybxbg{~~
e^{tA} = 
Ee^{tD}E^{-1} = 
E
\begin{bmatrix}
e^{\lambda_1 t} & 0 & \cdots & 0 \\
0 & e^{\lambda_2 t} & \cdots & 0 \\
\vdots & \vdots & \ddots & \vdots \\
0 & 0 & \cdots & e^{\lambda_n t}
\end{bmatrix} 
E^{-1} .
~~}
\end{equation}
Jadi, rumus \eqref{matexp:diagfundsol} memberikan rumus
untuk menghitung solusi matriks fundamental $e^{tA}$ bagi
sistem ${\vec{x}}' = A \vec{x}$ dalam kasus kita memiliki
$n$ vektor eigen yang bebas linear.

Perhitungan ini tetap berlaku ketika nilai eigen dan
vektor eigennya kompleks, meskipun saat itu Anda harus menghitung dengan bilangan
kompleks.  Dari definisi jelas bahwa jika $A$ real,
maka $e^{tA}$ real.  Jadi bilangan kompleks hanya diperlukan dalam
perhitungan, bukan dalam hasilnya.  Anda mungkin perlu menerapkan
\hyperref[eulersformula]{rumus Euler} untuk menyederhanakan
hasil.  Jika disederhanakan dengan benar, matriks akhir tidak akan memuat bilangan
kompleks.

\begin{example}
Hitunglah solusi matriks fundamental dengan menggunakan eksponensial matriks
untuk sistem
\begin{equation*}
\begin{bmatrix}
x \\ y
\end{bmatrix} '
=
\begin{bmatrix}
1 & 2 \\
2 & 1
\end{bmatrix}
\begin{bmatrix}
x \\ y
\end{bmatrix} .
\avoidbreak
\end{equation*}
Kemudian hitung solusi khusus untuk syarat awal
$x(0) = 4$ dan $y(0) = 2$.

Misalkan $A$ matriks koefisien $\left[ \begin{smallmatrix}
1 & 2 \\
2 & 1
\end{smallmatrix} \right]$.
Pertama kita hitung (latihan) bahwa nilai eigennya adalah 3 dan $-1$, serta 
vektor eigen yang bersesuaian adalah
$\left[ \begin{smallmatrix} 1 \\ 1 \end{smallmatrix} \right]$ dan
$\left[ \begin{smallmatrix} 1 \\ -1 \end{smallmatrix} \right]$.
Dengan demikian, diagonalisasi $A$ adalah
\begin{equation*}
\underbrace{
\begin{bmatrix}
1 & 2 \\
2 & 1
\end{bmatrix}
}_{A}
=
\underbrace{
\begin{bmatrix}
1 & 1 \\
1 & -1
\end{bmatrix}
}_{E}
\underbrace{
\begin{bmatrix}
3 & 0 \\
0 & -1
\end{bmatrix}
}_{D}
\underbrace{
\begin{bmatrix}
1 & 1 \\
1 & -1
\end{bmatrix}^{-1}
}_{E^{-1}} .
\end{equation*}
Kita tuliskan
\begin{equation*}
\begin{split}
e^{t A}
=
E e^{tD} E^{-1}
& =
\begin{bmatrix}
1 & 1 \\
1 & -1
\end{bmatrix}
\begin{bmatrix}
e^{3t} & 0 \\
0 & e^{-t}
\end{bmatrix}
\begin{bmatrix}
1 & 1 \\
1 & -1
\end{bmatrix}^{-1}
\\
& =
\begin{bmatrix}
1 & 1 \\
1 & -1
\end{bmatrix}
\begin{bmatrix}
e^{3t} & 0 \\
0 & e^{-t}
\end{bmatrix}
\frac{-1}{2}
\begin{bmatrix}
-1 & -1 \\
-1 & 1
\end{bmatrix}
\\
& =
\frac{-1}{2}
\begin{bmatrix}
e^{3t} & e^{-t} \\
e^{3t} & -e^{-t}
\end{bmatrix}
\begin{bmatrix}
-1 & -1 \\
-1 & 1
\end{bmatrix}
\\
& =
\frac{-1}{2}
\begin{bmatrix}
-e^{3t}-e^{-t} & -e^{3t}+e^{-t} \\
-e^{3t}+e^{-t} & -e^{3t}-e^{-t}
\end{bmatrix}
 =
\begin{bmatrix}
\frac{e^{3t}+e^{-t}}{2} & \frac{e^{3t}-e^{-t}}{2} \\
\frac{e^{3t}-e^{-t}}{2} & \frac{e^{3t}+e^{-t}}{2}
\end{bmatrix} .
\end{split}
\end{equation*}

Syarat awalnya adalah $x(0) = 4$ dan $y(0) = 2$.  Jadi, berdasarkan sifat
$e^{0A} = I$, kita memperoleh bahwa solusi khusus yang kita cari
adalah $e^{tA} \vec{b}$ dengan $\vec{b}$ adalah $\left[
\begin{smallmatrix} 4 \\ 2 \end{smallmatrix} \right]$.
Solusi khusus yang kita cari adalah
\begin{equation*}
\begin{bmatrix}
x \\ y
\end{bmatrix}
=
\begin{bmatrix}
\frac{e^{3t}+e^{-t}}{2} & \frac{e^{3t}-e^{-t}}{2} \\
\frac{e^{3t}-e^{-t}}{2} & \frac{e^{3t}+e^{-t}}{2}
\end{bmatrix}
\begin{bmatrix}
4 \\ 2
\end{bmatrix}
=
\begin{bmatrix}
2e^{3t}+2e^{-t} + e^{3t}-e^{-t} \\
2e^{3t}-2e^{-t} + e^{3t}+e^{-t}
\end{bmatrix}
=
\begin{bmatrix}
3e^{3t}+e^{-t} \\
3e^{3t}-e^{-t}
\end{bmatrix} .
\end{equation*}
\end{example}

\subsection{Solusi matriks fundamental}

Jika Anda dapat menghitung solusi matriks fundamental
dengan cara lain, Anda dapat menggunakannya untuk mencari eksponensial matriks $e^{tA}$.
Solusi matriks fundamental suatu sistem PDB tidak tunggal.  Eksponensial tersebut
adalah solusi matriks fundamental yang memiliki sifat bahwa
untuk $t=0$ kita memperoleh matriks identitas.  Jadi kita
harus mencari solusi matriks fundamental yang tepat.  Misalkan $X$ sebarang solusi
matriks fundamental dari ${\vec{x}}' = A \vec{x}$.  Kita klaim
\begin{equation*}
e^{tA} = X(t) \left[ X(0) \right]^{-1} .
\end{equation*}
Jika kita menyubstitusikan $t=0$ ke 
$X(t) \left[ X(0) \right]^{-1}$, kita memperoleh identitas sebagaimana diperlukan.
Kita dapat mengalikan suatu solusi matriks fundamental di sebelah kanan dengan sebarang
matriks konstan invertibel dan tetap memperoleh solusi matriks fundamental.
Yang kita lakukan hanyalah mengubah konstanta-konstanta sebarang dalam solusi
umum $\vec{x}(t) = X(t)\, \vec{c}$.

\subsection{Hampiran}

Jika dipikirkan, perhitungan sebarang solusi matriks fundamental
$X$ dengan metode nilai eigen
sama sulitnya dengan perhitungan $e^{tA}$.
Jadi mungkin kita tidak memperoleh
banyak keuntungan dari alat baru ini.  Namun, ekspansi deret Taylor sebenarnya memberi
kita cara untuk menghampiri solusi, sesuatu yang tidak diberikan
oleh metode nilai eigen.

Hal paling sederhana yang dapat kita lakukan
adalah menghitung deret sampai sejumlah suku tertentu.  Ada
cara yang lebih baik untuk menghampiri eksponensial%
\footnote{C.\ Moler dan C.F.\ Van Loan, \emph{Nineteen Dubious Ways
to Compute the Exponential of a Matrix, Twenty-Five Years Later}, SIAM Review
45 (1), 2003, 3--49}.  Namun, dalam banyak kasus,
beberapa suku deret Taylor memberikan
hampiran yang wajar bagi eksponensial dan mungkin cukup untuk
penerapannya.  Sebagai contoh, mari kita hitung 4 suku pertama
deret untuk matriks $A = 
\left[ \begin{smallmatrix}
1 & 2 \\
2 & 1
\end{smallmatrix} \right]$.
\begin{multline*}
e^{tA}
\approx
I + tA + \frac{t^2}{2}A^2 + \frac{t^3}{6}A^3
=
I + t
\begin{bmatrix}
1 & 2 \\
\noalign{\smallskip}
2 & 1
\end{bmatrix}
+ t^2
\begin{bmatrix}
\frac{5}{2} & 2 \\
\noalign{\smallskip}
2 & \frac{5}{2}
\end{bmatrix}
+ t^3
\begin{bmatrix}
\frac{13}{6} & \frac{7}{3} \\
\noalign{\smallskip}
\frac{7}{3} & \frac{13}{6}
\end{bmatrix}
=
\\
=
\begin{bmatrix}
1 + t + \frac{5}{2}\, t^2 + \frac{13}{6}\, t^3 &
   2\,t + 2\, t^2   + \frac{7}{3}\, t^3 \\
\noalign{\smallskip}
   2\,t + 2\, t^2   + \frac{7}{3}\, t^3 &
1 + t + \frac{5}{2}\, t^2 + \frac{13}{6}\, t^3
\end{bmatrix} .
\end{multline*}
Sama seperti versi skalar dari hampiran deret Taylor,
hampiran akan lebih baik untuk $t$ kecil dan lebih buruk untuk $t$ yang lebih besar.  Untuk
$t$ yang lebih besar, kita harus menghitung lebih banyak suku.
Mari kita bandingkan hasilnya dengan solusi sebenarnya untuk $t=0.1$.  Solusi
hampirannya adalah kira-kira (dibulatkan hingga 8 tempat desimal)
\begin{equation*}
e^{0.1\,A} \approx
I + 0.1\,A + \frac{0.1^2}{2}A^2 + \frac{0.1^3}{6}A^3
=
\begin{bmatrix}
1.12716667 & 0.22233333 \\
0.22233333 & 1.12716667 \\
\end{bmatrix} .
\end{equation*}
Dengan menyubstitusikan $t=0.1$ ke solusi sebenarnya (dibulatkan hingga 8 tempat desimal), kita memperoleh
\begin{equation*}
e^{0.1\,A} = 
\begin{bmatrix}
1.12734811 & 0.22251069 \\
0.22251069 & 1.12734811
\end{bmatrix} .
\end{equation*}
Sama sekali tidak buruk!  Namun, jika kita menggunakan hampiran yang sama untuk $t=1$,
kita memperoleh
\begin{equation*}
I + A + \frac{1}{2}A^2 + \frac{1}{6}A^3
=
\begin{bmatrix}
6.66666667 & 6.33333333 \\
6.33333333 & 6.66666667
\end{bmatrix} ,
\end{equation*}
sedangkan nilai sebenarnya adalah (sekali lagi dibulatkan hingga 8 tempat desimal)
\begin{equation*}
e^{A} =
\begin{bmatrix}
10.22670818           & \phantom{0}9.85882874 \\
\phantom{0}9.85882874 & 10.22670818
\end{bmatrix} .
\end{equation*}
Jadi, hampiran tersebut tidak terlalu baik ketika kita mencapai $t=1$.  Untuk memperoleh
hampiran yang baik pada $t=1$ (misalnya hingga 2 tempat desimal), kita perlu mencapai
pangkat ${11}^{\text{th}}$ (latihan).

\subsection{Latihan}

\begin{exercise}
Dengan menggunakan eksponensial matriks,
carilah solusi matriks fundamental untuk sistem
$x' = 3x+y$, $y' = x+3y$.
\end{exercise}

\begin{exercise}
Carilah $e^{tA}$ untuk matriks $A =
\left[ \begin{smallmatrix}
2 & 3 \\
0 & 2
\end{smallmatrix} \right]$.
\end{exercise}

\begin{exercise}
Carilah solusi matriks fundamental untuk sistem
$x_1' = 7x_1+4x_2+ 12x_3$,
$x_2' = x_1+2x_2+x_3$,
$x_3' = -3x_1-2x_2- 5x_3$.  Kemudian carilah solusi yang
memenuhi $\vec{x}(0) = 
\left[ \begin{smallmatrix} 0 \\ 1 \\ -2 \end{smallmatrix} \right]$.
\end{exercise}

\begin{exercise}
Hitunglah eksponensial matriks $e^A$ untuk
$A = \left[ \begin{smallmatrix} 1 & 2 \\ 0 & 1 \end{smallmatrix} \right]$.
\end{exercise}

\begin{exercise}[menantang] \label{matexp:explawex}
Misalkan $AB = BA$.  Tunjukkan bahwa dengan asumsi ini, $e^{A+B} = e^A e^B$.
\end{exercise}

\begin{exercise} \label{matexp:expinvex}
Gunakan \exerciseref{matexp:explawex}
untuk menunjukkan bahwa ${(e^{A})}^{-1} = e^{-A}$.  Secara khusus,
$e^A$ invertibel meskipun $A$ tidak.
\end{exercise}

\begin{exercise}
Misalkan $A$ matriks $2 \times 2$ dengan nilai eigen $-1$, $1$, dan
vektor eigen yang bersesuaian
$\left[ \begin{smallmatrix}
1 \\
1
\end{smallmatrix} \right]$,
$\left[ \begin{smallmatrix}
0 \\
1
\end{smallmatrix} \right]$.
\begin{tasks}
\task Carilah matriks $A$ dengan sifat-sifat ini.
\task Carilah solusi matriks fundamental dari ${\vec{x}}' = A \vec{x}$.
\task Selesaikan sistem dengan syarat awal $\vec{x}(0) =
\left[ \begin{smallmatrix}
2 \\
3
\end{smallmatrix} \right]$ .
\end{tasks}
\end{exercise}

\begin{exercise}
Misalkan $A$ matriks $n \times n$ dengan nilai eigen berulang
$\lambda$ bermultiplisitas $n$ dan memiliki
$n$ vektor eigen yang bebas linear.  Tunjukkan bahwa
matriks tersebut diagonal; bahkan, $A = \lambda I$.  Petunjuk: Gunakan
diagonalisasi dan fakta bahwa matriks identitas komutatif dengan setiap
matriks lain.
\end{exercise}

\begin{exercise}
Misalkan $A = \left[ \begin{smallmatrix}
-1 & -1 \\
1 & -3
\end{smallmatrix} \right]$.
\begin{tasks}(2)
\task Carilah $e^{tA}$.
\task Selesaikan
${\vec{x}}' = A \vec{x}$, $\vec{x}(0) =
\left[ \begin{smallmatrix}
1 \\
-2
\end{smallmatrix} \right]$.
\end{tasks}
\end{exercise}

\begin{exercise}
Misalkan $A = \left[ \begin{smallmatrix}
1 & 2 \\
3 & 4 
\end{smallmatrix} \right]$.
Hampiri $e^{tA}$ dengan mengembangkan deret pangkat
hingga orde ketiga.
\end{exercise}

\begin{exercise}
Untuk setiap bilangan bulat positif $n$,
carilah rumus (atau prosedur) untuk $A^n$ bagi matriks-matriks berikut:
\begin{tasks}(4)
\task
$
\begin{bmatrix}
3 & 0 \\ 0 & 9
\end{bmatrix}
$
\task
$
\begin{bmatrix}
5 & 2 \\ 4 & 7
\end{bmatrix}
$
\task
$
\begin{bmatrix}
0 & 1 \\ 0 & 0
\end{bmatrix}
$
\task
$
\begin{bmatrix}
2 & 1 \\ 0 & 2
\end{bmatrix}
$
\end{tasks}
\end{exercise}

\setcounter{exercise}{100}

\begin{exercise}
Hitunglah
$e^{tA}$ dengan
$A=\left[ \begin{smallmatrix}
1 & -2 \\
-2 & 1 
\end{smallmatrix}\right]$.
\end{exercise}
\exsol{%
$e^{tA}=\left[ \begin{smallmatrix}
\frac{{e}^{3t}+{e}^{-t}}{2} & \quad
\frac{{e}^{-t}-{e}^{3t}}{2}\\
\frac{{e}^{-t}-{e}^{3t}}{2} & \quad
\frac{{e}^{3t}+{e}^{-t}}{2}
\end{smallmatrix}\right]$
}

\begin{exercise}
Hitunglah
$e^{tA}$ dengan
$A=\left[ \begin{smallmatrix}
1 & -3 & 2 \\
-2 & 1 & 2 \\
-1 & -3 & 4
\end{smallmatrix}\right]$.
\end{exercise}
\exsol{%
$e^{tA}=\left[ \begin{smallmatrix}
2{e}^{3t}-4{e}^{2t}+3{e}^{t} & \quad
\frac{3{e}^{t}}{2}-\frac{3{e}^{3t}}{2} & \quad
-{e}^{3t}+4{e}^{2t}-3{e}^{t} \\
2{e}^{t}-2{e}^{2t} & \quad
{e}^{t} & \quad
2{e}^{2t}-2{e}^{t} \\
2{e}^{3t}-5{e}^{2t}+3{e}^{t} & \quad
\frac{3{e}^{t}}{2}-\frac{3{e}^{3t}}{2} & \quad
-{e}^{3t}+5{e}^{2t}-3{e}^{t}
\end{smallmatrix}\right]$
}


\begin{exercise}
\pagebreak[2]
\leavevmode
\begin{tasks}(2)
\task
Hitunglah
$e^{tA}$ dengan
$A=\left[ \begin{smallmatrix}
3 & -1 \\
1 & 1 
\end{smallmatrix}\right]$.
\task
Selesaikan $\vec{x}\,' = A \vec{x}$
untuk $\vec{x}(0) =
\left[ \begin{smallmatrix}
1 \\ 2
\end{smallmatrix}\right]$.
\end{tasks}
\end{exercise}
\exsol{%
a) $e^{tA}=\left[ \begin{smallmatrix}
( t+1) \,{e}^{2t} & -t{e}^{2t} \\
t{e}^{2t} & ( 1-t) \,{e}^{2t}
\end{smallmatrix}\right]$
\quad
b) $\vec{x}=\left[ \begin{smallmatrix}
(1-t) \,{e}^{2t} \\
(2-t) \,{e}^{2t}
\end{smallmatrix}\right]$
}

\begin{exercise}
Hitunglah 3 suku pertama (hingga derajat kedua) dari ekspansi Taylor
$e^{tA}$ dengan
$A=\left[ \begin{smallmatrix}
2 & 3 \\
2 & 2 
\end{smallmatrix}\right]$.  Tuliskan sebagai satu matriks.
Kemudian gunakan untuk menghampiri $e^{0.1A}$.
\end{exercise}
\exsol{%
$\left[ \begin{smallmatrix}
1+2t+5t^2 & 3t+6t^2 \\
2t+4t^2 & 1+2t+5t^2
\end{smallmatrix}\right]$
\quad
$e^{0.1A} \approx
\left[ \begin{smallmatrix}
1.25 & 0.36 \\
0.24 & 1.25
\end{smallmatrix}\right]$
}

\begin{exercise}
Untuk setiap bilangan bulat positif $n$,
carilah rumus (atau prosedur) untuk $A^n$ bagi matriks-matriks berikut:
\begin{tasks}(3)
\task
$\begin{bmatrix}
7 & 4 \\ -5 & -2
\end{bmatrix}$
\task
$\begin{bmatrix}
-3 & 4 \\ -6 & 7
\end{bmatrix}$
\task
$\begin{bmatrix}
0 & 1 \\ 1 & 0
\end{bmatrix}$
\end{tasks}
\end{exercise}
\exsol{%
a)
$\begin{bmatrix}
5 (3^n)-2^{n+2} & 4 (3^n)-2^{n+2}\\
5 (2^n)-5 (3^n) & 5 (2^n)-4 (3^n)
\end{bmatrix}$
\qquad
b)
$\begin{bmatrix}3-2 (3^n) & 2 (3^n)-2\\
3-3^{n+1} & 3^{n+1}-2\end{bmatrix}$
\\
c) 
$\begin{bmatrix}1 & 0 \\ 0 & 1\end{bmatrix}$
jika $n$ genap, dan
$\begin{bmatrix}0 & 1 \\ 1 & 0\end{bmatrix}$
jika $n$ ganjil.
}

%%%%%%%%%%%%%%%%%%%%%%%%%%%%%%%%%%%%%%%%%%%%%%%%%%%%%%%%%%%%%%%%%%%%%%%%%%%%%%

\sectionnewpage
\section{Sistem takhomogen}
\label{nonhomogsys:section}

%mbxINTROSUBSECTION

\sectionnotes{3 kuliah (mungkin perlu sedikit bagian dilewati)\EPref{, agak
berbeda dari \S5.7 dalam \cite{EP}}\BDref{,
\S7.9 dalam \cite{BD}}}

\subsection{Orde pertama berkoefisien konstan}

\subsubsection{Faktor integrasi}

Pertama-tama kita berfokus pada persamaan takhomogen orde pertama
\begin{equation*}
{\vec{x}}'(t) = A\vec{x}(t) + \vec{f}(t) ,
\end{equation*}
dengan $A$ suatu matriks konstan.  Metode pertama yang kita tinjau adalah
\emph{metode faktor integrasi\index{metode faktor integrasi!untuk sistem}}.
Untuk menyederhanakan, kita tuliskan kembali persamaan tersebut sebagai
\begin{equation*}
{\vec{x}}'(t) + P \vec{x}(t) = \vec{f}(t) ,
\end{equation*}
dengan $P = -A$.
Kita kalikan kedua ruas
persamaan dengan $e^{tP}$ (dengan mengingat bahwa kita menangani matriks-matriks yang
mungkin tidak komutatif) untuk memperoleh
\begin{equation*}
e^{tP}{\vec{x}}'(t) + e^{tP}P\vec{x}(t) = e^{tP}\vec{f}(t) .
\end{equation*}
Kita perhatikan bahwa $P e^{tP} = e^{tP} P$.  Fakta ini diperoleh dengan menuliskan
definisi deret dari $e^{tP}$:
\begin{equation*}
\begin{split}
P e^{tP} & = 
P \left(
I + tP + \frac{1}{2} {(tP)}^2 + \cdots \right)
=
P + tP^2 + \frac{1}{2} t^2P^3 + \cdots
=
\\
& =
\left(
I + tP + \frac{1}{2} {(tP)}^2 + \cdots \right) P 
= e^{tP} P . 
\end{split}
\end{equation*}
Jadi $\frac{d}{dt} \left( e^{tP} \right) = P e^{tP} = e^{tP} P$.  Aturan
hasil kali menyatakan
\begin{equation*}
\frac{d}{dt}
\Bigl( e^{tP} \vec{x}(t) \Bigr) =
e^{tP}{\vec{x}}'(t) + e^{tP}P\vec{x}(t),
\end{equation*}
sehingga
\begin{equation*}
\frac{d}{dt}
\Bigl( e^{tP} \vec{x}(t) \Bigr) = e^{tP}\vec{f}(t) .
\end{equation*}
Sekarang kita dapat mengintegralkan.  Artinya, kita integralkan setiap komponen vektor
secara terpisah
\begin{equation*}
e^{tP} \vec{x}(t) = \int e^{tP}\vec{f}(t) \, dt + \vec{c} .
\end{equation*}
Ingat dari \exerciseref{matexp:expinvex} bahwa ${(e^{tP})}^{-1} = e^{-tP}$.
Oleh karena itu, kita memperoleh
\begin{equation*}
\vec{x}(t) = e^{-tP} \int e^{tP}\vec{f}(t) \, dt + e^{-tP} \vec{c} .
\end{equation*}

Mungkin bentuk integral tentu lebih mudah dipahami.  Dalam kasus ini,
syarat awal juga mudah diselesaikan.
Tinjau
persamaan dengan syarat awal
\begin{equation*}
{\vec{x}}'(t) + P\vec{x}(t) = \vec{f}(t) ,
\qquad \vec{x}(0) = \vec{b} .
\end{equation*}
Solusinya kemudian dapat dituliskan sebagai
\begin{equation} \label{nhsys:intfacsoleq}
\mybxbg{~~
\vec{x}(t) = e^{-tP} \int_0^t e^{sP}\vec{f}(s) \, ds + e^{-tP} \vec{b} .
~~}
\end{equation}
Sekali lagi, integrasi berarti bahwa setiap komponen vektor 
$e^{sP}\vec{f}(s)$ diintegralkan secara terpisah.
Tidak sulit dilihat bahwa \eqref{nhsys:intfacsoleq} benar-benar memenuhi
syarat awal $\vec{x}(0) = \vec{b}$.
\begin{equation*}
\vec{x}(0) = e^{-0P} \int_0^0 e^{sP}\vec{f}(s) \, ds + e^{-0P} \vec{b}
= I \vec{b} = \vec{b} .
\end{equation*}

\begin{example}
Misalkan kita memiliki sistem
\begin{align*}
x_1' + 5x_1 - 3x_2 &= e^t , \\
x_2' + 3x_1 - x_2 &= 0 ,
\end{align*}
dengan syarat awal $x_1(0) = 1, x_2(0) = 0$.

Mari kita tuliskan sistem tersebut sebagai
\begin{equation*}
{\vec{x}}' +
\begin{bmatrix} 5 & -3 \\ 3 & -1 \end{bmatrix}
\vec{x} =
\begin{bmatrix} e^t \\ 0 \end{bmatrix} ,
\qquad
\vec{x}(0) = 
\begin{bmatrix} 1 \\ 0 \end{bmatrix} .
\end{equation*}
Matriks
$P = \left[
\begin{smallmatrix} 5 & -3 \\ 3 & -1 \end{smallmatrix} \right]$
memiliki nilai eigen ganda 2 dengan defek 1, dan kita serahkan sebagai latihan
untuk memeriksa kembali bahwa kita menghitung $e^{tP}$ dengan benar.  Setelah memperoleh
$e^{tP}$, kita mencari $e^{-tP}$ hanya dengan menegatifkan $t$.
\begin{equation*}
e^{tP} = 
\begin{bmatrix}
(1+3t)\,e^{2t} & -3te^{2t} \\
3te^{2t} & (1-3t)\,e^{2t}
\end{bmatrix}
, \qquad
e^{-tP} = 
\begin{bmatrix}
(1-3t)\,e^{-2t} & 3te^{-2t} \\
-3te^{-2t} & (1+3t)\,e^{-2t}
\end{bmatrix}
.
\end{equation*}
Alih-alih menghitung seluruh rumus sekaligus, mari kita kerjakan secara bertahap.
Pertama,
\begin{equation*}
\begin{split}
\int_0^t e^{sP}\vec{f}(s) \, ds & = 
\int_0^t
 \begin{bmatrix}
 (1+3s)\,e^{2s} & -3se^{2s} \\
 3se^{2s} & (1-3s)\,e^{2s}
 \end{bmatrix}
 \begin{bmatrix} e^{s} \\ 0 \end{bmatrix}
\, ds
\\
& =
\int_0^t
 \begin{bmatrix}
 (1+3s)\,e^{3s} \\
 3se^{3s}
 \end{bmatrix}
\, ds
\\
 &=
 \begin{bmatrix}
 \int_0^t (1+3s)\,e^{3s} \,ds \\
 \int_0^t 3se^{3s} \,ds
 \end{bmatrix}
\\
& =
\begin{bmatrix}
t e^{3t} \\
\frac{(3t-1) \,e^{3t} + 1}{3}
\end{bmatrix} %.
\qquad \qquad \text{(menggunakan integrasi parsial).}
\end{split}
\end{equation*}
Kemudian
\begin{equation*}
\begin{split}
\vec{x}(t)
& = e^{-tP} \int_0^t e^{sP}\vec{f}(s) \, ds + e^{-tP} \vec{b} \\
& =
\begin{bmatrix}
(1-3t)\,e^{-2t} & 3te^{-2t} \\
-3te^{-2t} & (1+3t)\,e^{-2t}
\end{bmatrix}
\begin{bmatrix}
t e^{3t} \\
\frac{(3t-1) \,e^{3t} + 1}{3}
\end{bmatrix}
+
\begin{bmatrix}
(1-3t)\,e^{-2t} & 3te^{-2t} \\
-3te^{-2t} & (1+3t)\,e^{-2t}
\end{bmatrix}
\begin{bmatrix} 1 \\ 0 \end{bmatrix} \\
& =
\begin{bmatrix}
te^{-2t} \\
-\frac{e^t}{3}+\left( \frac{1}{3} + t \right) \, e^{-2t}
\end{bmatrix}
+
\begin{bmatrix}
(1-3t)\,e^{-2t} \\
-3te^{-2t}
\end{bmatrix} \\
& =
\begin{bmatrix}
(1-2t)\,e^{-2t} \\
-\frac{e^t}{3}+\left( \frac{1}{3} -2 t \right) \, e^{-2t}
\end{bmatrix} .
\end{split}
\end{equation*}
Akhirnya!

Mari kita periksa bahwa hasil ini benar-benar berlaku.
\begin{equation*}
x_1' + 5 x_1 - 3x_2 = (4te^{-2t} - 4 e^{-2t}) + 5
(1-2t)\,e^{-2t} 
+e^t-( 1 -6 t ) \, e^{-2t} = e^t .
\end{equation*}
Demikian pula (latihan), $x_2' + 3 x_1 - x_2 = 0$.   Syarat awal
juga dipenuhi (latihan).
\end{example}

Untuk sistem, metode faktor integrasi
hanya berlaku jika $P$ tidak bergantung pada $t$, yakni
$P$ konstan.  Masalahnya
adalah bahwa secara umum
\begin{equation*}
\frac{d}{dt} \left[ e^{\int P(t)\,dt} \right] \neq P(t) \, e^{\int P(t)\,dt} ,
\end{equation*}
karena perkalian matriks tidak komutatif.

\subsubsection{Dekomposisi vektor eigen}
\index{eigenvector decomposition}

Untuk metode berikutnya, perhatikan bahwa vektor eigen suatu matriks memberikan 
arah-arah tempat matriks tersebut bertindak seperti skalar.  Jika kita menyelesaikan
sistem sepanjang arah-arah ini, perhitungannya lebih sederhana karena kita
memperlakukan matriks sebagai skalar.  Kemudian kita
menggabungkan solusi-solusi tersebut untuk memperoleh solusi umum
sistem.

Ambil persamaan
\begin{equation} \label{nhsys:ednhsys}
{\vec{x}}' (t) = A \vec{x}(t) + \vec{f}(t) .
\end{equation}
Andaikan $A$ memiliki $n$ vektor eigen yang bebas linear,
$\vec{v}_1, \vec{v}_2, \ldots, \vec{v}_n$.  Tuliskan
\begin{equation} \label{nhsys:decompx}
\vec{x}(t) =
\vec{v}_1 \, \xi_1(t) + 
\vec{v}_2 \, \xi_2(t) + \cdots +
\vec{v}_n \, \xi_n(t) .
\end{equation}
Artinya, kita ingin menuliskan solusi sebagai kombinasi linear dari
vektor-vektor eigen $A$.
Jika kita menyelesaikan fungsi skalar $\xi_1$ sampai $\xi_n$, kita memperoleh
solusi $\vec{x}$.
Mari kita uraikan pula $\vec{f}$ dalam vektor-vektor eigen.  Kita ingin menuliskan
\begin{equation} \label{nhsys:decompf}
\vec{f}(t) =
\vec{v}_1 \, g_1(t) + 
\vec{v}_2 \, g_2(t) + \cdots +
\vec{v}_n \, g_n(t) .
\end{equation}
Artinya, kita ingin mencari $g_1$ sampai $g_n$ yang memenuhi
\eqref{nhsys:decompf}.  Karena semua vektor eigen
bebas linear, matriks
$E = [\, \vec{v}_1 \quad \vec{v}_2 \quad \cdots \quad \vec{v}_n \,]$
invertibel.  Tuliskan persamaan \eqref{nhsys:decompf}
sebagai $\vec{f} = E \vec{g}$, dengan komponen-komponen
$\vec{g}$ adalah fungsi $g_1$ sampai $g_n$.
Maka $\vec{g} = E^{-1} \vec{f}$.
Jadi, $\vec{g}$ selalu dapat dicari apabila terdapat
$n$ vektor eigen yang bebas linear.

Kita substitusikan \eqref{nhsys:decompx} ke \eqref{nhsys:ednhsys}, dan perhatikan bahwa
$A \vec{v}_k = \lambda_k \vec{v}_k$:
\begin{equation*}
\begin{split}
%{\vec{x}}' & =
\overbrace{
\vec{v}_1 \xi_1' + 
\vec{v}_2 \xi_2' + \cdots +
\vec{v}_n \xi_n'
}^{{\vec{x}}'}
%\\
& =
\overbrace{
A \left( \vec{v}_1 \xi_1 + 
\vec{v}_2 \xi_2 + \cdots +
\vec{v}_n \xi_n \right)
}^{A\vec{x}}
+
\overbrace{
\vec{v}_1 g_1 + 
\vec{v}_2 g_2 + \cdots +
\vec{v}_n g_n
}^{\vec{f}}
\\
& = 
A \vec{v}_1 \xi_1 + 
A \vec{v}_2 \xi_2 + \cdots +
A \vec{v}_n \xi_n
+
\vec{v}_1 g_1 + 
\vec{v}_2 g_2 + \cdots +
\vec{v}_n g_n
\\
& =
\vec{v}_1 \lambda_1 \xi_1 + 
\vec{v}_2 \lambda_2 \xi_2 + \cdots +
\vec{v}_n \lambda_n \xi_n
+
\vec{v}_1 g_1 + 
\vec{v}_2 g_2 + \cdots +
\vec{v}_n g_n
\\
& =
\vec{v}_1 ( \lambda_1 \xi_1 + g_1 ) +
\vec{v}_2 ( \lambda_2 \xi_2 + g_2 ) + \cdots + 
\vec{v}_n ( \lambda_n \xi_n + g_n ) .
\end{split}
\end{equation*}
Dengan menyamakan koefisien vektor $\vec{v}_1$ sampai
$\vec{v}_n$, kita memperoleh
persamaan
\begin{align*}
\xi_1' & = \lambda_1 \xi_1 + g_1 , \\
\xi_2' & = \lambda_2 \xi_2 + g_2 , \\
& ~~ \vdots \\
\xi_n' & = \lambda_n \xi_n + g_n .
\end{align*}
Setiap persamaan ini tidak bergantung pada persamaan lainnya.  Semuanya merupakan
persamaan linear orde pertama dan dapat diselesaikan dengan mudah menggunakan
metode faktor integrasi standar untuk persamaan tunggal.
Artinya, untuk persamaan $k^{\text{th}}$,
kita tuliskan
\begin{equation*}
\xi_k'(t) - \lambda_k \xi_k(t) = g_k(t) .
\end{equation*}
Kita gunakan faktor integrasi $e^{-\lambda_k t}$ untuk memperoleh
\begin{equation*}
\frac{d}{dt}\Bigl[ \xi_k(t) \, e^{-\lambda_k t} \Bigr] = 
e^{-\lambda_k t} g_k(t) .
\end{equation*}
Kita integralkan dan selesaikan untuk $\xi_k$ sehingga memperoleh
\begin{equation*}
\xi_k(t) =  e^{\lambda_k t} 
\int e^{-\lambda_k t} g_k(t) \,dt + C_k e^{\lambda_k t} .
\end{equation*}
Jika kita hanya mencari sebarang solusi khusus, kita dapat menetapkan
$C_k$ sama dengan nol.  Jika konstanta-konstanta ini dipertahankan, kita memperoleh
solusi umum.  Tuliskan
$\vec{x}(t) =
\vec{v}_1 \xi_1(t) + 
\vec{v}_2 \xi_2(t) + \cdots +
\vec{v}_n \xi_n(t)$, dan selesailah kita.

\medskip

Seperti biasa,
mungkin lebih baik menuliskan integral-integral ini sebagai integral tentu.  
Misalkan kita memiliki syarat awal $\vec{x}(0) = \vec{b}$.
Ambil $\vec{a} = E^{-1} \vec{b}$ untuk memperoleh
$\vec{b} = \vec{v}_1 a_1 + \vec{v}_2 a_2 + \cdots + \vec{v}_n a_n$,
seperti sebelumnya.
Kemudian, jika kita tuliskan 
\begin{equation*}
\mybxbg{~~
\xi_k(t) =  e^{\lambda_k t} 
\int_0^t e^{-\lambda_k s} g_k(s) \,ds + a_k e^{\lambda_k t} ,
~~}
\end{equation*}
kita memperoleh solusi khusus
$\vec{x}(t) =
\vec{v}_1 \xi_1(t) + 
\vec{v}_2 \xi_2(t) + \cdots +
\vec{v}_n \xi_n(t)$ yang memenuhi $\vec{x}(0) = \vec{b}$,
karena $\xi_k(0) = a_k$.

Kita perhatikan bahwa teknik yang baru saja diuraikan adalah metode
nilai eigen yang diterapkan pada sistem takhomogen.  Jika suatu sistem homogen,
yakni jika $\vec{f}=\vec{0}$, maka persamaan yang kita peroleh adalah
$\xi_k' = \lambda_k \xi_k$, sehingga $\xi_k = C_k e^{\lambda_k t}$ merupakan
solusinya, dan itulah tepatnya yang kita peroleh dalam
\sectionref{eigenmethod:section}.

\begin{example}
Misalkan $A = \left[
\begin{smallmatrix}
1 & 3 \\
3 & 1
\end{smallmatrix} \right]$.
Selesaikan ${\vec{x}}' = A \vec{x} +
\vec{f}$ dengan $\vec{f}(t) = 
\left[ \begin{smallmatrix}
2e^t \\
2t
\end{smallmatrix} \right]$ untuk $\vec{x}(0) =
\left[ \begin{smallmatrix}
3/16 \\
-5/16
\end{smallmatrix} \right]$.

Nilai eigen $A$ adalah $-2$ dan 4, dan vektor eigen yang bersesuaian
adalah
$\left[ \begin{smallmatrix}
1 \\
-1
\end{smallmatrix} \right]$ dan
$\left[ \begin{smallmatrix}
1 \\
1
\end{smallmatrix} \right]$ berturut-turut.  Perhitungan ini diserahkan sebagai
latihan.  Kita tuliskan matriks $E$ dari vektor-vektor eigen dan hitung
inversnya (dengan menggunakan rumus invers untuk matriks $2 \times 2$)
\begin{equation*}
E = \begin{bmatrix}
1 & 1 \\
-1 & 1
\end{bmatrix} ,
\qquad
E^{-1}
=
\frac{1}{2}
\begin{bmatrix}
1 & -1 \\
1 & 1
\end{bmatrix} .
\end{equation*}

Kita mencari solusi berbentuk $\vec{x} = 
\left[ \begin{smallmatrix}
1 \\
-1
\end{smallmatrix} \right] \xi_1 +
\left[ \begin{smallmatrix}
1 \\
1
\end{smallmatrix} \right] \xi_2$.  Pertama, kita perlu menuliskan $\vec{f}$
dalam vektor-vektor eigen.
Artinya, kita ingin menuliskan $\vec{f} = 
\left[ \begin{smallmatrix}
2e^t \\
2t
\end{smallmatrix} \right] = 
\left[ \begin{smallmatrix}
1 \\
-1
\end{smallmatrix} \right] g_1 +
\left[ \begin{smallmatrix}
1 \\
1
\end{smallmatrix} \right] g_2$.  Dengan demikian,
\begin{equation*}
\begin{bmatrix}
g_1 \\
g_2
\end{bmatrix} = 
E^{-1}
\begin{bmatrix}
2e^t \\
2t
\end{bmatrix}
=
\frac{1}{2}
\begin{bmatrix}
1 & -1 \\
1 & 1
\end{bmatrix}
\begin{bmatrix}
2e^t \\
2t
\end{bmatrix}
=
\begin{bmatrix}
e^t-t \\
e^t+t
\end{bmatrix} .
\end{equation*}
Jadi $g_1 = e^t-t$ dan $g_2 = e^t+t$.

Selanjutnya kita perlu menuliskan $\vec{x}(0)$ dalam vektor-vektor eigen.
Artinya, kita ingin menuliskan $\vec{x}(0) = 
\left[ \begin{smallmatrix}
3/16 \\
-5/16
\end{smallmatrix} \right] = 
\left[ \begin{smallmatrix}
1 \\
-1
\end{smallmatrix} \right] a_1 +
\left[ \begin{smallmatrix}
1 \\
1
\end{smallmatrix} \right] a_2$.  Oleh karena itu,
\begin{equation*}
\begin{bmatrix}
a_1 \\
\noalign{\smallskip}
a_2
\end{bmatrix} = 
E^{-1}
\begin{bmatrix}
\nicefrac{3}{16} \\
\noalign{\smallskip}
\nicefrac{-5}{16}
\end{bmatrix}
=
\begin{bmatrix}
\nicefrac{1}{4} \\
\noalign{\smallskip}
\nicefrac{-1}{16}
\end{bmatrix} .
\end{equation*}
Jadi $a_1 = \nicefrac{1}{4}$ dan $a_2 = \nicefrac{-1}{16}$.
Kita substitusikan $\vec{x}$ ke persamaan dan memperoleh
\begin{equation*}
\begin{split}
\overbrace{
\begin{bmatrix}
1 \\
-1
\end{bmatrix} \xi_1' +
\begin{bmatrix}
1 \\
1
\end{bmatrix} \xi_2'
}^{\vec{x}'}
& =
\overbrace{
A
\begin{bmatrix}
1 \\
-1
\end{bmatrix} \xi_1 +
A
\begin{bmatrix}
1 \\
1
\end{bmatrix} \xi_2
}^{A\vec{x}}
+
\overbrace{
\begin{bmatrix}
1 \\
-1
\end{bmatrix} g_1 +
\begin{bmatrix}
1 \\
1
\end{bmatrix} g_2
}^{\vec{f}}
\\
& =
\begin{bmatrix}
1 \\
-1
\end{bmatrix} (-2\xi_1) +
\begin{bmatrix}
1 \\
1
\end{bmatrix} 4\xi_2
+
\begin{bmatrix}
1 \\
-1
\end{bmatrix} (e^t - t)
+
\begin{bmatrix}
1 \\
1
\end{bmatrix} (e^t + t) .
\end{split}
\end{equation*}
Kita memperoleh dua persamaan
\begin{align*}
& \xi_1' = -2\xi_1 + e^t -t, & & \text{dengan } \xi_1(0) = a_1 = \frac{1}{4} , \\
& \xi_2' = 4\xi_2 + e^t + t, & & \text{dengan } \xi_2(0) = a_2 = \frac{-1}{16} .
\end{align*}
Kita selesaikan dengan metode faktor integrasi.  Perhitungan integral diserahkan
sebagai latihan bagi mahasiswa.  Anda akan memerlukan integrasi parsial.
\begin{equation*}
\xi_1 = e^{-2t}\int e^{2t} \, (e^t-t) \, dt + C_1 e^{-2t} = 
\frac{e^t}{3}-\frac{t}{2}+\frac{1}{4}+C_1 e^{-2t} ,
\end{equation*}
dengan $C_1$ konstanta integrasi.
Karena $\xi_1(0) = \nicefrac{1}{4}$, maka $\nicefrac{1}{4}= \nicefrac{1}{3}
+ \nicefrac{1}{4} + C_1$ dan
$C_1 = \nicefrac{-1}{3}$.
Demikian pula,
\begin{equation*}
\xi_2 = e^{4t}\int e^{-4t} \, (e^t+ t) \, dt + C_2 e^{4t} = 
-\frac{e^t}{3}-\frac{t}{4}-\frac{1}{16} + C_2 e^{4t} .
\end{equation*}
Karena $\xi_2(0) = \nicefrac{-1}{16}$, kita memiliki $\nicefrac{-1}{16}= \nicefrac{-1}{3}
-\nicefrac{1}{16} + C_2$, sehingga
$C_2 = \nicefrac{1}{3}$.
Solusinya adalah
\begin{equation*}
\vec{x}(t)=
\begin{bmatrix}
1 \\
-1
\end{bmatrix}
\underbrace{\left( \frac{e^t-e^{-2t}}{3}+\frac{1-2t}{4} \right)}%
_{\xi_1} +
\begin{bmatrix}
1 \\
1
\end{bmatrix}
\underbrace{\left( \frac{e^{4t}-e^t}{3}-\frac{4t+1}{16} \right)}%
_{\xi_2}
=
\begin{bmatrix}
\frac{e^{4t}-e^{-2t}}{3}+\frac{3-12t}{16} \\
\frac{e^{-2t}+e^{4t}-2e^t}{3}+\frac{4t-5}{16}
\end{bmatrix} .
\end{equation*}
Artinya,
$x_1 = \frac{e^{4t}-e^{-2t}}{3}+\frac{3-12t}{16}$
dan
$x_2 = \frac{e^{-2t}+e^{4t}-2e^t}{3}+\frac{4t-5}{16}$.
\end{example}

\begin{exercise}
Periksalah bahwa $x_1$ dan $x_2$ menyelesaikan masalah.  Periksalah bahwa keduanya
memenuhi persamaan diferensial dan syarat awal.
\end{exercise}

\subsubsection{Koefisien tak tentu}

Metode
koefisien tak tentu\index{koefisien tak tentu!untuk sistem} juga
berlaku untuk sistem.
Satu-satunya
perbedaan adalah bahwa kita menggunakan vektor yang tidak diketahui, bukan sekadar
bilangan.  Peringatan yang sama bagi koefisien tak tentu pada sistem
berlaku seperti pada persamaan tunggal.  Metode ini tidak selalu berhasil.
Selain itu, jika ruas kanan rumit, kita harus menyelesaikan
banyak
variabel.  Setiap elemen dari
vektor yang tidak diketahui merupakan bilangan yang tidak diketahui.  Dalam sistem 3 persamaan dengan,
misalnya, 4 vektor yang tidak diketahui (hal ini tidak jarang), kita sudah memiliki 12
bilangan yang harus dicari.
Metode ini dapat menjadi pekerjaan yang sangat
\myindex{melelahkan} % a bit of fun
jika dilakukan dengan tangan.
Karena metode tersebut pada dasarnya sama seperti untuk persamaan tunggal,
mari kita langsung kerjakan sebuah contoh.

\begin{example}
Misalkan $A = \left[
\begin{smallmatrix}
-1 & 0 \\
-2 & 1
\end{smallmatrix} \right]$.
Carilah solusi khusus dari ${\vec{x}}' = A \vec{x} +
\vec{f}$ dengan $\vec{f}(t) = 
\left[ \begin{smallmatrix}
e^t \\
t
\end{smallmatrix} \right]$.

Sistem ini dapat diselesaikan dengan cara yang lebih mudah (dapatkah Anda melihat caranya?), tetapi
untuk keperluan contoh, kita menggunakan metode nilai eigen dan
koefisien tak tentu.
Nilai eigen $A$ adalah 
$-1$ dan $1$.
Vektor eigen yang bersesuaian adalah
$\left[ \begin{smallmatrix}
1 \\
1
\end{smallmatrix} \right]$ dan
$\left[ \begin{smallmatrix}
0 \\
1
\end{smallmatrix} \right]$ berturut-turut.
Jadi, solusi komplementer kita adalah
\begin{equation*}
\vec{x}_c = 
\alpha_1
\begin{bmatrix}
1 \\ 1
\end{bmatrix}
e^{-t}
+
\alpha_2
\begin{bmatrix}
0 \\ 1
\end{bmatrix}
e^{t} ,
\end{equation*}
untuk konstanta sebarang $\alpha_1$ dan $\alpha_2$.

Selanjutnya kita ingin menebak solusi khusus berbentuk
\begin{equation*}
\vec{x} = 
\vec{a}
e^{t}
+
\vec{b}
t +
\vec{c} .
\end{equation*}
Namun, sesuatu yang berbentuk $\vec{a} e^t$ muncul dalam solusi
komplementer.  Karena kita belum mengetahui apakah vektor $\vec{a}$
merupakan kelipatan dari $\left[ \begin{smallmatrix}
0 \\
1
\end{smallmatrix} \right]$, kita tidak mengetahui apakah terjadi
konflik.  Mungkin saja tidak ada konflik,
tetapi agar aman kita juga mencoba
$\vec{b} t e^t$.
Di sini kita menemukan
inti perbedaan antara persamaan tunggal dan sistem.  Kita mencoba
\emph{kedua} suku $\vec{a} e^t$ dan $\vec{b} t e^t$ dalam solusi,
bukan hanya suku $\vec{b} t e^t$.
Oleh karena itu, kita mencoba
\begin{equation*}
\vec{x} = 
\vec{a}
e^{t}
+
\vec{b}
t
e^{t}
+
\vec{c}
t +
\vec{d}.
\end{equation*}
Kita memiliki 8 besaran yang tidak diketahui karena
$\vec{a} =
\Bigl[ \begin{smallmatrix} a_1 \\ a_2 \end{smallmatrix} \Bigr]$,
$\vec{b} =
\Bigl[ \begin{smallmatrix} b_1 \\ b_2 \end{smallmatrix} \Bigr]$,
$\vec{c} =
\Bigl[ \begin{smallmatrix} c_1 \\ c_2 \end{smallmatrix} \Bigr]$,
dan
$\vec{d} =
\Bigl[ \begin{smallmatrix} d_1 \\ d_2 \end{smallmatrix} \Bigr]$.
Kita substitusikan $\vec{x}$ ke persamaan.  Pertama,
mari kita hitung ${\vec{x}}'$.
\begin{equation*}
{\vec{x}}' = 
\left( \vec{a} + \vec{b} \right)
e^{t}
+
\vec{b}
t
e^{t}
+
\vec{c} =
\begin{bmatrix}
a_1 + b_1 \\ a_2+b_2
\end{bmatrix}
e^{t}
+
\begin{bmatrix}
b_1 \\ b_2
\end{bmatrix}
t e^{t}
+
\begin{bmatrix}
c_1 \\ c_2
\end{bmatrix} .
\end{equation*}
Sekarang ${\vec{x}}'$ harus sama dengan $A\vec{x} + \vec{f}$, yaitu
\begin{equation*}
\begin{split}
A \vec{x} + \vec{f} &=
A \vec{a}
e^{t}
+
A \vec{b}
t e^{t}
+
A \vec{c}
t
+
A \vec{d}
+ \vec{f}
\\
& =
\begin{bmatrix}
-a_1 \\ -2a_1+a_2
\end{bmatrix}
e^{t}
+
\begin{bmatrix}
-b_1 \\ -2b_1+b_2
\end{bmatrix}
t e^{t}
+
\begin{bmatrix}
-c_1 \\ -2c_1+c_2
\end{bmatrix}
t
+
\begin{bmatrix}
-d_1 \\ -2d_1+d_2
\end{bmatrix}
+
\begin{bmatrix}
1 \\ 0
\end{bmatrix} 
e^t
+
\begin{bmatrix}
0 \\ 1
\end{bmatrix} 
t 
\\
&=
\begin{bmatrix}
-a_1+1 \\ -2a_1+a_2
\end{bmatrix}
e^{t}
+
\begin{bmatrix}
-b_1 \\ -2b_1+b_2
\end{bmatrix}
t e^{t}
+
\begin{bmatrix}
-c_1 \\ -2c_1+c_2+1
\end{bmatrix}
t
+
\begin{bmatrix}
-d_1 \\ -2d_1+d_2
\end{bmatrix} .
\end{split}
\end{equation*}
Kita samakan koefisien $e^t$, $te^t$, $t$, dan setiap vektor
konstan dalam $\vec{x}'$ dan $A\vec{x}+\vec{f}$ untuk memperoleh persamaan:
\begin{align*}
a_1+b_1 & = -a_1+1 , &
0 & = -c_1 , \\
a_2+b_2 & = -2a_1+a_2 , &
0 & = -2c_1+c_2 + 1 , \\
b_1 & = -b_1 , &
c_1 & = -d_1 , \\
b_2 & = -2b_1+b_2 , &
c_2 & = -2d_1+d_2 .
\end{align*}
Kita dapat menuliskan matriks diperbesar $8 \times 9$ dan melakukan
reduksi baris, tetapi lebih mudah menyelesaikannya secara ad hoc.
Kita segera melihat $b_1 = 0$, $c_1 = 0$, $d_1 = 0$.  Dengan menyubstitusikan
nilai-nilai ini
kembali, kita memperoleh $c_2 = -1$ dan $d_2 = -1$.  Persamaan sisanya
yang memberi informasi adalah
\begin{align*}
a_1 & = -a_1+1 , \\
a_2+b_2 & = -2a_1+a_2 .
\end{align*}
Jadi $a_1 = \nicefrac{1}{2}$ dan $b_2 = -1$.  Terakhir, $a_2$ dapat sebarang dan tetap memenuhi
persamaan.  Kita hanya mencari satu solusi, sehingga mungkin yang
paling sederhana
adalah ketika $a_2 = 0$.
Oleh karena itu,
\begin{equation*}
\vec{x} = 
\vec{a}
e^{t}
+
\vec{b}
t
e^{t}
+
\vec{c}
t +
\vec{d}
=
\begin{bmatrix}
\nicefrac{1}{2} \\ 0
\end{bmatrix}
e^t
+
\begin{bmatrix}
0 \\ -1
\end{bmatrix}
te^t
+
\begin{bmatrix}
0 \\ -1
\end{bmatrix}
t
+
\begin{bmatrix}
0 \\ -1
\end{bmatrix}
=
\begin{bmatrix}
\frac{1}{2}\,e^t \\
-te^t - t - 1
\end{bmatrix} .
\end{equation*}
Artinya, $x_1 = \frac{1}{2}\,e^t$,
$x_2 = -te^t - t - 1$.
Kita tambahkan solusi khusus ini ke solusi komplementer untuk memperoleh
solusi umum masalah:
\begin{equation*}
x_1 = \frac{1}{2}\,e^t + 
\alpha_1 e^{-t} \quad \text{dan} \quad
x_2 = -te^t - t - 1 + \alpha_1 e^{-t} + \alpha_2 e^{t}.
\avoidbreak
\end{equation*}
Perhatikan bahwa baik $\vec{a} e^t$
maupun $\vec{b} te^t$ memang diperlukan.
\end{example}

\begin{exercise}
Periksalah bahwa $x_1$ dan $x_2$ menyelesaikan masalah.  Kemudian tetapkan $a_2 = 1$
dan periksa pula solusi ini.  Apa perbedaan antara kedua
solusi (satu dengan $a_2=0$ dan satu dengan $a_2=1$)?
\end{exercise}

Seperti dapat Anda lihat, selain penanganan konflik, metode koefisien
tak tentu bekerja persis sama seperti pada
persamaan tunggal.  Namun, perhitungannya dapat dengan cepat menjadi
sangat rumit untuk sistem.  Persamaan yang kita tinjau cukup sederhana.

\subsection{Orde pertama berkoefisien variabel}

\subsubsection{Variasi parameter}

Seperti untuk persamaan tunggal, terdapat metode
variasi parameter\index{variasi parameter!untuk sistem}.
Untuk sistem berkoefisien konstan, metode ini pada dasarnya
sama dengan metode faktor integrasi yang telah dibahas.
Namun, metode ini berlaku untuk setiap sistem linear, sekalipun tidak
berkoefisien konstan, asalkan kita entah bagaimana menyelesaikan masalah
homogen yang terkait.

Tinjau persamaan
\begin{equation} \label{nhsys:nhvceq}
{\vec{x}}' = A(t) \, \vec{x} + \vec{f}(t) .
\end{equation}
Misalkan kita entah bagaimana telah menyelesaikan persamaan homogen terkait,
${\vec{x}}' = A(t) \, \vec{x}$, dan menemukan
solusi matriks fundamental
$X(t)$.  Solusi umum persamaan homogen terkait
adalah $X(t) \vec{c}$ untuk vektor konstan $\vec{c}$.  Seperti dalam
variasi parameter untuk persamaan tunggal, kita mencoba solusi
persamaan takhomogen berbentuk
\begin{equation*}
\vec{x}_p = X(t)\, \vec{u}(t) ,
\end{equation*}
dengan $\vec{u}(t)$ fungsi bernilai vektor, bukan konstanta.
Kita substitusikan $\vec{x}_p$ ke \eqref{nhsys:nhvceq}:
\begin{equation*}
%{\vec{x}_p}'(t)
%=
\underbrace{X'(t)\, \vec{u}(t) + X(t)\, {\vec{u}}'(t)}%
_{{\vec{x}_p}'(t)}
=
\underbrace{A(t)\, X(t)\, \vec{u}(t)}%
_{A(t) \vec{x}_p (t)} 
 + \vec{f}(t) .
\end{equation*}
Karena $X(t)$ adalah solusi matriks fundamental dari masalah homogen,
yakni
$X'(t) = A(t)X(t)$, kita memperoleh
\begin{equation*}
\cancel{X'(t)\, \vec{u}(t)} + X(t)\, {\vec{u}}'(t)
=
\cancel{X'(t)\, \vec{u}(t)} + \vec{f}(t) .
\end{equation*}
Jadi,
$X(t)\, {\vec{u}}'(t) = \vec{f}(t)$.  Kita hitung
$\left[X(t)\right]^{-1}$,
lalu
${\vec{u}}'(t) = \left[X(t)\right]^{-1}\vec{f}(t)$.  Kita integralkan
untuk memperoleh $\vec{u}$, sehingga kita memiliki solusi khusus
$\vec{x}_p = X(t)\, \vec{u}(t)$.
Dengan demikian, kita memperoleh rumus
\begin{equation*}
\mybxbg{~~
\vec{x}_p = 
X(t)
\int \left[X(t)\right]^{-1}\vec{f}(t) \, dt .
~~}
\end{equation*}

Jika $A$ adalah matriks konstan dan $X(t) = e^{tA}$, maka
$\left[X(t)\right]^{-1} = e^{-tA}$.  
Kita memperoleh solusi
$\vec{x}_p = 
e^{tA}
\int e^{-tA}\,\vec{f}(t) \, dt$,
yang tepat sama dengan hasil metode faktor integrasi.

\begin{example}
Carilah solusi khusus dari
\begin{equation} \label{nhsys:vcexeq}
{\vec{x}}'
=
\frac{1}{t^2+1}
\begin{bmatrix}
t & -1 \\
1 & t
\end{bmatrix}
\vec{x}
+ \begin{bmatrix} t \\ 1 \end{bmatrix} \,(t^2+1) .
\end{equation}

Di sini, $A = 
\frac{1}{t^2+1}
\left[ \begin{smallmatrix}
t & -1 \\
1 & t
\end{smallmatrix} \right]$ jelas tidak konstan.
Mungkin melalui tebakan yang beruntung, kita menemukan bahwa
$X = 
\left[ \begin{smallmatrix}
1 & -t \\
t & 1
\end{smallmatrix} \right]$ menyelesaikan
$X'(t) = A(t) X(t)$.  Setelah mengetahui solusi komplementer, kita dapat mencari
solusi untuk \eqref{nhsys:vcexeq}.  Pertama, kita hitung
\begin{equation*}
\left[ X(t) \right]^{-1}
=
\frac{1}{t^2+1}
\begin{bmatrix}
1 & t \\
-t & 1
\end{bmatrix} .
\end{equation*}
Selanjutnya, kita mengetahui bahwa suatu solusi khusus dari
\eqref{nhsys:vcexeq} adalah
\begin{equation*}
\begin{split}
\vec{x}_p & = 
X(t)
\int \left[X(t)\right]^{-1}\vec{f}(t) \, dt 
\\
& =
\begin{bmatrix}
1 & -t \\
t & 1
\end{bmatrix} 
\int
\frac{1}{t^2+1}
\begin{bmatrix}
1 & t \\
-t & 1
\end{bmatrix}
\begin{bmatrix} t \\ 1 \end{bmatrix} \,(t^2+1) 
\,dt
\\
& =
\begin{bmatrix}
1 & -t \\
t & 1
\end{bmatrix} 
\int
\begin{bmatrix}
2t \\
-t^2 + 1
\end{bmatrix} 
\,dt
\\
& =
\begin{bmatrix}
1 & -t \\
t & 1
\end{bmatrix} 
\begin{bmatrix}
t^2 \\
-\frac{1}{3}\,t^3 + t
\end{bmatrix} 
\\
& =
\begin{bmatrix}
\frac{1}{3}\,t^4 \\
\frac{2}{3}\,t^3 + t
\end{bmatrix}  .
\end{split}
\end{equation*}
Dengan menambahkan solusi komplementer, kita memperoleh solusi umum
dari \eqref{nhsys:vcexeq}:
\begin{equation*}
\vec{x} =
\begin{bmatrix}
1 & -t \\
t & 1
\end{bmatrix}
\begin{bmatrix}
c_1 \\ c_2
\end{bmatrix}
+
\begin{bmatrix}
\frac{1}{3}\,t^4 \\
\frac{2}{3}\,t^3 + t
\end{bmatrix}
=
\begin{bmatrix}
c_1 - c_2 t
+
\frac{1}{3}\,t^4 \\
c_2 + 
(c_1 + 1)\, t
+
\frac{2}{3}\,t^3
\end{bmatrix} .
\end{equation*}
\end{example}

\begin{exercise}
Periksalah bahwa $x_1 = 
\frac{1}{3}\,t^4$ dan $x_2 = 
\frac{2}{3}\,t^3 + t$ benar-benar menyelesaikan
\eqref{nhsys:vcexeq}.
\end{exercise}

Dalam variasi parameter, seperti dalam metode faktor integrasi,
kita dapat memperoleh solusi umum dengan menambahkan konstanta integrasi.
Hal itu menambahkan $X(t) \vec{c}$ untuk vektor konstanta sebarang $\vec{c}$.  Namun,
itulah tepatnya solusi komplementer.

\subsection{Orde kedua berkoefisien konstan}

\subsubsection{Koefisien tak tentu}

Kita telah melihat contoh sederhana metode
koefisien
tak tentu\index{koefisien tak tentu!untuk sistem orde kedua}
untuk sistem orde kedua dalam \sectionref{sol:section}.
Metode ini pada dasarnya sama dengan metode koefisien tak tentu
untuk sistem orde pertama.
Ada beberapa penyederhanaan yang dapat kita lakukan, seperti dalam 
\sectionref{sol:section}.  Tinjau persamaan
\begin{equation*}
{\vec{x}}'' = A \vec{x} + \vec{F}(t) ,
\end{equation*}
dengan $A$ matriks konstan.  Jika $\vec{F}(t)$ berbentuk
$\vec{F}_0 \cos (\omega t)$, maka karena dua kali turunan kosinus kembali menjadi
kosinus,
kita tidak perlu memasukkan sinus, dan
kita mencoba solusi berbentuk
\begin{equation*}
\vec{x}_p = \vec{c} \cos (\omega t) .
\end{equation*}

Jika $\vec{F}$ merupakan jumlah
kosinus, ingat prinsip superposisi.
Jika $\vec{F}(t) =
\vec{F}_0 \cos (\omega_0 t) + 
\vec{F}_1 \cos (\omega_1 t)$, maka kita mencoba
$\vec{a} \cos (\omega_0 t)$ untuk masalah
${\vec{x}}'' = A \vec{x} + \vec{F}_0 \cos (\omega_0 t)$, dan kita mencoba
$\vec{b} \cos (\omega_1 t)$ untuk masalah
${\vec{x}}'' = A \vec{x} + \vec{F}_1 \cos (\omega_1 t)$.  Kemudian kita jumlahkan
solusi-solusinya.

Jika terdapat pengulangan dengan solusi komplementer, atau persamaannya
berbentuk
${\vec{x}}'' = A{\vec{x}}'+ B \vec{x} + \vec{F}(t)$, maka kita perlu
melakukan hal yang sama seperti pada sistem orde pertama.

Anda tidak akan keliru dengan memasukkan lebih banyak suku daripada yang diperlukan ke dalam
tebakan.  Koefisien tambahan tersebut akan ternyata
bernilai nol.  Namun, menghemat waktu dan usaha tetaplah berguna.

\subsubsection{Dekomposisi vektor eigen}

Jika kita memiliki sistem
\begin{equation*}
{\vec{x}}'' = A \vec{x} + \vec{f}(t) ,
\end{equation*}
kita dapat melakukan \emph{\myindex{dekomposisi vektor eigen}},
seperti untuk sistem orde pertama.

Misalkan $\lambda_1, \lambda_2, \ldots, \lambda_n$ nilai-nilai eigen
dan $\vec{v}_1, \vec{v}_2, \ldots, \vec{v}_n$ vektor-vektor eigen.
Sekali lagi bentuklah matriks %$E = [\, \vec{v}_1 \cdots \vec{v}_n \,]$.
$E = [\, \vec{v}_1 \quad \vec{v}_2 \quad \cdots \quad \vec{v}_n \,]$.
Tuliskan
\begin{equation*}
\vec{x}(t) =
\vec{v}_1 \, \xi_1(t) + 
\vec{v}_2 \, \xi_2(t) + \cdots +
\vec{v}_n \, \xi_n(t) .
\end{equation*}
Uraikan $\vec{f}$ dalam vektor-vektor eigen
\begin{equation*}
\vec{f}(t) =
\vec{v}_1 \, g_1(t) + 
\vec{v}_2 \, g_2(t) + \cdots +
\vec{v}_n \, g_n(t) ,
\end{equation*}
dengan, sekali lagi, $\vec{g} = E^{-1} \vec{f}$.

Kita substitusikan dan, seperti sebelumnya, memperoleh
\begin{equation*}
\begin{split}
%{\vec{x}}'' & =
\overbrace{
\vec{v}_1 \xi_1'' + 
\vec{v}_2 \xi_2'' + \cdots +
\vec{v}_n \xi_n''
}^{{\vec{x}}''}
%\\
& =
\overbrace{
A \left( \vec{v}_1 \xi_1 + 
\vec{v}_2 \xi_2 + \cdots +
\vec{v}_n \xi_n \right)
}^{A\vec{x}}
+
\overbrace{
\vec{v}_1 g_1 + 
\vec{v}_2 g_2 + \cdots +
\vec{v}_n g_n
}^{\vec{f}}
\\
& = 
A \vec{v}_1 \xi_1 + 
A \vec{v}_2 \xi_2 + \cdots +
A \vec{v}_n \xi_n
+
\vec{v}_1 g_1 + 
\vec{v}_2 g_2 + \cdots +
\vec{v}_n g_n
\\
& =
\vec{v}_1 \lambda_1 \xi_1 + 
\vec{v}_2 \lambda_2 \xi_2 + \cdots +
\vec{v}_n \lambda_n \xi_n
+
\vec{v}_1 \, g_1 + 
\vec{v}_2 \, g_2 + \cdots +
\vec{v}_n \, g_n
\\
& =
\vec{v}_1 ( \lambda_1 \xi_1 + g_1 ) +
\vec{v}_2 ( \lambda_2 \xi_2 + g_2 ) + \cdots + 
\vec{v}_n ( \lambda_n \xi_n + g_n ) .
\end{split}
\end{equation*}
Kita samakan koefisien vektor-vektor eigen untuk memperoleh persamaan
\begin{align*}
\xi_1'' & = \lambda_1 \xi_1 + g_1 , \\
\xi_2'' & = \lambda_2 \xi_2 + g_2 , \\
& ~~ \vdots \\
\xi_n'' & = \lambda_n \xi_n + g_n .
\end{align*}
Setiap persamaan ini tidak bergantung pada persamaan lainnya.
Kita selesaikan setiap persamaan dengan metode dalam \chapterref{ho:chapter}.
Kita tuliskan
$\vec{x}(t) =
\vec{v}_1 \xi_1(t) +
\vec{v}_2 \xi_2(t) + \cdots +
\vec{v}_n \xi_n(t)$ untuk mencari solusi khusus.
Jika kita mencari solusi umum untuk $\xi_1$ sampai $\xi_n$, maka
$\vec{x}(t) =
\vec{v}_1 \xi_1(t) + \vec{v}_2 \, \xi_2(t) + \cdots +
\vec{v}_n \xi_n(t)$ juga merupakan solusi umum.

\begin{example}
Mari kita kerjakan contoh dari \sectionref{sol:section} dengan metode ini.
Persamaannya adalah
\begin{equation*}
{\vec{x}}'' =
\begin{bmatrix}
-3 & 1 \\
2 & -2
\end{bmatrix}
\vec{x} 
+ 
\begin{bmatrix}
0 \\ 2
\end{bmatrix}
\cos (3 t) .
\end{equation*}
Nilai eigennya adalah $-1$ dan $-4$, dengan vektor eigen 
$\left[ \begin{smallmatrix} 1 \\ 2 \end{smallmatrix} \right]$ dan
$\left[ \begin{smallmatrix} 1 \\ -1 \end{smallmatrix} \right]$.
Oleh karena itu, $E =
\left[ \begin{smallmatrix} 1 & 1 \\ 2 & -1 \end{smallmatrix} \right]$
dan
$E^{-1} =
\frac{1}{3}
\left[ \begin{smallmatrix} 1 & 1 \\ 2 & -1 \end{smallmatrix}
\right]$.
Dengan demikian,
\begin{equation*}
\begin{bmatrix}
g_1 \\
\noalign{\smallskip}
g_2
\end{bmatrix}
=
E^{-1} \vec{f}(t)
=
\frac{1}{3}
\begin{bmatrix} 1 & 1 \\
\noalign{\smallskip}
2 & -1 \end{bmatrix}
\begin{bmatrix} 0 \\
\noalign{\smallskip}
2\cos (3t) \end{bmatrix}
=
\begin{bmatrix} \frac{2}{3}\cos (3t)
\\
\noalign{\smallskip}
\frac{-2}{3}\cos (3t) \end{bmatrix} .
\end{equation*}
Jadi, setelah seluruh proses substitusi, persamaan yang kita peroleh adalah
\begin{equation*}
\xi_1'' = - \xi_1 + \frac{2}{3} \cos (3t) , \qquad
\xi_2'' = -4 \, \xi_2 - \frac{2}{3} \cos (3t) .
\end{equation*}
Untuk setiap persamaan, kita gunakan metode koefisien tak tentu.
Kita 
mencoba $C_1 \cos (3t)$ untuk persamaan pertama dan $C_2 \cos (3t)$ untuk persamaan
kedua.  Kita substitusikan untuk memperoleh
\begin{align*}
- 9 C_1 \cos (3t) & = - C_1 \cos (3t) + \frac{2}{3} \cos (3t) , \\
- 9 C_2 \cos (3t) & = - 4 C_2 \cos (3t) - \frac{2}{3} \cos (3t) .
\end{align*}
Kita selesaikan 
masing-masing persamaan secara terpisah.  Kita memperoleh
$- 9 C_1 = - C_1 + \nicefrac{2}{3}$ dan
$- 9 C_2 = - 4C_2 - \nicefrac{2}{3}$.  Dengan demikian $C_1 = \nicefrac{-1}{12}$
dan $C_2 = \nicefrac{2}{15}$.
Jadi solusi khusus kita adalah
\begin{equation*}
\vec{x} =
\begin{bmatrix} 1 \\
2 \end{bmatrix} \,
\left( \frac{-1}{12} \, \cos (3t) \right)
+
\begin{bmatrix} 1 \\
-1 \end{bmatrix} \,
\left( \frac{2}{15} \, \cos (3t) \right)
=
\begin{bmatrix} \nicefrac{1}{20} \\
\nicefrac{-3}{10} \end{bmatrix} \,
\cos (3t) .
\end{equation*}
Solusi ini sama dengan yang sebelumnya kita peroleh dalam \sectionref{sol:section}.
\end{example}

\subsection{Latihan}

\begin{exercise}
Carilah solusi khusus dari
$x' = x+ 2y +2t$, 
$y' = 3x + 2y -4$
\begin{tasks}(2)
\task dengan metode faktor integrasi,
\task dengan dekomposisi vektor eigen,
\task dengan koefisien tak tentu.
\end{tasks}
\end{exercise}

\begin{exercise}
Carilah solusi umum dari
$x' = 4x+ y -1$, 
$y' = x + 4y -e^t$
\begin{tasks}(2)
\task dengan metode faktor integrasi,
\task dengan dekomposisi vektor eigen,
\task dengan koefisien tak tentu.
\end{tasks}
\end{exercise}

\begin{exercise}
Carilah solusi umum dari
$x_1'' = -6x_1+ 3x_2 + \cos (t)$, 
$x_2'' = 2x_1 -7x_2 + 3\cos (t)$
\begin{tasks}(2)
\task dengan dekomposisi vektor eigen,
\task dengan koefisien tak tentu.
\end{tasks}
\end{exercise}

\begin{exercise}
Carilah solusi umum dari
$x_1'' = -6x_1+ 3x_2 + \cos (2t)$, 
$x_2'' = 2x_1 -7x_2 + 3\cos (2t)$
\begin{tasks}(2)
\task dengan dekomposisi vektor eigen,
\task dengan koefisien tak tentu.
\end{tasks}
\end{exercise}

\begin{exercise}
Ambil persamaan
$\displaystyle
{\vec{x}}'
=
\begin{bmatrix}
\frac{1}{t} & -1 \\
1 & \frac{1}{t}
\end{bmatrix}
\vec{x}
+ \begin{bmatrix} t^2 \\ -t \end{bmatrix} .
$
\begin{tasks}
\task
Periksalah bahwa
$\displaystyle
\vec{x}_c =
c_1
\begin{bmatrix}
t\, \sin t \\
- t \, \cos t
\end{bmatrix}
+
c_2
\begin{bmatrix}
t\, \cos t \\
t \, \sin t
\end{bmatrix}
$
adalah solusi komplementer.
\task
Gunakan variasi parameter untuk
mencari solusi khusus.
\end{tasks}
\end{exercise}

\setcounter{exercise}{100}

\begin{exercise}
Carilah solusi khusus dari $x' = 5x + 4y + t$, $y' = x + 8y - t$
\begin{tasks}(2)
\task dengan metode faktor integrasi,
\task dengan dekomposisi vektor eigen,
\task dengan koefisien tak tentu.
\end{tasks}
\end{exercise}
\exsol{%
Solusi umumnya adalah
(solusi khusus harus bersesuaian dengan salah satu bentuk ini):\\
$x(t) = C_1 e^{9 t}+4 C_2 e^{4t}-\nicefrac{t}{3}-\nicefrac{5}{54}$,
\quad
$y(t) = C_1 e^{9 t}-C_2 e^{4t}+\nicefrac{t}{6}+\nicefrac{7}{216}$
}

\begin{exercise}
Carilah solusi khusus dari $x' = y + e^t$, $y' = x +e^t$
\begin{tasks}(2)
\task dengan metode faktor integrasi,
\task dengan dekomposisi vektor eigen,
\task dengan koefisien tak tentu.
\end{tasks}
\end{exercise}
\exsol{%
Solusi umumnya adalah
(solusi khusus harus bersesuaian dengan salah satu bentuk ini):\\
$x(t) = C_1 e^t + C_2 e^{-t}+te^t$,
\quad
$y(t) = C_1 e^t - C_2 e^{-t}+te^t$
}

\begin{exercise}
Selesaikan
$x_1' = x_2 + t$, $x_2' = x_1 +t$ dengan syarat awal
$x_1(0) = 1$, $x_2(0) = 2$
menggunakan dekomposisi vektor eigen.
\end{exercise}
\exsol{%
$\vec{x} = 
\left[ \begin{smallmatrix}
1 \\ 1
\end{smallmatrix}\right]
\left(\frac{5}{2} e^t-t-1\right)
+
\left[ \begin{smallmatrix}
1 \\ -1
\end{smallmatrix}\right]
\frac{-1}{2} e^{-t}$
}

\begin{exercise}
Selesaikan
$x_1'' = -3x_1 + x_2 + t$, $x_2'' = 9x_1 + 5x_2
+\cos(t)$ dengan syarat awal
$x_1(0) = 0$, $x_2(0) = 0$,
$x_1'(0) = 0$, $x_2'(0) = 0$
menggunakan dekomposisi vektor eigen.
\end{exercise}
\exsol{%
%a) 
%eig:6,-4
%[1,9], [1,-1]
%
%xi(0)
%
%xi1 = c_1 e^(sqrt(6) t)+c_2 e^(-sqrt(6) t)-t\/60-(cos(t))\/70
%xi2 = c_2 sin(2 t)+c_1 cos(2 t)+(9 t)\/40-(cos(t))\/30
%
$\vec{x} = 
\left[ \begin{smallmatrix}
1 \\ 9
\end{smallmatrix}\right]
\left(\left(\frac{1}{140} + \frac{1}{120\sqrt{6}}\right) e^{\sqrt{6} t}
+
\left(\frac{1}{140} - \frac{1}{120\sqrt{6}}\right) e^{-\sqrt{6} t}
-\frac{t}{60}-\frac{\cos(t)}{70} \right)
$
\\
$
+
\left[ \begin{smallmatrix}
1 \\ -1
\end{smallmatrix}\right]
\left(\frac{-9}{80} \sin(2 t)+ \frac{1}{30} \cos(2 t)+\frac{9 t}{40}-\frac{\cos(t)}{30}\right)$
%
%$\vec{x} (0) = 
%\left[ \begin{smallmatrix}
%1 \\ 9
%\end{smallmatrix}\right]
%(c_1 +c_2 -1/70)
%+
%\left[ \begin{smallmatrix}
%1 \\ -1
%\end{smallmatrix}\right]
%(d_1-1/30)$
%
%$0 = (c_1+c_2-1/70) + (d_1-1/30)$
%$0 = 9(c_1+c_2-1/70) - (d_1-1/30)$
%$0 = c_1+c_2-1/70$
%$0 = (d_1-1/30)$
%
%$d_1 = 1/30$
%
%$\vec{x}' = 
%\left[ \begin{smallmatrix}
%1 \\ 9
%\end{smallmatrix}\right]
%(c_1 sqrt(6) e^(sqrt(6) t)- sqrt(6)c_2 e^(-sqrt(6) t)-1/60+(sin(t))\/70)
%+
%\left[ \begin{smallmatrix}
%1 \\ -1
%\end{smallmatrix}\right]
%(2d_2 cos(2 t)-2d_1 sin(2 t)+9/40+(sin(t))/30)$
%
%$\vec{x}'(0) = 
%\left[ \begin{smallmatrix}
%1 \\ 9
%\end{smallmatrix}\right]
%(c_1 sqrt(6) - sqrt(6)c_2 -1/60)
%+
%\left[ \begin{smallmatrix}
%1 \\ -1
%\end{smallmatrix}\right]
%(2d_2 +9/40)$
%
%$0 = (c_1 sqrt(6) - c_2 sqrt(6) - 1/60) + (2d_2 + 9/40)$
%$0 = 9(c_1 sqrt(6) - c_2 sqrt(6) - 1/60) - (2d_2 + 9/40)$
%
%$0 = c_1 sqrt(6) - c_2 sqrt(6) - 1/60$
%$d_2 = -9/80$
%
%$0 = c_1 sqrt(6) - c_2 sqrt(6) - 1/60$
%$0 = c_1+c_2-1/70$
%
%$c_1 + c_2 = 1/70$
%$c_1 - c_2 = 1/(60sqrt(6))$
%
%$c_1 = 1/140 + 1/(120sqrt(6))$
%$c_2 = 1/140 - 1/(120sqrt(6))$
}
